% PDF page 100
\source{gilbarg-trudinger:page-0100}




let $[A^{ij}]$ be a constant matrix such that

\[
\lambda |\xi|^2 \le A^{ij}\xi_i \xi_j \le \Lambda |\xi|^2, \qquad \forall \xi \in \R^n
\]

for positive constants $\lambda$, $\Lambda$.

\medskip

\noindent\textbf{(a)} \emph{Let $u \in C^2(\Omega)$, $f \in C^\alpha(\Omega)$ satisfy $L_0u=f$ in an open set $\Omega$ of $\R^n$. Then}

\[
(6.4)\qquad |u|_{2,\alpha;\Omega}^* \le C(|u|_{0;\Omega}+|f|_{0,\alpha;\Omega}^{(2)})
\]

\emph{where $C = C(n,\alpha,\lambda,\Lambda)$.}

\medskip

\noindent\textbf{(b)} \emph{Let $\Omega$ be an open subset of $\Rnp$ with a boundary portion $T$ on $x_n=0$, and let $u \in C^2(\Omega)\cap C^0(\Omega \cup T)$, $f \in C^\alpha(\Omega \cup T)$ satisfy $L_0u=f$ in $\Omega$, $u=0$ on $T$. Then}

\[
(6.5)\qquad |u|_{2,\alpha;\Omega\cup T}^* \le C(|u|_{0;\Omega}+|f|_{0,\alpha;\Omega\cup T}^{(2)}),
\]

\emph{where $C = C(n,\alpha,\lambda,\Lambda)$.}

\medskip

\emph{Proof.} Let $P$ be a constant matrix which defines a nonsingular linear transformation $y=xP$ from $\Rn$ onto $\Rn$. Letting $u(x)\to \tilde u(y)$ under this transformation one verifies easily that

\[
A^{ij}D_{ij}u(x)=\tilde A^{ij}D_{ij}\tilde u(y),
\]

\emph{where $\tilde A = P^tAP$ ($P^t = P$ transpose). For a suitable orthogonal matrix $P$, $\tilde A$ is a diagonal matrix whose diagonal elements are the eigenvalues $\lambda_1,\ldots,\lambda_n$ of $A$. If, further, $Q=PD$, where $D$ is the diagonal matrix $[\lambda_i^{-1/2}\delta_{ij}]$, then the transformation $y=xQ$ takes $L_0u=f(x)$ into the Poisson equation $\Delta \tilde u(y)=\tilde f(y)$, where $u(x)\to \tilde u(y)$, $f(x)\to \tilde f(y)$ under the transformation. By a further rotation we may assume that $Q$ takes the half-space $x_n>0$ onto the half-space $y_n>0$.}

\medskip

\emph{Since the orthogonal matrix $P$ preserves length, we have}

\[
\Lambda^{-1/2}|x| \le |xQ| \le \lambda^{-1/2}|x|.
\]

\emph{It follows that if $\Omega \to \tilde \Omega$, $v(x)\to \tilde v(y)$ under the transformation $y=xQ$ then the norms (4.17), (4.18) defined on $\Omega$ and $\tilde \Omega$ are related by the inequalities}

\[
\begin{aligned}
c^{-1}|v|_{k,\alpha;\Omega}^* &\le |\tilde v|_{k,\alpha;\tilde \Omega}^* \le c|v|_{k,\alpha;\Omega}^*, \qquad k=0,1,2,\ldots,\qquad 0\le \alpha \le 1, \\
c^{-1}|v|_{0,\alpha;\Omega}^{(k)} &\le |\tilde v|_{0,\alpha;\tilde \Omega}^{(k)} \le c|v|_{0,\alpha;\Omega}^{(k)}
\end{aligned}
\]

\emph{where $c=c(k,n,\lambda,\Lambda)$.}

\medskip

\emph{Similarly if $\Omega$ is an open subset of $\Rnp$ with a boundary portion $T$ on $x_n=0$ which is taken by $y=xQ$ into the set $\tilde \Omega$ in $\Rnp$ with boundary portion $\tilde T$, the norms}


% PDF page 101
\source{gilbarg-trudinger:page-0101}

\[
(4.29)\ \text{in }\Omega\text{ and }\widetilde{\Omega}\text{ satisfy the inequalities}
\]

\begin{equation}\tag{6.7}
c^{-1}\lvert v\rvert^{*}_{k,\alpha;\Omega\cup\widetilde{\Omega}}
\leq \lvert \widetilde{v}\rvert^{*}_{k,\alpha;\Omega\cup\widetilde{\Omega}}
\leq c\lvert v\rvert^{*}_{k,\alpha;\Omega\cup\widetilde{\Omega}},
\end{equation}
\begin{equation*}
c^{-1}\lvert v\rvert^{(k)}_{0,\alpha;\Omega\cup\widetilde{\Omega}}
\leq \lvert \widetilde{v}\rvert^{(k)}_{0,\alpha;\Omega\cup\widetilde{\Omega}}
\leq c\lvert v\rvert^{(k)}_{0,\alpha;\Omega\cup\widetilde{\Omega}},
\end{equation*}

where $c$ is the same constant as in (6.6).

To prove part (a) of the lemma we apply Theorem 4.8 in $\widetilde{\Omega}$ and the inequalities
(6.6) to obtain

\[
\lvert u\rvert^{*}_{2,\alpha;\Omega}
\leq C\lvert \widetilde{u}\rvert^{*}_{2,\alpha;\widetilde{\Omega}}
\leq C\bigl(\lvert \widetilde{u}\rvert_{0;\widetilde{\Omega}}
+\lvert \widetilde{f}\rvert^{(2)}_{0,\alpha;\widetilde{\Omega}}\bigr)
\leq C\bigl(\lvert u\rvert_{0;\Omega}
+\lvert f\rvert^{(2)}_{0,\alpha;\Omega}\bigr),
\]

which is the desired conclusion (6.4). (Here we have used the same letter $C$ to de-
note constants depending on $n$, $\alpha$, $\lambda$, $\Lambda$.)

Part (b) of the lemma is proved in the same way, using Theorem 4.12 and the
inequalities (6.7). $\square$

Lemma 6.1 provides an immediate extension of the estimates on balls in
Theorems 4.6 and 4.11 from Poisson's equation to the more general equation with
constant coefficients (6.3). Of course in the latter case the constant $C$ depends on
$\lambda$, $\Lambda$ as well as $n$, $\alpha$.

\section*{6.1. The Schauder Interior Estimates}

Our first objective in the study of the equation $Lu=f$ is the derivation of the Schauder
interior estimates, which play an essential part in the ensuing treatment of the
existence and regularity theory. These estimates are based on the same kind of
result already obtained in (6.4) for solutions of $L_0u=f$.

To obtain estimates for the interior norm $\lvert u\rvert^{*}_{2,\alpha;\Omega}$ of solutions of $Lu=f$ in $\Omega$,
it suffices to bound only $\lvert u\rvert_{0;\Omega}$ and the seminorm $[u]^{*}_{2,\alpha;\Omega}$ (defined in (4.17)). That
this is so is a consequence of the following interpolation inequalities.

\emph{Let $u\in C^{2,\alpha}(\Omega)$, where $\Omega$ is an open subset of $\mathbb{R}^n$. Then for any $\varepsilon>0$ there is a
constant $C=C(\varepsilon)$ such that}

\begin{equation}\tag{6.8}
[u]^{*}_{j;\beta;\Omega}\leq C\lvert u\rvert_{0;\Omega}+\varepsilon [u]^{*}_{2,\alpha;\Omega},
\qquad j=0,\,1,\,2;\ \ 0\leq \alpha,\beta\leq 1,\ j+\beta<2+\alpha;
\end{equation}
\begin{equation}\tag{6.9}
\lvert u\rvert^{*}_{j;\beta;\Omega}\leq C\lvert u\rvert_{0;\Omega}+\varepsilon [u]^{*}_{2,\alpha;\Omega}.
\end{equation}

These inequalities are proved in Lemma 6.32 of Appendix 1 to this chapter.
In order to state the Schauder estimates in a sharp form, and also for subse-
quent applications, we introduce the following additional interior seminorms and

% PDF page 102
\source{gilbarg-trudinger:page-0102}




norms on the spaces $C^k(\Omega)$, $C^{k,\alpha}(\Omega)$. For $\sigma$ a real number and $k$ a non-negative integer we define

\begin{equation}
\left[f\right]_{k,0;\Omega}^{(\sigma)}=\left[f\right]_{k;\Omega}^{(\sigma)}=\sup_{x\in\Omega \atop |\beta|=k} d_x^{k+\sigma}\abs{D^\beta f(x)};
\end{equation}

\begin{equation}
\left[f\right]_{k,\alpha;\Omega}^{(\sigma)}=\sup_{x,y\in\Omega \atop |\beta|=k} d_{x,y}^{k+\alpha+\sigma}\frac{\abs{D^\beta f(x)-D^\beta f(y)}}{\abs{x-y}^\alpha},\qquad 0<\alpha\le 1;
\end{equation}

\begin{equation}
\tag{6.10}
\abs{f}_{k;\Omega}^{(\sigma)}=\sum_{j=0}^k \left[f\right]_{j;\Omega}^{(\sigma)},
\end{equation}

\begin{equation}
\abs{f}_{k,\alpha;\Omega}^{(\sigma)}=\abs{f}_{k;\Omega}^{(\sigma)}+\left[f\right]_{k,\alpha;\Omega}^{(\sigma)}.
\end{equation}

In this notation, when $\sigma=0$ these quantities are identical with those defined in (4.17), so that $\left[\cdot\right]^{(0)}=\left[\cdot\right]^*$ and $\abs{\cdot}^{(0)}=\abs{\cdot}^*$.

It is easy to verify that

\begin{equation}
\tag{6.11}
\abs{fg}_{0,x;\Omega}^{(\sigma+\tau)}\le \abs{f}_{0,x;\Omega}^{(\sigma)}\abs{g}_{0,x;\Omega}^{(\tau)} \qquad \text{for }\sigma+\tau\ge 0.
\end{equation}

We now establish the basic Schauder interior estimates.

\begin{theorem}
\label{gt-ch06-thm-6-2}
\dcref{Theorem 6.2}
\group{gt-ch06}
\source{gilbarg-trudinger:page-0103}
\textit{Let $\OO$ be an open subset of $\R^n$, and let $u\in C^{2,\alpha}(\OO)$ be a bounded solution in $\OO$ of the equation}

\[
Lu=a^{ij}D_{ij}u+b^iD_i u+cu=f.
\]

\textit{where $f\in C^\alpha(\OO)$ and there are positive constants $\lambda$, $\Lambda$ such that the coefficients satisfy}

\begin{equation}
\tag{6.12}
a^{ij}\xi_i\xi_j\ge \lambda\abs{\xi}^2 \qquad \forall x\in \OO,\ \xi\in \R^n
\end{equation}

\textit{and}

\begin{equation}
\tag{6.13}
\abs{a^{ij}}_{0,x;\OO}^{(0)}+\abs{b^i}_{0,x;\OO}^{(1)}+\abs{c}_{0,x;\OO}^{(2)}\le \Lambda.
\end{equation}

\textit{Then}

\begin{equation}
\tag{6.14}
\left|u\right|_{2,x;\OO}^*\le C\bigl(\left|u\right|_{0,\OO}+\left|f\right|_{0,x;\OO}^{(2)}\bigr)
\end{equation}

\textit{where $C=C(n,\alpha,\lambda,\Lambda)$.}

\textit{Proof.} By virtue of (6.9) it suffices to prove the inequality (6.14) for $\left[u\right]^*_{2,x;\OO}$, and it suffices to prove the latter for compact subsets of $\OO$. Namely, let $\{\OO_i\}$ be a sequence of open subsets of $\OO$ such that $\OO_i\subset \OO_{i+1}\subset \OO$ and $\bigcup \OO_i=\OO$. We have that


% PDF page 103
\source{gilbarg-trudinger:page-0103}



\noindent 6.1. The Schauder Interior Estimates \hfill 91

\[
[u]^{*}_{2,\alpha;\Omega_i} \text{ is finite for each } i \text{ and if (6.14) holds in } \Omega_i, \text{ we may assert for any pair of}
\]
\[
\text{points } x,y \in \Omega,\ \text{for all sufficiently large } i,\text{ and for any second derivative } D^2u,
\]
\[
\left(d^{(i)}_{x,y}\right)^{2+\alpha}\frac{\lvert D^2u(x)-D^2u(y)\rvert}{\lvert x-y\rvert^\alpha}
\le [u]^{*}_{2,\alpha;\Omega_i}
\le C\!\left(\lvert u\rvert_{0;\Omega_i}+\lvert f\rvert^{(2)}_{0;\alpha;\Omega_i}\right)
\]
\[
\le C\!\left(\lvert u\rvert_{0;\Omega}+\lvert f\rvert^{(2)}_{0;\alpha;\Omega}\right)
\]

\noindent where $d^{(i)}_{x,y}=\min[\operatorname{dist}(x,\partial\Omega_i),\operatorname{dist}(y,\partial\Omega_i)]$. Letting $i\to\infty$, we obtain the inequality

\[
d_{x,y}^{2+\alpha}\frac{\lvert D^2u(x)-D^2u(y)\rvert}{\lvert x-y\rvert^\alpha}
\le C\!\left(\lvert u\rvert_{0;\Omega}+\lvert f\rvert^{(2)}_{0;\alpha;\Omega}\right)
\qquad \forall x,y\in\Omega,
\]

\noindent which implies the same bound for $[u]^{*}_{2,\alpha;\Omega}$. We may therefore assume in the following that $[u]^{*}_{2,\alpha;\Omega}$ is finite.

\noindent For notational convenience the same letter $C$ will be used to denote constants depending on $n,\alpha,\lambda,\Lambda$.

\noindent Let $x_0,y_0$ be any two distinct points in $\Omega$ and suppose $d_{x_0}=d_{x_0,y_0}=\min(d_{x_0},d_{y_0})$.
Let $\mu\le \tfrac12$ be a positive constant to be specified later, and set $d=\mu d_{x_0}$, $B=B_d(x_0)$.
We rewrite $Lu=f$ in the form

\begin{equation}
(6.15)\qquad a^{ij}(x_0)D_{ij}u=\bigl(a^{ij}(x_0)-a^{ij}(x)\bigr)D_{ij}u-b^iD_iu-cu+f\equiv F(x),
\end{equation}

\noindent and we consider this as an equation in $B$ with constant coefficients $a^{ij}(x_0)$. Lemma 6.1(a), applied to this equation, asserts that if $y_0\in B_{d/2}(x_0)$, then for any second derivative $D^2u$

\[
\left(\frac d2\right)^{2+\alpha}\frac{\lvert D^2u(x_0)-D^2u(y_0)\rvert}{\lvert x_0-y_0\rvert^\alpha}
\le C\!\left(\lvert u\rvert_{0;B}+\lvert F\rvert^{(2)}_{0;\alpha;B}\right);
\]

\noindent and thus

\[
d_{x_0}^{2+\alpha}\frac{\lvert D^2u(x_0)-D^2u(y_0)\rvert}{\lvert x_0-y_0\rvert^\alpha}
\le \frac{C}{\mu^{2+\alpha}}\left(\lvert u\rvert_{0;B}+\lvert F\rvert^{(2)}_{0;\alpha;B}\right).
\]

\noindent On the other hand, if $\lvert x_0-y_0\rvert\ge d/2$,

\[
d_{x_0}^{2+\alpha}\frac{\lvert D^2u(x_0)-D^2u(y_0)\rvert}{\lvert x_0-y_0\rvert^\alpha}
\le \left(\frac{2}{\mu}\right)^\alpha \bigl[d_{x_0}^2\lvert D^2u(x_0)\rvert+d_{y_0}^2\lvert D^2u(y_0)\rvert\bigr]
\]
\[
\le \frac{4}{\mu^\alpha}[u]^{*}_{2;\Omega},
\]

\noindent so that, combining these two inequalities, we obtain

\begin{equation}
(6.16)\qquad d_{x_0}^{2+\alpha}\frac{\lvert D^2u(x_0)-D^2u(y_0)\rvert}{\lvert x_0-y_0\rvert^\alpha}
\le \frac{C}{\mu^{2+\alpha}}\left(\lvert u\rvert_{0;\Omega}+\lvert F\rvert^{(2)}_{0;\alpha;B}\right)
+ \frac{4}{\mu^\alpha}[u]^{*}_{2;\Omega}.
\end{equation}


% PDF page 104
\source{gilbarg-trudinger:page-0104}

We proceed to estimate $\abs{F}^{(2)}_{0,\alpha;B}$ in terms of $\abs{u}_{0;\Omega}$ and $\Hs$. We have

\begin{equation}\tag{6.17}
\abs{F}^{(2)}_{0,\alpha;B}\leq \sum_{i,j}\abs{(a^{ij}(x_0)-a^{ij}(x))D_{ij}u}^{(2)}_{0,\alpha;B}+\sum_i \abs{b^iD_i u}^{(2)}_{0,\alpha;B}
\end{equation}
\[
\qquad +\abs{cu}^{(2)}_{0,\alpha;B}+\abs{f}^{(2)}_{0,\alpha;B}.
\]

It will be useful in estimating these terms to have the following inequality.
Recalling that for all $x\in B$, $d_x(=\dist(x,\partial\Omega))>(1-\mu)d_{x_0}\geq \frac12 d_{x_0}$, we have for
$g\in C^\alpha(\Omega)$

\begin{equation}\tag{6.18}
\abs{g}^{(2)}_{0,\alpha;B}\leq d^2\abs{g}_{0;B}+d^{2+\alpha}\abs{g}_{\alpha;B}
\leq \frac{\mu^2}{(1-\mu)^2}\abs{g}^{(2)}_{0;\Omega}+\frac{\mu^{2+\alpha}}{(1-\mu)^{2+\alpha}}\abs{g}^{(2)}_{0,\alpha;\Omega}
\end{equation}
\[
\qquad \leq 4\mu^2\abs{g}^{(2)}_{0;\Omega}+8\mu^{2+\alpha}\abs{g}^{(2)}_{0,\alpha;\Omega}\leq 8\mu^2\abs{g}^{(2)}_{0,\alpha;\Omega}.
\]

Writing $(a(x_0)-a(x))D^2u=(a^{ij}(x_0)-a^{ij}(x))D_{ij}u$ for each pair $i,j$, we obtain
from (6.11) and (6.18)

\[
\abs{(a(x_0)-a(x))D^2u}^{(2)}_{0,\alpha;B}\leq \abs{a(x_0)-a(x)}^{(0)}_{0,\alpha;B}\abs{D^2u}^{(2)}_{0,\alpha;B}
\]
\[
\qquad \leq \abs{a(x_0)-a(x)}^{(0)}_{0,\alpha;B}\bigl(4\mu^2[u]^{\ast}_{2;\Omega}+8\mu^{2+\alpha}[u]^{\ast}_{2,\alpha;\Omega}\bigr).
\]

Since

\[
\abs{a(x_0)-a(x)}^{(0)}_{0,\alpha;B}\leq \sup_{x\in B}\abs{a(x_0)-a(x)}+d^\alpha[a]_{\alpha;B}\leq 2d^\alpha[a]_{\alpha;B}
\]
\[
\qquad \leq 2^{1+\alpha}\mu^\alpha[u]^{\ast}_{0,\alpha;\Omega}\leq 4\Lambda\mu^\alpha,
\]

we arrive at the following estimate for the principal term in (6.17),

\begin{equation}\tag{6.19}
\sum_{i,j}\abs{(a^{ij}(x_0)-a^{ij}(x))D_{ij}u}^{(2)}_{0,\alpha;B}\leq 32n^2\Lambda\mu^{2+\alpha}\bigl([u]^{\ast}_{2;\Omega}+\mu^\alpha[u]^{\ast}_{2,\alpha;\Omega}\bigr)
\end{equation}
\[
\qquad \leq 32n^2\Lambda\mu^{2+\alpha}\bigl(C(\mu)\abs{u}_{0;\Omega}+2\mu^\alpha[u]^{\ast}_{2,\alpha;\Omega}\bigr).
\]

The last inequality is obtained by setting $\varepsilon=\mu^\alpha$ in the interpolation inequality (6.8).
Letting $bDu=b^iD_i u$ for each $i$, we obtain from (6.18) and (6.13)

\[
\abs{bDu}^{(2)}_{0,\alpha;B}\leq 8\mu^2\abs{bDu}^{(2)}_{0,\alpha;\Omega}\leq 8\mu^2\abs{b}^{(1)}_{0,\alpha;\Omega}\abs{Du}^{(1)}
\]
\[
\qquad \leq 8\mu^2\Lambda\abs{u}^{\ast}_{1;\alpha;\Omega}\leq 8\mu^2\Lambda(C(\mu))\abs{u}_{0;\Omega}+\mu^{2\alpha}[u]^{\ast}_{2,\alpha;\Omega}.
\]

The last inequality is obtained by setting $\varepsilon=\mu^{2\alpha}$ in (6.9). Thus we have

\begin{equation}\tag{6.20}
\abs{b^iD_i u}^{(2)}_{0,\alpha;B}\leq 8n\Lambda\mu^2(C(\mu))\abs{u}_{0;\Omega}+\mu^{2\alpha}[u]^{\ast}_{2,\alpha;\Omega}.
\end{equation}

Similarly, from (6.18), (6.11) and (6.9), we obtain

\begin{equation}\tag{6.21}
\abs{cu}^{(2)}_{0,\alpha;B}\leq 8\mu^2\abs{c}^{(2)}_{0,\alpha;\Omega}\abs{u}^{(0)}_{0;\Omega}\leq 8\Lambda\mu^2(C(\mu))\abs{u}_{0;\Omega}+\mu^{2\alpha}[u]^{\ast}_{2,\alpha;\Omega}.
\end{equation}


% PDF page 105
\source{gilbarg-trudinger:page-0105}



\noindent 6.1. The Schauder Interior Estimates \hfill 93

\noindent Finally,

\begin{equation*}
(6.22)\qquad \lvert f\rvert^{(2)}_{0;\alpha;B}\le 8\mu^2 \lvert f\rvert^{(2)}_{0;\alpha;\Omega}.
\end{equation*}

\noindent Letting $C$ denote constants depending only on $n$, $\alpha$, $\lambda$, $\Lambda$, and $C(\mu)$ constants
depending also on $\mu$, we find after combining (6.19)--(6.22),

\[
\lvert F\rvert^{(2)}_{0;\alpha;B}\le C\mu^{2+2\alpha}[u]^{*}_{2,\alpha;\Omega}+C(\mu)(\lvert u\rvert_{0;\Omega}+\lvert f\rvert^{(2)}_{0;\alpha;\Omega}).
\]

\noindent Inserting this into the right member of (6.16), and using (6.8) with $\varepsilon=\mu^{2\alpha}$ to esti-
mate $[u]^{*}_{2;\Omega}$, we obtain from (6.16)

\[
d^{2+\alpha}_{x_0,y_0}\frac{\lvert D^2u(x_0)-D^2u(y_0)\rvert}{\lvert x_0-y_0\rvert^\alpha}
\le C\mu^{\alpha}[u]^{*}_{2,\alpha;\Omega}+C(\mu)(\lvert u\rvert_{0;\Omega}+\lvert f\rvert^{(2)}_{0;\alpha;\Omega}).
\]

\noindent The right member of this inequality is independent of $x_0,y_0$. Taking the supremum
over all $x_0, y_0\in \Omega$, we obtain

\[
[u]^{*}_{2;\Omega}\le C\mu^{\alpha}[u]^{*}_{2,\alpha;\Omega}+C(\mu)(\lvert u\rvert_{0;\Omega}+\lvert f\rvert^{(2)}_{0;\alpha;\Omega}).
\]

\noindent We now choose and fix $\mu=\mu_0$ so that $C\mu_0^\alpha\le \tfrac12$. Then we arrive at the desired
estimate

\[
[u]^{*}_{2,\alpha;\Omega}\le C(\lvert u\rvert_{0;\Omega}+\lvert f\rvert^{(2)}_{0;\alpha;\Omega}).
\qquad \Box
\]
\end{theorem}

\noindent The preceding form of the interior estimates for $Lu=f$ permits the coefficients
and $f$ to be unbounded subject to the condition (6.13). In typical applications of the
interior estimates to convergence results, it suffices to know equicontinuity of
solutions and their derivatives up to second order on compact subsets. For this
purpose the following corollary is usually adequate.

\noindent \textbf{Corollary 6.3.} \textit{Let $u\in C^{2,\alpha}(\Omega)$, $f\in C^{\alpha}(\bar{\Omega})$ satisfy $Lu=f$ in a bounded domain $\Omega$
where $L$ satisfies (6.2) and its coefficients are in $C^{\alpha}(\bar{\Omega})$. Then if $\Omega' \Subset \Omega$ with
$\operatorname{dist}(\Omega',\partial\Omega)\ge d$, there is a constant $C$ such that}

\begin{equation*}
(6.23)\qquad d\lvert Du\rvert_{0;\Omega'}+d^2\lvert D^2u\rvert_{0;\Omega'}+d^{2+\alpha}[D^2u]_{\alpha;\Omega'}\le C(\lvert u\rvert_{0;\Omega}+\lvert f\rvert_{0;\alpha;\Omega})
\end{equation*}

\noindent \textit{where $C$ depends only on the ellipticity constant $\lambda$ and the $C^{\alpha}(\bar{\Omega})$ norms of the coef-
ficients of $L$ (as well as on $n$, $\alpha$ and the diameter of $\Omega$).}

\noindent \textit{Remark.} An immediate consequence of this result is that uniformly bounded
solutions of elliptic equations $Lu=f$ with locally H\"older continuous coefficients
and locally H\"older continuous $f$ are equicontinuous with their first and second
derivatives on compact subsets. This is true as well for the solutions of any \textit{family of}


% PDF page 106
\source{gilbarg-trudinger:page-0106}



\noindent 6. Classical Solutions; the Schauder Approach \hfill 94

\medskip

\noindent such equations with ellipticity constant $\lambda$ uniformly bounded away from zero in
compact subsets $\Omega' \Subset \Omega$ and with coefficients and inhomogeneous term $f$ having
uniformly bounded $C^{\alpha}(\bar{\Omega}')$ norms.

\medskip

\noindent \textbf{6.2. Boundary and Global Estimates}

\medskip

\noindent To extend the preceding interior estimates to the entire domain it is necessary to
have estimates that are meaningful near the boundary. These can be obtained
provided the boundary values of the solution and the boundary itself are sufficiently
smooth. In many respects the proof of these boundary estimates follows closely
that of the interior estimates.

\noindent \hspace*{2em}The global estimates which are the principal goal of this section will be estab-
lished in domains of class $C^{2,\alpha}$.

\medskip

\noindent \textbf{Definition.} A bounded domain $\Omega$ in $\mathbb{R}^n$ and its boundary are of class $C^{k,\alpha}$,
$0\le \alpha \le 1$, if at each point $x_0 \in \partial\Omega$ there is a ball $B=B(x_0)$ and a one-to-one mapping
$\psi$ of $B$ onto $D\subset \mathbb{R}^n$ such that:

\begin{equation*}
\text{(i) } \psi(B\cap \Omega)\subset \mathbb{R}^n_{+}; \qquad
\text{(ii) } \psi(B\cap \partial\Omega)\subset \partial\mathbb{R}^n_{+}; \qquad
\text{(iii) } \psi \in C^{k,\alpha}(B), \ \psi^{-1}\in C^{k,\alpha}(D).
\end{equation*}

\noindent A domain $\Omega$ will be said to have a boundary portion $\Gamma\subset \partial\Omega$ of class $C^{k,\alpha}$ if at each
point $x_0 \in \Gamma$ there is a ball $B=B(x_0)$ in which the above conditions are satisfied
and such that $B\cap \partial\Omega\subset \Gamma$. We shall say that the diffeomorphism $\psi$ straightens the
boundary near $x_0$.

\medskip

\noindent \hspace*{2em}We note in particular that $\Omega$ is a $C^{k,\alpha}$ domain if each point of $\partial\Omega$ has a neighbor-
hood in which $\partial\Omega$ is the graph of a $C^{k,\alpha}$ function of $n-1$ of the coordinates
$x_1,\ldots,x_n$. The converse is also true if $k\ge 1$.

\noindent \hspace*{2em}It follows from the above definition that a domain of class $C^{k,\alpha}$ is also of class
$C^{j,\beta}$ provided $j+\beta<k+\alpha$, $0\le \alpha,\beta \le 1$.

\noindent \hspace*{2em}A function $\varphi$ defined on a $C^{k,\alpha}$ boundary portion $\Gamma$ of a domain $\Omega$ will be said
to be in class $C^{k,\alpha}(\Gamma)$ if $\varphi \circ \psi^{-1}\in C^{k,\alpha}(D\cap \partial\mathbb{R}^n_{+})$ for each $x_0 \in \Gamma$. It is important to
note that if $\partial\Omega$ is in $C^{k,\alpha}(k\ge 1)$, then a function $\varphi \in C^{k,\alpha}(\partial\Omega)$ can be extended to
a function in $C^{k,\alpha}(\bar{\Omega})$, and conversely, a function in $C^{k,\alpha}(\bar{\Omega})$ has boundary values in
$C^{k,\alpha}(\partial\Omega)$ (see Lemma 6.38). Hence, in what follows it is immaterial whether we
consider boundary values $\varphi$ belonging to $C^{k,\alpha}(\partial\Omega)$ or $C^{k,\alpha}(\bar{\Omega})$.

\noindent \hspace*{2em}It is also possible to define a boundary norm on $C^{k,\alpha}(\partial\Omega)$, in various ways. For
example, if $\varphi \in C^{k,\alpha}(\partial\Omega)$ let $\Phi$ denote an extension of $\varphi$ to $\bar{\Omega}$ and define
$\|\varphi\|_{C^{k,\alpha}(\partial\Omega)}=\inf_{\Phi}\|\Phi\|_{C^{k,\alpha}(\bar{\Omega})}$, where the infimum is taken over the set of all global extensions $\Phi$.
Equipped with this norm, $C^{k,\alpha}(\partial\Omega)$ becomes a Banach space. In this work we shall
not use boundary norms for functions on curved boundaries but instead we
generally consider a boundary function as the restriction of a globally defined
function with its appropriate norm.

\noindent \hspace*{2em}In obtaining boundary estimates for $Lu=f$ in domains with a $C^{2,\alpha}(\alpha>0)$
boundary portion we first establish such an estimate in domains with a hyperplane


% PDF page 107
\source{gilbarg-trudinger:page-0107}



boundary portion. For this purpose we use the following interpolation inequality,
analogous to (6.8), (6.9). For its statement we recall the partially interior norms and
seminorms defined in (4.29).

\medskip

\noindent\textit{Let $\Omega$ be an open subset of $\mathbb{R}_+^n$ with a boundary portion $T$ on $x_n=0$, and assume
$u \in C^{2,\alpha}(\Omega \cup T)$. Then for any $\epsilon>0$ and some constant $C(\epsilon)$ we have}

\begin{equation}
(6.24)\qquad [u]^*_{j;\beta;\Omega \cup T}\le C|u|_{0;\Omega}+\epsilon [u]^*_{2,\alpha;\Omega \cup T},
\end{equation}

\[
j=0,1,2;\qquad 0\le \alpha,\beta\le 1,\ j+\beta<2+\alpha;
\]

\begin{equation}
(6.25)\qquad |u|^*_{j;\beta;\Omega \cup T}\le C|u|_{0;\Omega}+\epsilon [u]^*_{2,\alpha;\Omega \cup T}.
\end{equation}

\noindent These inequalities are proved in Lemma 6.34 of Appendix 1 to this chapter.
We can now assert the following essential local boundary estimate.

\medskip

\noindent\textbf{Lemma 6.4.} \quad \textit{Let $\Omega$ be an open subset of $\mathbb{R}_+^n$, with a boundary portion $T$ on $x_n=0$.
Suppose that $u \in C^{2,\alpha}(\Omega \cup T)$ is a bounded solution in $\Omega$ of $Lu=f$ satisfying the
boundary condition $u=0$ on $T$. In addition to (6.2) it is assumed that}

\begin{equation}
(6.26)\qquad |a^{ij}|^{(0)}_{0,\alpha;\Omega \cup T},\, |b^i|^{(1)}_{0,\alpha;\Omega \cup T},\, |c|^{(2)}_{0,\alpha;\Omega \cup T}\le \Lambda;\quad |f|^{(2)}_{0,\alpha;\Omega \cup T}<\infty.
\end{equation}

\noindent \textit{Then}

\begin{equation}
(6.27)\qquad |u|^*_{2,x;\Omega \cup T}\le C\bigl(|u|_{0;\Omega}+|f|^{(2)}_{0,\alpha;\Omega \cup T}\bigr)
\end{equation}

\noindent \textit{where $C=C(n,\alpha,\lambda,\Lambda)$.}

\medskip

\noindent\textit{Proof.} \quad The proof is identical with that of Theorem 6.2 if in the latter $d_x$ is replaced
by $\bar d_x$ and Lemma 6.1(a) and the inequalities (6.8), (6.9) are replaced when necessary
by Lemma 6.1(b) and the inequalities (6.24), (6.25). $\square$

\medskip

\noindent This lemma provides a bound on the first and second derivatives of $u$ and
the H\"older coefficients of its second derivatives in any subset $\Omega' \subset \Omega$ for which
dist $(\Omega',\,\partial \Omega - T)>0$. In particular $\partial \Omega'$ may contain any portion of $T$ at a non-
zero distance from $\partial \Omega - T$.

\medskip

\noindent In order to extend the preceding lemma to domains with a curved boundary
portion, we introduce the relevant seminorms and norms, in obvious generaliza-
tion of (4.29). Let $\Omega$ be an open set in $\mathbb{R}^n$ with $C^{k,\alpha}$ boundary portion $T$. For
$x,y \in \Omega$ we set

\[
\bar d_x = \operatorname{dist}(x,\partial \Omega - T),\qquad \bar d_{x,y}=\min(\bar d_x,\bar d_y)
\]


% PDF page 108
\source{gilbarg-trudinger:page-0108}


and for functions $u \in C^{k,\alpha}(\Omega \cup T)$ we define the following quantities:

\begin{equation*}
(6.28)
\end{equation*}

\[
[u]^{\ast}_{k,0;\Omega\cup T}=[u]^{\ast}_{k,\Omega\cup T}
= \sup_{\substack{x \in \Omega\\|\beta|=k}} d_x^k |D^\beta u(x)|,\qquad k=0,1,2,\ldots;
\]

\[
[u]^{\ast}_{k,\alpha;\Omega\cup T}
= \sup_{\substack{x,y \in \Omega\\|\beta|=k}} d_{x,y}^{k+\alpha}\frac{|D^\beta u(x)-D^\beta u(y)|}{|x-y|^\alpha},
\qquad 0<\alpha\le 1;
\]

\[
|u|^{\ast}_{k,0;\Omega\cup T}=|u|^{\ast}_{k;\Omega\cup T}
=\sum_{j=0}^k [u]^{\ast}_{j;\Omega\cup T};
\]

\[
|u|^{\ast}_{k,\alpha;\Omega\cup T}=|u|^{\ast}_{k;\Omega\cup T}+[u]^{\ast}_{k,\alpha;\Omega\cup T};
\]

\[
|u|^{(k)}_{0,\alpha;\Omega\cup T}
= \sup_{x\in\Omega} d_x^k|u(x)|+\sup_{x,y\in\Omega} d_{x,y}^{k+\alpha}\frac{|u(x)-u(y)|}{|x-y|^\alpha}.
\]

\noindent When $T=\phi$ and $\Omega \cup T=\Omega$, these quantities reduce to the interior seminorms and
norms already defined in (4.17) and (4.18).
Let $\Omega$ be a bounded domain with $C^{k,\alpha}$ boundary portion $T$, $k\ge 1$, $0\le \alpha \le 1$.
Suppose that $\Omega \subset \subset D$, where $D$ is a domain that is mapped by a $C^{k,\alpha}$ diffeomor-
phism $\psi$ onto $D'$. Letting $\psi(\Omega)=\Omega'$ and $\psi(T)=T'$, we can define the quantities in
(6.28) with respect to $\Omega'$ and $T'$. If $x'=\psi(x)$, $y'=\psi(y)$, one sees that

\begin{equation}
(6.29)\qquad K^{-1}|x-y|\le |x'-y'|\le K|x-y|
\end{equation}

\noindent for all points $x,y\in\Omega$, where $K$ is a constant depending on $\psi$ and $\Omega$. Letting
$u(x)\to \tilde{u}(x')$ under the mapping $x\to x'$, we find after a calculation using (6.29): for
$0\le j\le k$, $0\le \beta \le 1$, $j+\beta\le k+\alpha$,

\[
K^{-1}|u(x)|_{j,\beta;\Omega} \le |\tilde{u}(x')|_{j,\beta;\Omega'}
\le K|u(x)|_{j,\beta;\Omega};
\]

\begin{equation}
(6.30)\qquad
K^{-1}|u(x)|^{\ast}_{j,\beta;\Omega\cup T} \le |\tilde{u}(x')|^{\ast}_{j,\beta;\Omega'\cup T'}
\le K|u(x)|^{\ast}_{j,\beta;\Omega\cup T};
\end{equation}

\[
K^{-1}|u(x)|^{(\sigma)}_{0,\beta;\Omega\cup T} \le |\tilde{u}(x')|^{(\sigma)}_{0,\beta;\Omega'\cup T'}
\le K|u(x)|^{(\sigma)}_{0,\beta;\Omega\cup T}.
\]

\noindent In these inequalities $K$ denotes constants depending on the mapping $\psi$ and the
domain $\Omega$.
Lemma 6.4 can now be applied together with (6.30) to obtain a local boundary
estimate for curved boundaries. For this it is convenient to use the global norms
(4.6).

\noindent \textbf{Lemma 6.5.} \textit{Let $\Omega$ be a $C^{2,\alpha}$ domain in $\R^n$, and let $u \in C^{2,\alpha}(\bar{\Omega})$ be a solution of}
\textit{$Lu=f$ in $\Omega$, $u=0$ on $\partial\Omega$, where $f \in C^{\alpha}(\bar{\Omega})$. It is assumed that the coefficients of $L$}
\textit{satisfy (6.2) and}

\begin{equation}
(6.31)\qquad |a^{ij}|_{0,x;\Omega}, |b^i|_{0,x;\Omega}, |c|_{0,x;\Omega}\le \Lambda.
\end{equation}


% PDF page 109
\source{gilbarg-trudinger:page-0109}

Then for some $\delta$ there is a ball $B = B_\delta(x_0)$ at each point $x_0 \in \partial \Omega$ such that

\begin{equation}
\tag{6.32}
|u|_{2,\alpha;B\cap\Omega} \le C\bigl(|u|_{0;\Omega} + |f|_{0,\alpha;\Omega}\bigr)
\end{equation}

\noindent where $C = C(n,\alpha,\lambda,\Lambda,\Omega)$.

Proof. By the definition of a $C^{2,\alpha}$ domain, at each point $x_0 \in \partial \Omega$ there is a neighbourhood $N$ of $x_0$ and a $C^{2,\alpha}$ diffeomorphism that straightens the boundary in $N$. Let $B_\rho(x_0)\subset N$ and set $B' = B_\rho(x_0)\cap \Omega$, $D'=\psi(B')$, $T=B_\rho(x_0)\cap \partial\Omega\subset \partial B'$ and $T'=\psi(T)\subset \partial D'$ ($T'$ is a hyperplane portion of $\partial D'$). Under the mapping $y=\psi(x)=(\psi_1(x),\ldots,\psi_n(x))$, let $\tilde u(y)=u(x)$ and $\tilde L\tilde u(y)=Lu(x)$, where

\begin{equation*}
\tilde L\tilde u \equiv \tilde a^{ij}D_{ij}\tilde u+\tilde b^iD_i\tilde u+\tilde c\,\tilde u=\tilde f(y)
\end{equation*}

\noindent and

\begin{equation*}
\tilde a^{ij}(y)=\frac{\partial\psi_i}{\partial x_r}\frac{\partial\psi_j}{\partial x_s}a^{rs}(x),\qquad
\tilde b^i(y)=\frac{\partial^2\psi_i}{\partial x_r\partial x_s}a^{rs}(x)+\frac{\partial\psi_i}{\partial x_r}b^r(x),
\end{equation*}

\begin{equation*}
\tilde c(y)=c(x),\qquad \tilde f(y)=f(x).
\end{equation*}

\noindent We observe that in $D'$

\begin{equation*}
\lambda |\xi|^2 \le \tilde a^{ij}\xi_i\xi_j \qquad \forall \xi \in \mathbf{R}^n,
\end{equation*}

\noindent where

\begin{equation}
\tag{6.33}
\tilde\lambda = \lambda/K
\end{equation}

\noindent for a suitable positive constant $K$ depending only on the mapping $\psi$ on $B'$. By virtue of (6.30) we have also (for appropriate choice of $K$ in (6.33))

\begin{equation}
\tag{6.34}
|\tilde a^{ij}|_{0,\alpha;D'},\,|\tilde b^i|_{0,\alpha;D'},\,|\tilde c|_{0,\alpha;D'}\le \tilde\Lambda=K\Lambda,\qquad |\tilde f|_{0,\alpha;D'}<\infty.
\end{equation}

\noindent Thus the conditions of Lemma 6.4 are satisfied for the equation $\tilde L\tilde u=\tilde f$ in $D'$ with the hyperplane portion $T'$. We can therefore assert

\begin{equation*}
|\tilde u|_{2,\alpha;D'\cup T'}^* \le C\bigl(|\tilde u|_{0;D'}+|\tilde f|^{(2)}_{0,\alpha;D'\cup T'}\bigr)
\end{equation*}

\noindent where the constant $C=C(n,\alpha,\tilde\lambda,\tilde\Lambda)$. It follows from (6.30) that

\begin{equation*}
|u|_{2,\alpha;B'\cup T}^* \le C\bigl(|u|_{0;B'}+|f|^{(2)}_{0,\alpha;B'\cup T}\bigr)\le C\bigl(|u|_{0;B'}+|f|_{0,\alpha;B'}\bigr)
\end{equation*}

\begin{equation*}
\le C\bigl(|u|_{0;\Omega}+|f|_{0,\alpha;\Omega}\bigr),
\end{equation*}

\noindent where $C$ now depends on $n,\alpha,\lambda,\Lambda$ and $B'$. Letting $B''=B_{\rho/2}(x_0)\cap \Omega$ and observing that

\begin{equation*}
\min\bigl(1,(\rho/2)^{2+\alpha}\bigr)|u|_{2,\alpha;B''}\le |u|_{2,\alpha;B'\cup T}^*,
\end{equation*}

% PDF page 110
\source{gilbarg-trudinger:page-0110}
we obtain
\[
\tag{6.35}
\lvert u\rvert_{2,\alpha;B''}\le C\bigl(\lvert u\rvert_{0;\Omega}+\lvert f\rvert_{0,\alpha;\Omega}\bigr).
\]

The radius \(\rho\) appearing in the previous estimates depends in general on the point \(x_0\in \partial\Omega\). Consider now the collection of balls \(B_{\rho_i/4}(x)\) for all \(x\in \partial\Omega\). A finite subset \(B_{\rho_i/4}(x_i)\), \(i=1,\dots,N\), of this collection covers \(\partial\Omega\). Let \(\delta=\min \rho_i/4\) be the minimum radius of the balls in this finite covering. We assert that for this \(\delta\) the conclusion of the lemma is true. Namely, let \(C_i\) be the constant in (6.35) corresponding to \(x_i\), and let \(C=\max C_i\). Consider any point \(x_0\in \partial\Omega\) and the ball \(B_\delta(x_0)\). For some \(i\) we must have \(x_0\in B_{\rho_i/4}(x_i)\) and hence \(B=B_\delta(x_0)\subset B_{\rho_i/2}(x_i)=B_i\). From (6.35) we obtain the desired conclusion
\[
\lvert u\rvert_{2,\alpha;B\cap\Omega}\le \lvert u\rvert_{2,\alpha;B_i\cap\Omega}\le C\bigl(\lvert u\rvert_{0,\Omega}+\lvert f\rvert_{0,\alpha;\Omega}\bigr)
\]
where \(C\) depends on \(n\), \(\alpha\), \(\lambda\), \(\Lambda\), and \(\Omega\). \(\square\)

We remark that in the preceding lemma the dependence of the constant \(C\) on the domain \(\Omega\) is through the constant \(K\) in (6.30), (6.33) and (6.34), and \(K\) in turn depends only on the \(C^{2,\alpha}\) bounds on the mapping \(\psi\) which defines the local representation of the boundary \(\partial\Omega\). If the bounds on the mapping \(\psi\) can be stated uniformly over the boundary (which is always possible for \(C^{2,\alpha}\) domains), then \(K\) can replace \(\Omega\) in the statement of the estimate (6.32) and the domain \(\Omega\) may also be unbounded.

The principal result of this section is the following apriori global estimate for solutions with \(C^{2,\alpha}\) boundary values defined on \(C^{2,\alpha}\) domains.

\textbf{Theorem 6.6.} \emph{Let \(\Omega\) be a \(C^{2,\alpha}\) domain in \(\mathbf{R}^n\) and let \(u\in C^{2,\alpha}(\overline{\Omega})\) be a solution of \(Lu=f\) in \(\Omega\), where \(f\in C^\alpha(\overline{\Omega})\) and the coefficients of \(L\) satisfy, for positive constants \(\lambda\), \(\Lambda\),}
\[
a^{ij}\xi_i\xi_j\ge \lambda\lvert\xi\rvert^2 \qquad \forall x\in\Omega,\ \xi\in\mathbf{R}^n,
\]
\emph{and}
\[
\lvert a^{ij}\rvert_{0,\alpha;\Omega},\ \lvert b^i\rvert_{0,\alpha;\Omega},\ \lvert c\rvert_{0,\alpha;\Omega}\le \Lambda.
\]

\emph{Let \(\varphi(x)\in C^{2,\alpha}(\overline{\Omega})\), and suppose \(u=\varphi\) on \(\partial\Omega\). Then}
\[
\tag{6.36}
\lvert u\rvert_{2,\alpha;\Omega}\le C\bigl(\lvert u\rvert_{0;\Omega}+\lvert \varphi\rvert_{2,\alpha;\Omega}+\lvert f\rvert_{0,\alpha;\Omega}\bigr)
\]
\emph{where \(C=C(n,\alpha,\lambda,\Lambda,\Omega)\).}

\emph{Proof.} It suffices to prove the theorem for the case \(u=0\) on \(\partial\Omega\) and \(\varphi=0\). Namely, if we set \(v=u-\varphi\), then \(v=0\) on \(\partial\Omega\) and \(Lv=f-L\varphi\equiv f'\in C^\alpha(\overline{\Omega})\). The conclusion (6.36), applied to \(v\) with \(\varphi=0\), asserts
\[
\lvert v\rvert_{2,\alpha;\Omega}\le C\bigl(\lvert v\rvert_{0;\Omega}+\lvert f'\rvert_{0,\alpha;\Omega}\bigr).
\]

% PDF page 111
\source{gilbarg-trudinger:page-0111}



Since $|L\varphi|_{0,x;\Omega}\le C|\varphi|_{2,x;\Omega}$, it follows

\[
|u|_{2,\alpha;\Omega}\le |v|_{2,\alpha;\Omega}+|\varphi|_{2,\alpha;\Omega}\le C\bigl(|u|_{0;\Omega}+|\varphi|_{2,\alpha;\Omega}+|f|_{0,\alpha;\Omega}\bigr)
\]

as asserted in the theorem. In the remainder of the proof we assume $u=0$ on $\partial\Omega$.
Let $x \in \Omega$. We consider the two possibilities: (i) $x \in B_0 = B_{2\sigma}(x_0)\cap \Omega$ for some
$x_0 \in \partial\Omega$, where $\delta = 2\sigma$ is the radius in Lemma 6.5; (ii) $x \in \Omega_\sigma = \{x \in \Omega \mid$
dist $(x,\partial\Omega)>\sigma\}$. In case (i) Lemma 6.5 implies that

\begin{equation}
\tag{6.37}
|Du(x)|+|D^2u(x)|\le C(|u|_{0}+|f|_{0,\alpha}).
\end{equation}

\noindent (Here and in the following we omit the subscript $\Omega$ without ambiguity.) In case (ii)
we obtain the same inequality, with a different constant $C$, from Corollary 6.3 after
setting $d=\sigma$ in (6.23). Choosing the larger of the two constants $C$, we may assume
(6.37) for any point $x$ in $\Omega$, and hence for $|u|_2$.
Now let $x,y$ be distinct points in $\Omega$ and consider the three possibilities: (i) $x,y \in B_0$
for some $x_0$; (ii) $x,y \in \Omega_\sigma$; (iii) $x$ or $y$ is in $\Omega-\Omega_\sigma$, but not both $x$ and $y$ are in the
same ball $B_0$ for any $x_0$. These exhaust all possibilities. We consider the Hölder
quotient $|D^2u(x)-D^2u(y)|/|x-y|^\alpha$. In case (i), Lemma 6.5 gives the inequality

\[
\frac{|D^2u(x)-D^2u(y)|}{|x-y|^\alpha}\le C_1(|u|_{0}+|f|_{0,\alpha}).
\]

\noindent In case (ii) we obtain the same inequality, with a different constant $C_2$, from
Corollary 6.3. In case (iii), dist $(x,y)>\sigma$, so that

\[
\frac{|D^2u(x)-D^2u(y)|}{|x-y|^\alpha}\le \sigma^{-\alpha}(|D^2u(x)|+|D^2u(y)|)
\]

\[
\le C_3(|u|_{0}+|f|_{0,\alpha})\quad \text{by (6.37)}.
\]

\noindent Letting $C=\max(C_1,C_2,C_3)$, and taking the supremum over all $x,y \in \Omega$, we
obtain

\[
[D^2u]_\alpha\le C(|u|_{0}+|f|_{0,\alpha}).
\]

\noindent Combining this result with the bound for $|u|_2$ given by (6.37), we get

\[
|u|_{2,\alpha}\le C(|u|_{0}+|f|_{0,\alpha}),
\]

\noindent which establishes the theorem. $\square$

\medskip

\noindent \textit{Remark.} \quad The typical application of Theorem 6.6 concerns a set of solutions of an
equation or family of equations, each solution satisfying a uniform estimate (6.36).
The boundedness of the set of solutions in $C^{2,\alpha}(\overline{\Omega})$ then assures their precompactness
in $C^2(\overline{\Omega})$; (see Lemma 6.36).


% PDF page 112
\source{gilbarg-trudinger:page-0112}

A simple modification of the argument in Theorem 6.6 yields the following local estimate in domains with a $C^{2,\alpha}$ boundary portion.

\begin{corollary}[6.7]
Let $\Omega \subset \R^n$ be a domain with a $C^{2,\alpha}$ boundary portion $T \subset \partial\Omega$. Let $u \in C^{2,\alpha}(\Omega \cup T)$ be a solution of $Lu=f$ in $\Omega$, $u=\varphi$ on $T$, where $L$ and $f$ satisfy the conditions of Theorem 6.6, and $\varphi \in C^{2,\alpha}(\bar{\Omega})$. Then, if $x_0 \in T$ and $B = B_\rho(x_0)$ is a ball with radius $\rho < \dist(x_0,\partial\Omega - T)$, we have
\end{corollary}

\begin{equation}
\tag{6.38}
\lvert u \rvert_{2,\alpha;B \cap \Omega} \le C\bigl(\lvert u \rvert_{0;\Omega} + \lvert \varphi \rvert_{2,\alpha;\Omega} + \lvert f \rvert_{0,\alpha;\Omega}\bigr),
\end{equation}

where $C = C(n,\alpha,\lambda,\Lambda,B \cap \Omega)$.

It is a straightforward matter to extend the estimates of this and the preceding section to equations whose coefficients and inhomogeneous term are in $C^{k,\alpha}$, where $k > 0$; (see Problems 6.1, 6.2).

\section*{6.3. The Dirichlet Problem}

We consider the Dirichlet problem for $Lu=f$ in a bounded domain $\Omega$ of $\R^n$. Our procedure for solving this variable coefficient equation is to reduce it to the case of constant coefficients by the method of continuity (Theorem 5.2). Briefly summarized, this method as applied here starts with the solution for Poisson's equation $\Delta u=f$ and then arrives at a solution for $Lu=f$ through solutions for a continuous family of equations connecting $\Delta u=f$ and $Lu=f$.

We treat first the Dirichlet problem for sufficiently smooth domains and boundary values. In this case the connection between the solvability of Poisson's equation and of $Lu=f$ when $c\le 0$ is contained in the following theorem.

\begin{theorem}[6.8]
Let $\Omega$ be a $C^{2,\alpha}$ domain in $\R^n$, and let the operator $L$ be strictly elliptic in $\Omega$ with coefficients in $C^{\alpha}(\bar{\Omega})$ and with $c\le 0$. Then if the Dirichlet problem for Poisson's equation, $\Delta u=f$ in $\Omega$, $u=\varphi$ on $\partial\Omega$, has a $C^{2,\alpha}(\bar{\Omega})$ solution for all $f \in C^{\alpha}(\bar{\Omega})$ and all $\varphi \in C^{2,\alpha}(\bar{\Omega})$, the problem,
\end{theorem}

\begin{equation}
\tag{6.39}
Lu=f \text{ in } \Omega,\qquad u=\varphi \text{ on } \partial\Omega,
\end{equation}

also has a (unique) $C^{2,\alpha}(\bar{\Omega})$ solution for all such $f$ and $\varphi$.

\begin{proof}
By hypothesis we may assume that the coefficients of $L$ satisfy the conditions
\end{proof}

\begin{equation}
\tag{6.40}
\lambda \lvert \xi\rvert^2 \le a^{ij}\xi_i\xi_j \qquad \forall x\in\Omega,\ \xi\in\R^n,
\end{equation}

\begin{equation}
\lvert a^{ij}\rvert_{0,\alpha}, \lvert b^i\rvert_{0,\alpha}, \lvert c\rvert_{0,\alpha} \le \Lambda,
\end{equation}

% PDF page 113
\source{gilbarg-trudinger:page-0113}



with positive constants $\lambda, \Lambda$. (In writing norms, we omit the subscript $\Omega$, which will be implicitly understood.) It suffices to restrict consideration to zero boundary values, since the problem (6.39) is equivalent to $Lu=f-L\varphi=f'$ in $\Omega$, $u=0$ on $\partial\Omega$.
We consider the family of equations,

\begin{equation}
\tag{6.41}
L_tu \equiv tLu + (1-t)\Delta u = f, \qquad 0 \le t \le 1.
\end{equation}

We note that $L_0 = \Delta$, $L_1 = L$, and that the coefficients of $L_t$ satisfy (6.40) with

\[
\lambda_t = \min(1,\lambda), \qquad \Lambda_t = \max(1,\Lambda).
\]

The operator $L_t$ may be considered a bounded linear operator from the Banach space $\mathbb{B}_1 = \{u \in C^{2,\alpha}(\bar{\Omega}) \mid u = 0 \text{ on } \partial\Omega\}$ into the Banach space $\mathbb{B}_2 = C^\alpha(\bar{\Omega})$. The solvability of the Dirichlet problem, $L_tu=f$ in $\Omega$, $u=0$ on $\partial\Omega$, for arbitrary $f \in C^\alpha(\bar{\Omega})$ is then equivalent to the invertibility of the mapping $L_t$. Let $u_t$ denote a solution of this problem. By virtue of Theorem 3.7, we have the bound

\[
\lvert u_t \rvert_{0} \le C \sup_{\Omega} \lvert f \rvert \le C \lvert f \rvert_{0,\alpha},
\]

where $C$ depends only on $\lambda, \Lambda$ and the diameter of $\Omega$. Hence from (6.36), we have

\begin{equation}
\tag{6.42}
\lvert u_t \rvert_{2,\alpha} \le C \lvert f \rvert_{0,\alpha},
\end{equation}

that is,

\[
\lVert u \rVert_{\mathbb{B}_1} \le C \lVert L_tu \rVert_{\mathbb{B}_2},
\]

the constant $C$ being independent of $t$. Since, by hypothesis, $L_0 = \Delta$ maps $\mathbb{B}_1$ onto $\mathbb{B}_2$, the method of continuity (Theorem 5.2) is applicable and the theorem follows. $\square$

The preceding theorem presupposes the solvability in $C^{2,\alpha}(\bar{\Omega})$ of the Dirichlet problem for Poisson's equation when $\Omega$ and the boundary values are in class $C^{2,\alpha}$. Although this result--Kellogg's theorem--can be established independently by potential theoretic methods, we shall not assume it or prove it here, but rather derive it later as a consequence of elliptic theory. However, in the special case that $\Omega$ is a ball, Kellogg's theorem has already been proved, in Corollary 4.14. This provides the following existence theorem for balls.

\begin{quote}
\textbf{Corollary 6.9.} \textit{In Theorem 6.8 let $\Omega$ be a ball $B$, and let the operator $L$ satisfy the same conditions as in that theorem. Then if $f \in C^\alpha(\bar{B})$ and $\varphi \in C^{2,\alpha}(\bar{B})$, the Dirichlet problem, $Lu=f$ in $B$, $u=\varphi$ on $\partial B$, has a (unique) solution $u \in C^{2,\alpha}(\bar{B})$.}
\end{quote}

The conditions on the boundary data can be weakened to give the following generalization, which will be useful later.


% PDF page 114
\source{gilbarg-trudinger:page-0114}

\begin{lemma}[6.10]
Let $T$ be a boundary portion (possibly empty) of a ball $B$ in $\mathbf{R}^n$, and let $\varphi \in C^0(\partial B) \cap C^{2,\alpha}(T)$. Then if $L$ satisfies the conditions of Theorem 6.8 in $B$ and $f \in C^\alpha(\bar B)$, the Dirichlet problem, $Lu=f$ in $B$, $u=\varphi$ on $\partial B$, has a (unique) solution $u \in C^{2,\alpha}(B \cup T) \cap C^0(\bar B)$.
\end{lemma}

\begin{proof}
Let $x_0 \in T$ if $T$ is non-empty. We can assume that the boundary function $\varphi$ is continued by radial extension to a function $\varphi \in C^0(B') \cap C^{2,\alpha}(\bar G)$, where $B'$ is a ball containing $\bar B$ and $G=B_{\rho}(x_0) \subset B'$; (see Remark 2 after Lemma 6.38). Let $\{\varphi_k\}$ be a sequence of sufficiently smooth (say $C^3$) functions in $B'$ such that

\begin{equation}
\tag{6.43}
|\varphi_k-\varphi|_{0;B} \to 0 \quad \text{and} \quad |\varphi_k|_{2,\alpha;G} \leq C|\varphi|_{2,\alpha;G} \qquad \text{as } k \to \infty,
\end{equation}

where $C$ is a constant independent of $k$. (For the existence of such an approximation see the discussion after Lemma 7.1.) For each $k$ let $u_k$ be the corresponding solution of the Dirichlet problem, $Lu=f$ in $B$, $u=\varphi_k$ on $\partial B$. By Corollary 6.9 the functions $u_k$ are known to exist in $C^{2,\alpha}(\bar B)$. From the maximum principle it follows that the sequence $\{u_k\}$ converges uniformly to a function $u \in C^0(\bar B)$ such that $u=\varphi$ on $\partial B$. The compactness of $\{u_k\}$ provided by Corollary 6.3 assures that this sequence converges to a solution of $Lu=f$ on compact subsets of $B$ and hence the limit function $u$ is a solution in $B$ lying in $C^0(B)$. Furthermore, by Corollary 6.7, in $D=B_{\rho/2}(x_0) \cap B$ the functions $u_k$ satisfy the estimate

\[
|u_k|_{2,\alpha;D} \leq C\bigl(|u_k|_{0;B} + |\varphi_k|_{2,\alpha;G} + |f|_{0,\alpha;B}\bigr).
\]

It follows from (6.43) and Arzela's theorem that $u$ satisfies such an estimate in $D$ (with $\varphi$ replacing $\varphi_k$) and in particular that $u \in C^{2,\alpha}(\bar D)$. Thus $u \in C^{2,\alpha}(B \cup T)$ and the lemma is proved. $\square$
\end{proof}

We note especially that the preceding lemma provides a solution of the Dirichlet problem in balls when the boundary values are only continuous; the solution is then in class $C^0(\bar B) \cap C^{2,\alpha}(B)$.

It is now possible to imitate the Perron method of subharmonic functions of Chapter 2 and to extend the results obtained there for harmonic functions to the Dirichlet problem for $Lu=f$. A function $u \in C^0(\Omega)$ will be called a subsolution (supersolution) of $Lu=f$ in $\Omega$ if for every ball $B \Subset \Omega$ and every solution $v$ such that $Lv=f$ in $B$, the inequality $u \leq v$ ($u \geq v$) on $\partial B$ implies also $u \leq v$ ($u \geq v$) in $B$. If we suppose that $L$ satisfies the strong maximum principle and that the Dirichlet problem for $Lu=f$ is solvable in balls for continuous boundary values, then subsolutions and subharmonic functions are seen to share many of the same properties. In particular, we assert the following propositions without proof, the details being essentially the same as for subharmonic functions. When $f$ and the coefficients of $L$ are in $C^\alpha(\Omega)$, with $c \leq 0$, we need only observe that in the proofs Theorem 3.5 and Lemma 6.10 replace Theorems 2.2 and 2.6, respectively.

% PDF page 115
\source{gilbarg-trudinger:page-0115}


\section*{6.3. The Dirichlet Problem}

(i) A function $u \in C^2(\Omega)$ is a subsolution if and only if $Lu \ge f$.
(ii) If $u$ is a subsolution in a bounded domain $\Omega$ and $v$ is a supersolution such that
$v \ge u$ on $\partial\Omega$, then either $v > u$ throughout $\Omega$ or $v \equiv u$.
(iii) Let $u$ be a subsolution in $\Omega$ and $B$ a ball such that $\overline{B} \subset \Omega$. We denote by $\bar{u}$
the solution of $Lu = f$ in $B$ satisfying the condition $\bar{u} = u$ on $\partial B$. Then the function
$U$ defined by

\[
U(x) =
\begin{cases}
\bar{u}(x), & x \in B \\
u(x), & x \in \Omega - B
\end{cases}
\]

is a subsolution in $\Omega$.
(iv) Let $u_1, u_2, \ldots, u_N$ be subsolutions in $\Omega$, then the function
$u(x) = \max\{u_1(x), u_2(x), \ldots, u_N(x)\}$ is also a subsolution in $\Omega$.

In (i), (iii), and (iv) corresponding statements can obviously be made for supersolutions.

Now let $\Omega$ be a bounded domain and $\varphi$ a bounded function on $\partial\Omega$. A function
$u \in C^0(\bar{\Omega})$ will be called a subfunction (superfunction) relative to $\varphi$ if $u$ is a subsolution
(supersolution) in $\Omega$ and $u \le \varphi$ ($u \ge \varphi$) on $\partial\Omega$. By (ii) above every subfunction is less
than or equal to every superfunction. We denote by $S_\varphi$ the set of subfunctions in
$\Omega$ relative to $\varphi$. Let us assume that $S_\varphi$ is non-empty and is bounded from above.
This is the case, for example, when $L$ is strictly elliptic in $\Omega$ and its coefficients and
$f$ are bounded. Namely, if $\Omega$ lies in the slab $0 < x_1 < d$, then the functions

\begin{equation}
\tag{6.44}
\begin{aligned}
v^+ &= \sup |\varphi| + (e^{\gamma d} - e^{\gamma x_1}) \frac{\sup |f|}{\lambda} \\
v^- &= -\sup |\varphi| - (e^{\gamma d} - e^{\gamma x_1}) \frac{\sup |f|}{\lambda}
\end{aligned}
\end{equation}

are respectively superfunction and subfunction if the positive constant $\gamma$ is sufficiently
large; (see Theorem 3.7). The superfunction $v^+$ provides an upper bound for the
functions of $S_\varphi$, and the existence of the subfunction $v^-$ assures that $S_\varphi$ is
non-empty.

We can now assert the basic existence result of the Perron process for $Lu = f$,
assuming that $f$ and the coefficients of $L$ are in $C^\alpha(\Omega)$ and that $c \le 0$.

\begin{theorem}[6.11]
The function $u(x) = \sup_{v \in S_\varphi} v(x)$ lies in $C^{2,\alpha}(\Omega)$ and satisfies $Lu = f$
in $\Omega$ provided $u$ is bounded.
\end{theorem}

The proof differs only in minor details from that of Theorem 2.12 and is left to
the reader. We call attention to the fact that the compactness of solutions required
in the proof is provided by the interior estimate of Corollary 6.3, and that the
appropriate form of the maximum principle in the argument is given by Theorem 3.5.


% PDF page 116
\source{gilbarg-trudinger:page-0116}

We wish now to determine conditions that the solution defined in Theorem 6.11 assumes the boundary values $\varphi$ continuously. As in the case of harmonic functions, this problem can be treated by means of the barrier concept, which we define for $Lu=f$ with $c\le 0$ in a bounded domain $\Omega$. Let $\varphi$ be a bounded function on $\partial\Omega$, continuous at $x_0\in \partial\Omega$. Then a sequence of functions $\{w_i^+(x)\}$ ($\{w_i^-(x)\}$) in $C^0(\bar{\Omega})$ is an upper (lower) barrier in $\Omega$ relative to $L$, $f$, and $\varphi$ at $x_0$ if:

\begin{enumerate}
\item[(i)] $w_i^+ (w_i^-)$ is a super(sub)function relative to $\varphi$ in $\Omega$;
\item[(ii)] $w_i^\pm(x_0)\to \varphi(x_0)$ as $i\to \infty$.
\end{enumerate}

If both an upper and lower barrier exist at a point, it will be convenient to speak simply of a barrier at the point.

The basic property of barriers is contained in the following:

\begin{lemma}[6.12]
Let $u$ be the solution of $Lu=f$ in $\Omega$ defined by Theorem 6.11, where $\varphi$ is a bounded function on $\partial\Omega$, continuous at $x_0$. If there exists a barrier at $x_0$, then $u(x)\to \varphi(x_0)$ as $x\to x_0$.
\end{lemma}

\begin{proof}
From the definition of $u$ and the fact that every subfunction is dominated by every superfunction, we have for all $i$,

\[
w_i^-(x)\le u(x)\le w_i^+(x) \qquad \text{in } \Omega.
\]

For any $\varepsilon>0$ and all sufficiently large $i$, condition (ii) above implies

\[
\lim_{x\to x_0} w_i^-(x)>\varphi(x_0)-\varepsilon, \qquad \lim_{x\to x_0} w_i^+(x)<\varphi(x_0)+\varepsilon.
\]

It follows that

\[
\limsup_{x\to x_0} |u(x)-\varphi(x_0)|<\varepsilon,
\]

and hence we conclude $u(x)\to \varphi(x_0)$ as $x\to x_0$. $\square$
\end{proof}

We amplify on the barrier concept with the following remarks.

\begin{remark}[1]
In many cases of interest the special structure of the equation simplifies the determination of a barrier. Thus, if $c\equiv 0$, $f\equiv 0$ in the equation $Lu=f$, the situation is the same as for Laplace's equation in that a barrier at $x_0$ is determined by a single supersolution $w\in C^0(\bar{\Omega})$ with the property that $w>0$ on $\partial\Omega-x_0$, $w(x_0)=0$. To see this, let $\varepsilon>0$; then by the boundedness of $\varphi$ and its continuity at $x_0$, there is a positive constant $k_\varepsilon$ such that
\end{remark}

\[
w_\varepsilon^+\equiv \varphi(x_0)+\varepsilon+k_\varepsilon w, \qquad w_\varepsilon^-\equiv \varphi(x_0)-\varepsilon-k_\varepsilon w
\]

are respectively superfunction and subfunction in $\Omega$ with respect to $\varphi$, and obviously $w_\varepsilon^\pm(x_0)\to \varphi(x_0)$ as $\varepsilon\to 0$. Thus the family $w_\varepsilon^\pm(x)$ determines a barrier.

% PDF page 117
\source{gilbarg-trudinger:page-0117}

\setcounter{section}{6}
\setcounter{subsection}{2}
\noindent\textbf{6.3.\ The Dirichlet Problem}\hfill 105

\medskip

\noindent To consider another class of equations, let $f$ and the coefficients of $L$ be bounded in $\Omega$. Again a single function $w \in C^0(\overline{\Omega}) \cap C^2(\Omega)$ determines a barrier at $x_0$ if it satisfies the conditions: (a) $Lw \leq -1$ in $\Omega$; (b) $w > 0$ on $\partial\Omega - x_0$, $w(x_0)=0$. Namely, given $\varepsilon > 0$, then as above there is a positive constant $k_\varepsilon$ such that

\[
\varphi(x_0)+\varepsilon+k_\varepsilon w(x)\geq \varphi(x), \qquad
\varphi(x_0)-\varepsilon-k_\varepsilon w(x)\leq \varphi(x) \quad \text{on } \partial\Omega.
\]

\noindent If now we set $k_\varepsilon'=\max\bigl(k_\varepsilon,\sup_\Omega |f-c\varphi(x_0)|\bigr)$, the functions

\[
w_\varepsilon^+ \equiv \varphi(x_0)+\varepsilon+k_\varepsilon' w, \qquad
w_\varepsilon^- \equiv \varphi(x_0)-\varepsilon-k_\varepsilon' w
\]

\noindent are respectively superfunction and subfunction and define a barrier at $x_0$. In fact, $w_\varepsilon^+ \geq \varphi$, $w_\varepsilon^- \leq \varphi$ on $\partial\Omega$, and since $Lw \leq -1$,

\[
L[\varphi(x_0)+\varepsilon+k_\varepsilon' w]\leq c(\varphi(x_0)+\varepsilon)-k_\varepsilon' \leq c\varphi(x_0)-k_\varepsilon' \leq f;
\]

\noindent similarly, $L[\varphi(x_0)-\varepsilon-k_\varepsilon' w]\geq f$, and hence the family $w_\varepsilon^\pm$ determines a barrier with respect to $\varphi$ at $x_0$.

\noindent When, as in the above cases, a barrier at $x_0$ is constructed from a fixed super-solution $w$ of $Lu=0$ depending only on $L$ and the domain, we shall say that $w$ determines a barrier at $x_0$.

\medskip

\noindent\textit{Remark 2.} \quad The above definition of barrier is often difficult to apply, since it requires the construction of global sub- and supersolutions defined over all of $\Omega$. It may then become necessary to find \textit{local barriers} to achieve the desired results. To motivate the definition of this concept, let $M^+(M^-)$ be an upper (lower) bound for a solution in $\Omega$ whose boundary behavior is being studied at a point $x_0 \in \partial\Omega$. Then a sequence of functions $\{w_i^+(x)\}$ ($\{w_i^-(x)\}$) is a \textit{local upper (lower) barrier} relative to $L,f,\varphi$ and $M^+(M^-)$ at $x_0$ if there is an open neighborhood $\mathcal N$ of $x_0$ such that:

\begin{enumerate}
\item[(i)] $w_i^+ (w_i^-)$ is a super (sub) solution in $\mathcal N \cap \Omega$;
\item[(ii)] $w_i^+ \geq \varphi \quad (w_i^- \leq \varphi)$ on $\mathcal N \cap \partial\Omega$;
\item[(iii)] $w_i^+ \geq M^+ \quad (w_i^- \leq M^-)$ on $\Omega \cap \partial\mathcal N$;
\item[(iv)] $w_i^\pm(x_0) \to \varphi(x_0)$ as $i \to \infty$.
\end{enumerate}

\noindent (In particular, if $\overline{\Omega}\subset \mathcal N$, the functions $w_i^\pm$ define a barrier at $x_0$ in the previous global sense and condition (iii) can be omitted.) One sees immediately that if the solution $u$ defined in Theorem 6.11 satisfies $|u|\leq M$ in $\Omega$, then Lemma 6.12 remains valid if there exists a local barrier at $x_0$ with respect to the bounds $\pm M$. Later in this work local barriers will play an important part in the study of boundary behavior of solutions.

\medskip

\noindent\textit{Remark 3.} \quad The same argument as in Lemma 6.12 shows that a barrier determines a modulus of continuity at the boundary for any solution assuming its boundary values continuously. Thus in the cases considered in Remark 1, if $u$ is a bounded


% PDF page 118
\source{gilbarg-trudinger:page-0118}
\begin{quote}
solution of $Lu=f$ such that $u(x)\to\varphi(x_{0})$ as $x\to x_{0}$, we have for any $\varepsilon>0$ and a suitable positive constant $k_{\varepsilon}$
\end{quote}

\[
|u(x)-\varphi(x_{0})|\leq \varepsilon+k_{\varepsilon}w(x)\qquad \text{in }\Omega.
\]

If $w$ determines a local barrier at $x_{0}$ the same inequality holds in a fixed neighborhood of $x_{0}$ (independent of $\varepsilon$).

For the equation $Lu=f$, as for Laplace's equation, the existence and construction of barriers is closely connected to the local properties of the boundary. To illustrate by an example of interest in later application, let $L$ be strictly elliptic in a bounded domain $\Omega$, with $c\leq 0$, and let $f$ and the coefficients of $L$ be bounded. We suppose that $\Omega$ satisfies an exterior sphere condition at $x_{0}\in\partial\Omega$, so that $\overline{B}\cap\overline{\Omega}=x_{0}$ for some ball $B=B_{R}(y)$. We show that the function

\begin{equation}
w(x)=\tau\left(R^{-\sigma}-r^{-\sigma}\right),\qquad r=|x-y|,
\tag{6.45}
\end{equation}

satisfies $Lw\leq -1$ in $\Omega$ for suitable positive constants $\tau$ and $\sigma$ and hence (by Remark 1 above) $w$ determines a barrier at $x_{0}$. Taking $y=0$ for convenience, we have from (6.40) and the fact that $c\leq 0$, for $x\in\Omega$

\[
L(R^{-\sigma}-r^{-\sigma})\leq \sigma r^{-\sigma-4}\left[-(\sigma+2)a^{ij}x_{i}x_{j}+r^{2}\left(\sum a^{ii}+b^{i}x_{i}\right)\right]
\]
\[
\leq \sigma r^{-\sigma-2}\left[-(\sigma+2)\lambda+\left(\sum a^{ii}+b^{i}x_{i}\right)\right].
\]

Since $\Omega$ and the coefficients of $L$ are bounded, the right member is negative and bounded away from zero provided $\sigma$ is sufficiently large. Hence, for suitably large $\tau$ and $\sigma$, $Lw\leq -1$, as asserted.

Thus, under the above assumptions on the equation $Lu=f$, a bounded domain $\Omega$ satisfying an exterior sphere condition at every point $x_{0}\in\partial\Omega$ (e.g. any $C^{2}$ domain), has a barrier at every boundary point, and Lemma 6.12 is applicable whenever the prescribed boundary values are continuous. Combining this fact with Theorem 6.11, we are led to the following general existence theorem.

\textbf{Theorem 6.13.} \textit{Let $L$ be strictly elliptic in a bounded domain $\Omega$, with $c\leq 0$, and let $f$ and the coefficients of $L$ be bounded and belong to $C^{\alpha}(\Omega)$. Suppose that $\Omega$ satisfies an exterior sphere condition at every boundary point. Then, if $\varphi$ is continuous on $\partial\Omega$, the Dirichlet problem,}

\[
Lu=f\text{ in }\Omega,\qquad u=\varphi\text{ on }\partial\Omega,
\]

\textit{has a (unique) solution $u\in C^{0}(\overline{\Omega})\cap C^{2,\alpha}(\Omega)$.}

The preceding theorem can be extended to more general domains. In particular, under the same hypotheses on $L$ and $f$, it can be proved for domains satisfying an exterior cone condition (see Problem 6.3).

% PDF page 119
\source{gilbarg-trudinger:page-0119}

If the hypotheses of this theorem are strengthened so that $f$ and the coefficients
of $L$ are in $C^\alpha(\overline{\Omega})$, it can be shown that the domains for which the Dirichlet problem
is solvable for continuous boundary values are precisely the same for both the
Laplace operator and $L$ (see Notes).

We turn now to the question of global regularity of the above solutions when
the boundary data are sufficiently smooth. We have seen that under the hypotheses
of Theorem 6.8 the solution of the Dirichlet problem for $Lu=f$ lies in $C^{2,\alpha}(\overline{\Omega})$
provided the same result (Kellogg's theorem) is true for Poisson's equation. We proceed now to prove this regularity theorem directly from the results of this section.

\textbf{Theorem 6.14.} \textit{Let $L$ be strictly elliptic in a bounded domain $\Omega$, with $c\leq 0$, and
let $f$ and the coefficients of $L$ belong to $C^\alpha(\overline{\Omega})$. Suppose that $\Omega$ is a $C^{2,\alpha}$ domain
and that $\varphi\in C^{2,\alpha}(\overline{\Omega})$. Then the Dirichlet problem,}

\[
Lu=f \text{ in } \Omega, \qquad u=\varphi \text{ on } \partial\Omega,
\]

\textit{has a (unique) solution lying in $C^{2,\alpha}(\overline{\Omega})$.}

\textit{Proof.} Since the hypotheses of Theorem 6.13 are satisfied, let $u$ be the corresponding solution of the Dirichlet problem. We know from Theorem 6.13 that
$u\in C^0(\overline{\Omega})\cap C^{2,\alpha}(\Omega)$, so it remains only to show that at each point $x_0\in \partial\Omega$, we have
$u\in C^{2,\alpha}(D\cap \overline{\Omega})$, where $D$ is some neighbourhood of $x_0$. Since $\Omega$ is a $C^{2,\alpha}$ domain,
there is a neighborhood $N$ of $x_0$ that can be taken by a $C^{2,\alpha}$ mapping $y=\psi(x)$ with
$C^{2,\alpha}$ inverse into a neighborhood $\widetilde N$ in such a way that $\psi(N\cap \overline{\Omega})$ contains the
closure of a ball $B$, and a portion $T$ of $N\cap \partial\Omega$ containing $x_0$ is mapped by $\psi$ into a
boundary portion $\widetilde T$ of $B$. Under this mapping the equation $Lu(x)=f(x)$ transforms
into an equation $\widetilde L\widetilde u(y)=\widetilde f(y)$ defined in $B$. Because of the $C^{2,\alpha}$ character of the
mapping, we have $\varphi\to\widetilde\varphi\in C^{2,\alpha}(\overline{B})$, and $\widetilde L$ and $\widetilde f$ satisfy the same hypotheses in
$B$ as $L$ and $f$ do in $\Omega$; that is, $\widetilde L$ is strictly elliptic in $B$, with $\widetilde c\leq 0$, and $\widetilde f$ and the
coefficients of $\widetilde L$ are in $C^\alpha(\overline{B})$; (see Lemma 6.5.). Consider then the solution $v$ of the
Dirichlet problem, $\widetilde L v=\widetilde f$ in $B$, $v=\widetilde u$ on $\partial B$. Since $\widetilde u=\widetilde\varphi$ on $\widetilde T\subset \partial B$, we have
that $\widetilde u\in C^0(\partial B)\cap C^{2,\alpha}(\widetilde T)$ on $\partial B$. Hence, by uniqueness for the Dirichlet
problem and by Lemma 6.10 it follows $\widetilde u=v\in C^0(\overline{B})\cap C^{2,\alpha}(B\cup \widetilde T)$. Returning
to $\Omega$, let $D'=\psi^{-1}(B)$. We see that $u\in C^{2,\alpha}(D'\cup T)$, and since $x_0$ was arbitrary on
$\partial\Omega$, we conclude that $u\in C^{2,\alpha}(\overline{\Omega})$. \qed

The preceding result can be extended to conditions of lower regularity of the
coefficients, domain and boundary values (see Notes).

If the operator $L$ does not satisfy the condition $c\leq 0$, then, as is well known
from simple examples, the Dirichlet problem for $Lu=f$ no longer has a solution
in general. However, it is still possible to assert a Fredholm alternative, which we
formulate as follows.

\textbf{Theorem 6.15.} \textit{Let $L\equiv a^{ij}D_{ij}+b^iD_i+c$ be strictly elliptic with coefficients in
$C^\alpha(\overline{\Omega})$ in a $C^{2,\alpha}$ domain $\Omega$. Then either: (a) the homogeneous problem, $Lu=0$ in
$\Omega$, $u=0$ on $\partial\Omega$, has only the trivial solution, in which case the inhomogeneous problem,
}


% PDF page 120
\source{gilbarg-trudinger:page-0120}
Lu = f in $\Omega$, $u = \varphi$ in $\partial\Omega$, has a unique $C^{2,\alpha}(\overline{\Omega})$ solution for all $f \in C^\alpha(\overline{\Omega})$, $\varphi \in C^{2,\alpha}(\overline{\Omega})$;
or (b) the homogeneous problem has nontrivial solutions, which form a finite dimensional subspace of $C^{2,\alpha}(\overline{\Omega})$.

Proof. We have seen that under the stated assumptions concerning $f$ and $\varphi$, the inhomogeneous problem, $Lu = f$ in $\Omega$, $u = \varphi$ on $\partial\Omega$, is equivalent to the problem $Lv = f - L\varphi$, $v = 0$ on $\partial\Omega$. We shall therefore consider the Dirichlet problem with only the homogeneous boundary condition, $u = 0$ on $\partial\Omega$, and it will suffice to restrict the operator $L$ to the linear space

\[
\B = \{u \in C^{2,\alpha}(\overline{\Omega}) \mid u = 0 \text{ on } \partial\Omega\}.
\]

Let $\sigma$ be any constant such that $\sigma \ge \sup_\Omega c$, and define the operator $L_\sigma = L - \sigma$.
Then, by Theorem 6.14, the mapping

\[
L_\sigma : \B \to C^\alpha(\overline{\Omega})
\]

is invertible. Furthermore, the inverse mapping $L_\sigma^{-1}$, by the estimate (6.36) and Theorem 3.7, is a compact mapping from $C^\alpha(\overline{\Omega})$ into $C^{2}(\overline{\Omega})$ and hence is also compact as a mapping from $C^\alpha(\overline{\Omega})$ into $C^\alpha(\overline{\Omega})$. Consider then the equation

\begin{equation}
u + \sigma L_\sigma^{-1}u = L_\sigma^{-1}f, \qquad f \in C^\alpha(\overline{\Omega}),
\tag{6.46}
\end{equation}

in which $L_\sigma^{-1} : C^\alpha(\overline{\Omega}) \to C^\alpha(\overline{\Omega})$ is compact. By Theorem 5.3 on compact operators in a Banach space, the alternative applies and (6.46) always has a solution $u \in C^\alpha(\overline{\Omega})$ provided the homogeneous equation $u + \sigma L_\sigma^{-1}u = 0$ has only the trivial solution $u = 0$. When this condition is not satisfied, the null space of the operator $I + \sigma L_\sigma^{-1}$ ($I$ = identity) is a finite dimensional subspace of $C^\alpha(\overline{\Omega})$ (Theorem 5.5).

To interpret these statements in terms of the Dirichlet problem for $Lu = f$, we observe first that since $L_\sigma^{-1}$ maps $C^\alpha(\overline{\Omega})$ onto $\B$, any solution $u \in C^\alpha(\overline{\Omega})$ of (6.46) must also belong to $\B$. Hence, operating on (6.46) with $L_\sigma$, we obtain

\begin{equation}
Lu = L_\sigma(u + \sigma L_\sigma^{-1}u) = f, \qquad u \in \B.
\tag{6.47}
\end{equation}

Thus, the solutions of (6.46) are in one-to-one correspondence with the solutions of the boundary value problem (6.47), and we can therefore conclude the alternative stated in the theorem. $\square$

The importance of the alternative for the Dirichlet problem is that it shows uniqueness to be a sufficient condition for existence. We remark that by virtue of Lemma 6.18 (to be proved in the next section), the $C^{2}(\overline{\Omega})$ solutions of $Lu = f$ are also in $C^{2,\alpha}(\overline{\Omega})$, and hence the null space of $L$ in $C^{2}(\overline{\Omega})$ is also finite dimensional.
We note also that Theorem 5.5 implies that the set $\Sigma$ of real values $\sigma$ for which the homogeneous problem, $Lu - \sigma u = 0$, $u = 0$ on $\partial\Omega$, has nontrivial solutions is at most countable and discrete. Furthermore (by Theorem 5.3), if $\sigma \notin \Sigma$, any

% PDF page 121
\source{gilbarg-trudinger:page-0121}

solution of the Dirichlet problem $Lu=f$ in $\Omega$, $u=\varphi$ on $\partial\Omega$, satisfies an estimate
\[
|u|_{2,\alpha}\leq C(|\varphi|_{2,\alpha}+|f|_{0,\alpha})
\]
where the constant $C$ is independent of $u$, $f$ and $\varphi$.

\subsection*{6.4. Interior and Boundary Regularity}

In the preceding sections the inhomogeneous term $f$ and the coefficients of $L$ were assumed to be $C^\alpha$ functions and the corresponding solutions of the equation $Lu=f$ were in class $C^{2,\alpha}$. We now investigate the higher regularity properties of the solutions in their dependence on the smoothness of $f$ and the coefficients of $L$. The global regularity of the solutions will depend as well on the smoothness of the boundary and the boundary values.

First we show that any $C^2$ solution of $Lu=f$ must also be in class $C^{2,\alpha}$ if $f$ and the coefficients of $L$ are in $C^\alpha$.

\textbf{Lemma 6.16.} \emph{Let $u$ be a $C^2(\Omega)$ solution of the equation $Lu=f$ in an open set $\Omega$, where $f$ and the coefficients of the elliptic operator $L$ are in $C^\alpha(\Omega)$. Then $u\in C^{2,\alpha}(\Omega)$.}

\emph{Proof.} It obviously suffices to prove $u\in C^{2,\alpha}(B)$ for arbitrary balls $B\subset\subset\Omega$. Let $B$ be such a ball and in $B$ consider the Dirichlet problem for $v$:

\[
(6.48)\qquad L_0v=a^{ij}D_{ij}v+b^iD_iv=f' \equiv f-cu,\qquad v=u\ \text{on}\ \partial B.
\]

Since, by hypothesis, $u\in C^2(\bar B)$, we have that $f'$ and the coefficients of $L_0$ are in $C^\alpha(\bar B)$. Hence, by Lemma 6.10, there is a solution $v$ of (6.48) lying in $C^{2,\alpha}(B)\cap C^0(\bar B)$. By uniqueness we infer that the solutions $u$ and $v$ of (6.48) are identical in $B$ and thus $u\in C^{2,\alpha}(B)$. $\square$

We note that the preceding result and those that follow in this section make no assumption concerning the sign of the coefficient $c$.

The above lemma and the Schauder interior estimates yield the following interior regularity theorem.

\textbf{Theorem 6.17.} \emph{Let $u$ be a $C^2(\Omega)$ solution of the equation $Lu=f$ in an open set $\Omega$, where $f$ and the coefficients of the elliptic operator $L$ are in $C^{k,\alpha}(\Omega)$. Then $u\in C^{k+2,\alpha}(\Omega)$. If $f$ and the coefficients of $L$ lie in $C^\infty(\Omega)$, then $u\in C^\infty(\Omega)$.}

\emph{Proof.} By Lemma 6.16 the theorem is established for $k=0$. We now prove it for $k=1$. Let $v$ be a function on $\Omega$ and denote by $e_1\ (l=1,\ldots,n)$ the unit coordinate vector in the $x_1$ direction. We define the difference quotient of $v$ at $x$ in the direction $e_1$ by
\[
\Delta_1^h v(x)=\Delta_1^h v(x)=\frac{v(x+he_1)-v(x)}{h}.
\]

% PDF page 122
\source{gilbarg-trudinger:page-0122}
\begin{flushleft}
Taking the difference quotients of both sides of the equation
\end{flushleft}

\[
Lu = a^{ij}D_{ij}u + b^iD_i u + cu = f
\]

\begin{flushleft}
we obtain
\end{flushleft}

\begin{equation}
\tag{6.49}
L(\Delta^h u) = a^{ij}D_{ij}\Delta^h u + b^iD_i \Delta^h u + c\,\Delta^h u = F_h(x)
\end{equation}
\[
\equiv \Delta^h f - (\Delta^h a^{ij})D_{ij}\bar{u} - (\Delta^h b^i)D_i\bar{u} - (\Delta^h c)\bar{u},\qquad \bar{u}=u(x+he_l).
\]

\begin{flushleft}
All the difference quotients in this equation are assumed taken at $x\in\Omega$ in the direction $e_l$ for some $l=1,\ldots,n$. Since $f\in C^{1,\alpha}(\Omega)$ and
\end{flushleft}

\[
\Delta^h f(x)=\frac{1}{h}\int_0^1 \frac{d}{dt}f(x+the_l)\,dt
= \int_0^1 D_l f(x+the_l)\,dt.
\]

\begin{flushleft}
one sees that $\Delta^h f\in C^\alpha(\Omega')$ in any subset $\Omega' \Subset \Omega$ for which $|h|<\operatorname{dist}(\Omega',\partial\Omega)$. In particular, if $B$ and $B'$ are balls in $\Omega$ such that $B'\Subset B\Subset\Omega$ and $\operatorname{dist}(B',\partial B)=h_0>0$, then $\Delta^h f\in C^\alpha(\overline{B'})$ for $0<|h|<h_0$, and there is a uniform bound: $|\Delta^h f|_{0,\alpha;B'}\leq \operatorname{const}$ independent of $h$. Similar bounds hold for the difference quotients $\Delta^h a^{ij}$, $\Delta^h b^i$, $\Delta^h c$, which also belong to $C^\alpha(\overline{B'})$. Since $u\in C^{2,\alpha}(\Omega)$ (by Lemma 6.16), it follows that $F_h\in C^\alpha(\overline{B'})$ for $|h|<h_0$ and, furthermore, $|F_h|_{0,\alpha;B'}\leq \operatorname{const}$ independent of $h$. Observing also the bound, $\sup_{B'}|\Delta^h u|\leq \sup_B |Du|$, we can infer from the interior estimates of Corollary 6.3 that the set of functions $\Delta^h u$ and their first and second derivatives $D_i\Delta^h u$, $D_{ij}\Delta^h u$, $(i,j=1,\ldots,n)$, are bounded and equicontinuous on any ball $B''\Subset B'$, and hence every sequence in these sets contains a uniformly convergent subsequence on $B'$. Since $\Delta^h u\to D_lu$ as $h\to0$, we may therefore assert that $D_{ij}\Delta^h u\to D_{ijl}u$ as $h\to0$ and that $u\in C^{3,\alpha}(B'')$. Since $B''$ could be an arbitrary ball having its closure in $\Omega$, we infer that $u\in C^{3,\alpha}(\Omega)$, thereby proving the theorem for $k=1$.
\end{flushleft}

\begin{flushleft}
To prove the theorem for $k>1$, we proceed by induction on $k$. Under the stated hypotheses on $L$ and $f$ we may accordingly assume that $u\in C^{k+1,\alpha}(\Omega)$, and we wish to show that u $\in C^{k+2,\alpha}(\Omega)$. Since $f$ and the coefficients of $L$ are in $C^{k,\alpha}(\Omega)$, the equation $Lu=f$ may be differentiated $k-1$ times, and this yields an equation $L\hat{u}=\hat{f}$, where $\hat{u}=D^\beta u$ for some index $\beta$ such that $|\beta|=k-1$, and where $\hat{f}$ equals $D^\beta f$ plus a sum of terms which are products of derivatives of the coefficients of order $\leq k-1$ and of derivatives of $u$ of order $\leq k$. Thus $\hat{f}\in C^{1,\alpha}(\Omega)$. By the same argument as above for $k=1$, we see that $\hat{u}\in C^{3,\alpha}(\Omega)$ and hence $u\in C^{k+2,\alpha}(\Omega)$, as claimed in the theorem. The final assertion concerning solutions in $C^\infty(\Omega)$ follows immediately. $\square$
\end{flushleft}

\begin{flushleft}
It is also the case that if $f$ and the coefficients of $L$ are real analytic functions, then any solution of $Lu=f$ is likewise analytic. For the proof we refer to the literature (e.g. [HO 3]).
\end{flushleft}

\begin{flushleft}
To be able to assert analogous regularity properties up to the boundary, it is obviously necessary that the boundary itself and the boundary values of the
\end{flushleft}

% PDF page 123
\source{gilbarg-trudinger:page-0123}
\begin{flushleft}
solution be sufficiently smooth. To establish the appropriate results on regularity
up to the boundary we first prove an analogue of Lemma 6.16.
\end{flushleft}

\begin{flushleft}
\textbf{Lemma 6.18.} \emph{Let $\Omega$ be a domain with a $C^{2,\alpha}$ boundary portion $T$, and let $\varphi \in C^{2,\alpha}(\overline{\Omega})$.
Suppose that $u$ is a $C^0(\overline{\Omega}) \cap C^2(\Omega)$ function satisfying $Lu = f$ in $\Omega$, $u = \varphi$ on $T$,
where $f$ and the coefficients of the strictly elliptic operator $L$ belong to $C^\alpha(\overline{\Omega})$. Then
$u \in C^{2,\alpha}(\Omega \cup T)$.}
\end{flushleft}

\begin{flushleft}
\textbf{Proof.} Since $T$ is a $C^{2,\alpha}$ portion of $\partial \Omega$, it is possible to find at each point $x_0$ of
$T$ a boundary neighborhood $T' \subset \subset T$ and a $C^{2,\alpha}$ domain $D$ in $\Omega$ such that $x_0 \in
T' \subset \partial D$. In addition, $D$ can be chosen so small that Corollary 3.8 is applicable
and hence the Dirichlet problem for $Lu = f$ in $D$ has at most one solution in
$C^0(\overline{D}) \cap C^2(D)$.
\end{flushleft}

\begin{flushleft}
The argument is now very similar to that in Lemma 6.10. Since $u \in$
$C^0(\partial D) \cap C^{2,\alpha}(T')$, we can extend the boundary values of $u$ on $\partial D$ to a function
$v \in C^0(\overline{D}') \cap C^{2,\alpha}(\overline{B})$, where $D \subset \subset D'$ and $B = B_\rho(x_0) \subset D'$. (See Remark 1 after
Lemma 6.38 for the construction of $v$.) Let $\{v_k\}$ be a sequence of functions in
$C^3(D')$ such that
\end{flushleft}

\[
|v_k - v|_{0;D} \to 0 \quad \text{and} \quad |v_k|_{2,\alpha;B} \le C|v|_{2,\alpha;B} \quad \text{as } k \to \infty.
\]

\begin{flushleft}
By virtue of the Fredholm alternative (Theorem 6.15), the Dirichlet problem,
$Lu = f$ in $D$, $u = v_k$ on $\partial D$, has a unique solution $u_k$ in $C^{2,\alpha}(\overline{D})$ for each $k$. As a
consequence of Corollaries 6.3 and 3.8, the sequence $\{u_k\}$ converges to the solution
$u$ in $D$; and by Corollary 6.7 we have $u \in C^{2,\alpha}(\overline{B}' \cap \overline{D})$, where $B' = B_{\rho/2}(x_0)$. Since
$x_0$ was an arbitrary point on $T$, it follows that $u \in C^{2,\alpha}(\Omega \cup T)$. $\square$
\end{flushleft}

\begin{flushleft}
If $c \le 0$, the preceding result is essentially contained in the proof of Theorem 6.14.
However, when there are no restrictions on the sign of the coefficient $c$, the argu-
ment in Theorem 6.14 has to be suitably modified, as in the above proof. Another
proof, along different lines, is contained in the Notes.
\end{flushleft}

\begin{flushleft}
Lemma 6.18 remains valid if $C^{2,\alpha}$ is replaced by $C^{1,\alpha}$ in the hypotheses and
conclusion, and under more general circumstances as well (see [GH]).
\end{flushleft}

\begin{flushleft}
With Lemma 6.18 as a starting point we can establish the following global
regularity theorem.
\end{flushleft}

\begin{flushleft}
\textbf{Theorem 6.19.} \emph{Let $\Omega$ be a $C^{k+2,\alpha}$ domain $(k \geq 0)$ and let $\varphi \in C^{k+2,\alpha}(\overline{\Omega})$. Suppose
that $u$ is a $C^0(\overline{\Omega}) \cap C^2(\Omega)$ function satisfying $Lu = f$ in $\Omega$, $u = \varphi$ on $\partial \Omega$, where $f$ and
the coefficients of the strictly elliptic operator $L$ belong to $C^{k,\alpha}(\overline{\Omega})$. Then $u \in C^{k+2,\alpha}(\overline{\Omega})$.}
\end{flushleft}

\begin{flushleft}
\textbf{Proof.} For $k = 0$ the theorem is implied by Lemma 6.18. We now prove the result
for $k = 1$. Let $x_0$ be an arbitrary boundary point of $\Omega$, and consider a suitable
$C^{3,\alpha}$ diffeomorphism $\psi$ that straightens the boundary near $x_0$. As a consequence of
this mapping we may consider the equation $Lu = f$ as defined in a domain $G$ with
a hyperplane portion $T$ on $x_n = 0$, while the remaining hypotheses of the theorem
\end{flushleft}

% PDF page 124
\source{gilbarg-trudinger:page-0124}

are unchanged (see p. 91). By considering the function $u-\varphi$ in place of $u$ and noting that $L\varphi \in C^{1,\alpha}(\bar{G})$, we may assume that $\varphi = 0$ in the statement of the theorem.

As in Theorem 6.17, we take the difference quotient of the equation $Lu=f$ in the direction $e_l$ for any $l=1,\ldots,n-1$, obtaining thereby an equation of the form (6.49), which is satisfied by the difference quotient $\Delta^h u$ ($=\Delta_l^h u$). If $0<|h|<h_0$, then this equation is valid in the set $G'=\{x\in G\mid \dist(x,\partial G-T)>h_0\}$, which has a hyperplane boundary portion $T' \subset T$. Under the assumptions on $L$ and $f$, and since $u=0$ on $T$ and $u\in C^{2,\alpha}(\bar{G})$, the conditions of Lemma 6.4 are satisfied in $G'$ by equation (6.49) and the solution $\Delta^h u$. It follows from Lemma 6.4 that the function families $\Delta^h u$, $D_i\Delta^h u$, $D_{ij}\Delta^h u$, $(i,j=1,\ldots,n)$, are bounded and equicontinuous on compact subsets of $G' \cup T'$. Since $\Delta^h u \to D_lu$ as $h\to0$, we can assert that $D_{ij}\Delta^h u \to D_{ijl}u$ for $i,j=1,\ldots,n$ and $l=1,\ldots,n-1$, and, in addition, $D_lu \in C^{2,\alpha}(G' \cup T')$ for $l=1,\ldots,n-1$. It remains only to show that $D_nu \in C^{2,\alpha}(G' \cup T')$ as well. This follows at once by writing

\[
D_{nn}u=(1/a^{nn})(f-(L-a^{nn}D_{nn})u)
\]

and observing from the preceding results that the right hand side is contained in $C^{1,\alpha}(G' \cup T')$. Since $x_0$ was an arbitrary point on $\partial\Omega$, we conclude that $u\in C^{3,\alpha}(\bar{\Omega})$.

The proof of the theorem for $k>1$ proceeds by induction on $k$, as in Theorem 6.17, by considering the equation satisfied by any derivative of order $k-1$ and thereby reducing to the case $k=1$ treated above. $\square$

It is clear from the preceding argument, which is essentially local, that the regularity result remains true up to any $C^{k+2,\alpha}$ boundary portion $T$ provided the solution $u$ is continuous up to $T$ and takes on $C^{k+2,\alpha}$ boundary values on $T$.

\section*{6.5. An Alternative Approach}

An examination of the proof of Theorem 6.13 shows that this existence theorem follows by the Perron process from the solvability of the Dirichlet problem in balls for arbitrary continuous boundary values. The latter result (contained in Lemma 6.10) depended in an essential way for its proof on the boundary estimates of the Schauder theory. However, as we shall see below, it is possible to develop a theory of the Dirichlet problem for continuous boundary values that is based entirely on the Schauder interior estimates, without any use of boundary estimates.

The developments of this section will rest on the following extension of the interior estimates in Theorem 6.2. For its statement we make use of the seminorms and norms defined in (6.10).

\begin{lemma}[6.20]
Let $u\in C^{2,\alpha}(\Omega)$ satisfy the equation $Lu=f$ in an open set $\Omega$ of $\R^n$, where the coefficients of $L$ satisfy (6.12) and (6.13). Suppose that $\lvert u\rvert_{0;\Omega}^{(-\beta)}<\infty$ and $\lvert f\rvert_{0,\alpha;\Omega}^{(2-\beta)}<\infty$ for some $\beta\in\R$. Then we have
\end{lemma}

\begin{equation}
\tag{6.50}
\lvert u\rvert_{2,\alpha;\Omega}^{(-\beta)}\leq C\bigl(\lvert u\rvert_{0;\Omega}^{(-\beta)}+\lvert f\rvert_{0,\alpha;\Omega}^{(2-\beta)}\bigr).
\end{equation}

where $C=C(n,\alpha,\lambda,\Lambda,\beta)$.

% PDF page 125
\source{gilbarg-trudinger:page-0125}

Proof. Let $x$ be any point in $\Omega$, $d_x$ its distance from $\partial\Omega$ and $d=d_x/2$. Then by
(6.14), applied to the ball $B=B_d(x)$, we have

\begin{equation*}
d^{1-\beta}\lvert Du(x)\rvert+d^{2-\beta}\lvert D^2u(x)\rvert\leq Cd^{-\beta}\bigl(\lvert u\rvert_{0;B}+\lvert f\rvert_{0,\alpha;B}^{(2)}\bigr)
\end{equation*}
\begin{equation*}
\leq C\left[\sup_{y\in B}d_y^{-\beta}\lvert u(y)\rvert+\sup_{y\in B}d_y^{2-\beta}\lvert f(y)\rvert\right.
\end{equation*}
\begin{equation*}
\left.+\sup_{y\in B}d_{x,y}^{2+\alpha-\beta}\frac{\lvert f(x)-f(y)\rvert}{\lvert x-y\rvert^{\alpha}}\right]
\end{equation*}
\begin{equation*}
\leq C\bigl(\lvert u\rvert_{0;\Omega}^{(-\beta)}+\lvert f\rvert_{0,\alpha;\Omega}^{(2-\beta)}\bigr).
\end{equation*}

Hence

\begin{equation}
(6.51)\qquad \lvert u\rvert_{2;\Omega}^{(-\beta)}\leq C\bigl(\lvert u\rvert_{0;\Omega}^{(-\beta)}+\lvert f\rvert_{0,\alpha;\Omega}^{(2-\beta)}\bigr).
\end{equation}

To estimate $[u]_{2,\alpha;\Omega}^{(-\beta)}$, let $x, y$ be distinct points in $\Omega$, with $d_x\leq d_y$ and
$B=B_d(x)$ as above. Then, by considering the two cases $\lvert x-y\rvert\leq d/2$ and $\lvert x-y\rvert>d/2$, we
have for any second derivative $D^2u$.

\begin{equation*}
d_{x,y}^{2+\alpha-\beta}\frac{\lvert D^2u(x)-D^2u(y)\rvert}{\lvert x-y\rvert^{\alpha}}
\leq Cd^{-\beta}\bigl(\lvert u\rvert_{0;B}+\lvert f\rvert_{0,\alpha;B}^{(2)}\bigr)
\end{equation*}
\begin{equation*}
+d_x^{2+\alpha-\beta}\frac{\lvert D^2u(x)\rvert+\lvert D^2u(y)\rvert}{(d/2)^{\alpha}}
\end{equation*}
\begin{equation*}
\leq C\bigl(\lvert u\rvert_{0;\Omega}^{(-\beta)}+\lvert f\rvert_{0,\alpha;\Omega}^{(2-\beta)}\bigr)+8[u]_{2;\Omega}^{(-\beta)}.
\end{equation*}

Taking the supremum with respect to $x, y$, and applying (6.51), we obtain

\begin{equation*}
[u]_{2,\alpha;\Omega}^{(-\beta)}\leq C\bigl(\lvert u\rvert_{0;\Omega}^{(-\beta)}+\lvert f\rvert_{0,\alpha;\Omega}^{(2-\beta)}\bigr).
\end{equation*}

Combining this with (6.51), we get the desired estimate (6.50). $\square$

We observe that for $\beta=0$ the preceding result reduces to the previous estimate
(6.14). When $\beta>0$, the hypothesis that $\lvert u\rvert_{0;\Omega}^{(-\beta)}$ is finite obviously requires that
$u=0$ on $\partial\Omega$.

In this section we shall solve the Dirichlet problem for $Lu=f$ in balls using the
method of continuity. This will require an apriori estimate of solutions for suitably
unbounded $f$, given by the following lemma. The corresponding result for Poisson's
equation is contained in Theorem 4.9.

\begin{lemma}[6.21]
Let $L$ be strictly elliptic (satisfying (6.2)), with $c\leq 0$ and with coefficients bounded in magnitude by $\Lambda$ in a ball $B=B_R(x_0)$. Suppose that $u\in C^0(\bar{B})\cap C^2(B)$ is a solution of $Lu=f$ in $B$, $u=0$ on $\partial B$. Then, for any $\beta\in(0,1)$ we have
\end{lemma}

\begin{equation}
(6.52)\qquad \sup_{x\in B}d_x^{-\beta}\lvert u(x)\rvert\leq C\sup_{x\in B}d_x^{2-\beta}\lvert f(x)\rvert,
\end{equation}

where $C=C(\beta,n,R,\lambda,\Lambda)$.

% PDF page 126
\source{gilbarg-trudinger:page-0126}

\begin{flushleft}
\emph{Proof.} Let $\beta$ be fixed and assume that $\sup_{x\in B} d_x^{2-\beta}\lvert f(x)\rvert = N < \infty$. The required
estimate (6.52) will be obtained by the construction of a suitable comparison
function bounding $u$. For convenience we take $x_0=0$, and let us set
\end{flushleft}

\[
w_1(x)=(R^2-r^2)^\beta,\qquad r=\lvert x\rvert.
\]

\begin{flushleft}
\textbf{Then}
\end{flushleft}

\[
Lw_1(x)=\beta(R^2-r^2)^{\beta-2}\bigl[4(\beta-1)a^{ij}x_i x_j-2(R^2-r^2)\bigl(\sum a^{ii}+b^i x_i\bigr)
+(c/\beta)(R^2-r^2)^2\bigr]
\]

\[
\le -\beta(R^2-r^2)^{\beta-2}\bigl[4(1-\beta)\lambda r^2+2(R^2-r^2)(n\lambda-\sqrt{n}\,\Lambda r)\bigr].
\]

\begin{flushleft}
It is clear that for some $R_0$, $0\le R_0<R$, the expression in brackets is positive if
$R_0\le r\le R$. Hence
\end{flushleft}

\[
Lw_1(x)\le -c_1(R-r)^{\beta-2}\qquad\text{in }R_0\le r<R,
\]

\[
\le c_2(R-r)^{\beta-2}\qquad\text{in }0\le r<R_0,
\]

\begin{flushleft}
where $c_1$ and $c_2$ are positive constants depending only on $\beta$, $n$, $R$, $\lambda$, and $\Lambda$. (If $R_0=0$,
the second inequality is of course superfluous.)
\end{flushleft}

\begin{flushleft}
Now let $w_2(x)=e^{\alpha R}-e^{\alpha x_1}$, where $\alpha\ge 1+\sup\lvert b\rvert/\lambda$. Then, as in Theorem 3.7,
we have $Lw_2(x)\le -\lambda e^{-\alpha R}$ in $B$, and hence
\end{flushleft}

\[
Lw_2(x)\le -c_3(R-r)^{\beta-2}\qquad\text{in }0\le r<R_0,
\]

\[
\le 0\qquad\text{in }R_0\le r<R,
\]

\begin{flushleft}
where $c_3=\lambda e^{-\alpha R}(R-R_0)^{2-\beta}$. Since, by assumption, $\lvert f(x)\rvert\le Nd_x^{\beta-2}$ and $d_x=R-r$,
it follows that for the positive constants $\gamma_1=1/c_1$ and $\gamma_2=(1+c_2/c_1)/c_3$,
\end{flushleft}

\[
L(\gamma_1w_1+\gamma_2w_2)\le -(R-r)^{\beta-2}\le -\lvert f(x)\rvert/N\qquad\text{in }0\le\lvert x\rvert<R.
\]

\begin{flushleft}
Letting $w=\gamma_1w_1+\gamma_2w_2$, one sees that $w(x)\ge 0$ on $\partial B$, $w(x)=0$ at $x=(R,0,\ldots,0)$,
and
\end{flushleft}

\[
L(Nw+u)\le 0\qquad\text{in }B,\qquad Nw+u\ge 0\qquad\text{on }\partial B.
\]

\begin{flushleft}
From the maximum principle (Corollary 3.2) we infer
\end{flushleft}

\[
\tag{6.53}
\lvert u(x)\rvert\le Nw(x)\qquad\text{in }B.
\]

\begin{flushleft}
Consider now any point $x\in B$, which we assume without loss of generality to
lie on the $x_1$ axis. Then (6.53) implies the inequality
\end{flushleft}

\[
\lvert u(x)\rvert\le CN(R-r)^\beta = CNd_x^\beta
\]

\begin{flushleft}
for some constant $C = C(\beta,n,R,\lambda,\Lambda)$, and the lemma is proved. $\square$
\end{flushleft}

% PDF page 127
\source{gilbarg-trudinger:page-0127}


\section*{6.5. An Alternative Approach}

The preceding result can be extended by similar methods to more general
domains, for example, to \(C^{2}\) domains (see Problem 6.5).

With the help of the preceding two lemmas we can now prove the following
extension of Theorem 4.9 to the equation \(Lu=f\). We observe that the proof makes
no use of boundary estimates.

\textbf{Theorem 6.22.} \emph{Let \(B\) be a ball in \(\mathbf{R}^{n}\) and \(f\) a function in
\(C^{\alpha}(B)\) such that \(|f|_{0,\alpha;B}^{(2-\beta)}<\infty\) for some \(\beta\in(0,1)\).
Assume \(L\) to be strictly elliptic in \(B\), with \(c\le 0\) and coefficients
satisfying (6.2) and (6.31). Then there exists a (unique) solution \(u\in C^{0}(\overline{B})\cap C^{2,\alpha}(B)\)
of the Dirichlet problem, \(Lu=f\) in \(B\), \(u=0\) on \(\partial B\). In addition, \(|u|_{0;B}^{(-\beta)}<\infty\) and hence
\(u\) satisfies the estimate (6.50) in \(B\).}

\medskip

\textit{Proof.} The argument is based on the method of continuity. As in Theorem 6.8,
we consider the family of equations,

\[
L_{t}u \equiv tLu+(1-t)\Delta u=f,\qquad 0\le t\le 1,
\]

and observe that the coefficients of \(L_{t}\) also satisfy (6.2) and (6.31) in \(B\), with
\(\lambda_{t}=\min(1,\hat{\lambda})\), \(\Lambda_{t}=\max(1,\Lambda)\) replacing \(\lambda\) and \(\Lambda\) respectively. By virtue of (6.11) we
have

\[
|a^{ij}D_{ij}u|_{0,\alpha;B}^{(2-\beta)}\,,\ |b^{i}D_{i}u|_{0,\alpha;B}^{(2-\beta)}\,,\ |cu|_{0,\alpha;B}^{(2-\beta)}\le C|u|_{2,\alpha;B}^{(-\beta)}\,,
\]

and hence the operator \(L_{t}\), for each \(t\), is a bounded linear operator from the Banach
space

\[
\B_{1}=\{u\in C^{2,\alpha}(B)\mid |u|_{2,\alpha;B}^{(-\beta)}<\infty\}
\]

into the Banach space

\[
\B_{2}=\{f\in C^{\alpha}(B)\mid |f|_{0,\alpha;B}^{(2-\beta)}<\infty\}.
\]

The solvability of the Dirichlet problem, \(Lu=f\) in \(B\), \(u=0\) on \(\partial B\), for arbitrary
\(f\in \B_{2}\) is equivalent to the invertibility of the mapping \(\B_{1}\to \B_{2}\) defined by \(u\to Lu\).
Let \(u_{t}\) denote a solution of this problem for some \(t\in[0,1]\). Then from (6.52) we have

\[
|u_{t}|_{0}^{(-\beta)}\le C_{1}|f|_{0}^{(2-\beta)}\le C_{1}|f|_{0,\alpha}^{(2-\beta)}.
\]

It follows from (6.50) that

\[
|u_{t}|_{2,\alpha}^{(-\beta)}\le C|f|_{0,\alpha}^{(2-\beta)},
\]

or, equivalently,

\[
\|u\|_{\B_{1}}\le C\|L_{t}u\|_{\B_{2}}.
\]

% PDF page 128
\source{gilbarg-trudinger:page-0128}

\section*{6. Classical Solutions; the Schauder Approach}

the constant $C$ being independent of $t$. The fact that $L_0$ is onto is already contained in Theorem 4.9. The method of continuity (Theorem 5.2) is now applicable and the theorem is proved. $\square$

The preceding theorem can be extended to yield the following generalization of Lemma 6.10 for the case of continuous boundary values.

\textbf{Corollary 6.23.} \textit{Under the hypotheses of Theorem 6.22, if $\varphi \in C^0(\overline{B})$ there exists a (unique) solution $u \in C^0(\overline{B}) \cap C^{2,\alpha}(B)$ of the Dirichlet problem, $Lu=f$ in $B$, $u=\varphi$ on $\partial B$.}

\textit{Proof.} Let $\{\varphi_k\}$ be a sequence of functions in $C^3(\overline{B})$ converging uniformly to $\varphi$ on $\overline{B}$. By Theorem 6.22, the Dirichlet problem, $Lv_k=f-L\varphi_k$ in $B$, $v_k=0$ on $\partial B$, is uniquely solvable for each $k$, and defines a solution $u_k=v_k+\varphi_k$ of the corresponding problem with inhomogeneous boundary values, $Lu_k=f$ in $B$, $u_k=\varphi_k$ on $\partial B$. It follows in the usual way from the maximum principle that $u_k$ converges uniformly on $\overline{B}$ to a function $u \in C^0(\overline{B})$ such that $u=\varphi$ on $\partial B$. From the compactness provided by the interior estimates (Corollary 6.3), it follows also that $Lu=f$ in $B$, and hence $u$ is the required solution. $\square$

Starting with this existence theorem in balls, we can proceed as before to apply the Perron process in more general domains, obtaining in particular Theorem 6.13.

\section*{6.6. Non-Uniformly Elliptic Equations}

The existence results of the preceding sections, which were valid in arbitrary smooth bounded domains, were derived under the assumption of uniform ellipticity of the differential operator $L$. When the equation ceases to be uniformly elliptic, the conditions for solvability are greatly circumscribed and will in general require limitations on the geometry of the domain or connections between the differential operator and the geometry.

It is instructive to consider an example for which the Dirichlet problem is not solvable. Let us consider solutions $u(x, y)$ of the equation

\begin{equation*}
(6.54) \qquad u_{xx}+y^2u_{yy}=0
\end{equation*}

in the rectangle $R$, $0<x<\pi$, $0<y<Y$, such that $u\in C^0(\overline{R})\cap C^2(R)$ satisfies the boundary conditions $u(0,y)=u(\pi,y)=0$, $0\le y\le Y$. Any such solution $u(x,y)$ has a Fourier series expansion $\sum f_n(y)\sin nx$, in which the coefficients $f_n(y)$ satisfy the ordinary differential equation $y^2 f_n''-n^2 f_n=0$. This equation has the independent solutions $y^{\beta_n}$, $y^{\gamma_n}$, where $\beta_n=\tfrac12\bigl(1+\sqrt{1+4n^2}\bigr)>0$, $\gamma_n=\tfrac12\bigl(1-\sqrt{1+4n^2}\bigr)<0$. The fact that $u(x, y)$ is bounded at $y=0$ requires that $f_n(y)=\mathrm{const}\, y^{\beta_n}$, and hence $f_n(0)=0$. It follows that $u(x,0)=0$, and accordingly the only continuous solutions on $\overline{R}$ satisfying the prescribed boundary conditions must have zero boundary values on $y=0$ and therefore cannot take on prescribed non-zero boundary data.

% PDF page 129
\source{gilbarg-trudinger:page-0129}

\section*{6.6. Non-Uniformly Elliptic Equations}

The arguments leading to Theorem 6.13 can be extended to non-uniformly elliptic equations under suitable hypotheses on the coefficients and the domain. We observe to begin with that if $f$ and the coefficients of $L$ in the equation $Lu=f$ $(c\le 0)$ are locally H\"older continuous in the domain $\Omega$, and if the set of subfunctions of the Dirichlet problem with respect to the prescribed boundary function $\varphi$ is non-empty and bounded above, then the Perron process defines a bounded solution of $Lu=f$ in $\Omega$ (Theorem 6.11). In particular, this is the case if $\lvert b\rvert/\lambda$, $f/\lambda$ and the domain $\Omega$ are bounded, where $\lambda(x)$ is the minimum eigenvalue of the coefficient matrix $A(x)=[a^{ij}(x)]$; (see Theorem 3.7). This assumption will be understood in the following considerations.

To study whether $u(x)\to \varphi(x_0)$ at a specific boundary point $x_0\in\partial\Omega$ where $\varphi$ is continuous, let us suppose, as in the discussion preceding Theorem 6.13, that $\Omega$ satisfies an exterior sphere condition at $x_0$, and let $B=B_R(y)$ be a ball such that $\overline{B}\cap\overline{\Omega}=x_0$. While no longer assuming that $L$ is uniformly elliptic near $x_0$, we make the additional (but less restrictive) hypothesis that $\lvert A(x)\cdot(x-y)\rvert\ge \delta>0$ for all $x$ in the intersection $N\cap\Omega$ of $\Omega$ with some neighborhood $N$ of $x_0$. (In particular, if the coefficient matrix $A(x)$ is continuous at $x_0$, this condition will be satisfied if $A(x_0)\cdot(x_0-y)\ne 0$; that is, if the normal vector to $\partial\Omega$ is not in the null space of $A$ at $x_0$.) It then follows that $a^{ij}(x_i-y_i)(x_j-y_j)\ge \lambda' \lvert x-y\rvert^2$ for all $x\in N\cap\Omega$, where $\lambda'$ is a suitable positive constant, although the minimum eigenvalue $\lambda$ may tend to zero at $x_0$. Let us assume also that all coefficients of $L$ are bounded. The barrier argument preceding Theorem 6.13 now proceeds as before and we conclude that the function $w(x)=\tau(R^{-\sigma}-r^{-\sigma})$ defined in (6.45) determines a local barrier at $x_0$ for appropriate choices of $\tau$ and $\sigma$, and hence $u(x)\to\varphi(x_0)$. We note that in the boundary value problem considered for equation (6.54), the normal vector at each point of the boundary segment, $y=0$, $0<x<\pi$, is in the null-space of the coefficient matrix $[1,0/0,0]$, and thus the above hypothesis is not fulfilled.

Alternatively, let $[a^{ij}(x)]$ be an arbitrary positive matrix and assume that the functions $\lvert b\rvert/\lambda$, $c/\lambda$, $f/\lambda$, are bounded. By dividing by the minimum eigenvalue $\lambda$ we may assume without loss of generality that the operator $L$ is strictly elliptic in $\Omega$ with $\lambda=1$. The limit behaviour $u(x)\to \varphi(x_0)$ is now guaranteed if $\Omega$ satisfies a strict exterior plane condition at $x_0$. By this we shall mean that in some neighborhood of $x_0$, there is a hyperplane intersecting $\overline{\Omega}$ in the single point $x_0$. Such a condition will be satisfied, for example, if $\partial\Omega$ is strictly convex near $x_0$. To prove the assertion $u(x)\to \varphi(x_0)$, let us for convenience choose $x_0$ to be the origin and let the normal to the assumed exterior plane at $x_0$ be the $x_1$ axis, with $x_1>0$ in $\Omega$ near $x_0$. Under the stated conditions there is a slab $0<x_1<d$ whose intersection $D$ with $\Omega$ near $x_0$ is such that $x_1>0$ on $\overline{D}-x_0$. As in the proof of Theorem 3.7 we see that the function

\[
w(x)=e^{\gamma d}(1-e^{-\gamma x_1})
\]

satisfies $Lw\le -\lambda$ in $D$ provided $\gamma\ge 1+\sup_D(\lvert b\rvert/\lambda)$, and hence $Lw\le -1$ if we take $\lambda=1$. It follows as in Remark 1 following Lemma 6.12 that for a suitable constant $k=k(\varepsilon)$, the functions

\[
w_\varepsilon^+\equiv \varphi(x_0)+\varepsilon+kw,\qquad
w_\varepsilon^-\equiv \varphi(x_0)-\varepsilon-kw
\]

% PDF page 130
\source{gilbarg-trudinger:page-0130}

determine a local barrier at $x_0$ with respect to the upper and lower bounds on $u$ in $\Omega$.
We observe that if the boundary function $\varphi$ is constant near $x_0$, then $u(x)\to \varphi(x_0)$ follows even if $\Omega$ satisfies a non-strict exterior plane condition at $x_0$. In this connection one notes that in the boundary value problem considered previously for (6.54), the boundary segment, $y=0$, $0<x<\pi$, is convex but not strictly convex, and the stated boundary value problem is solvable only for zero data on the interval.

The preceding remarks immediately yield the following simple extension of Theorem 6.13 to non-uniformly elliptic equations.

\medskip

\noindent\textbf{Theorem 6.24.} \textit{Let $L$ be strictly elliptic (satisfying (6.2)) in a bounded domain $\Omega$, with $c\leq 0$, $a^{ij}$, $b^i$, $c$, $f\in C^\alpha(\Omega)$, and assume $b^i$, $c$, $f$ are bounded. Suppose that $\Omega$ satisfies the exterior sphere condition and, in addition, a strict exterior plane condition at those boundary points where any of the coefficients $a^{ij}$ are unbounded. Then, if $\varphi$ is continuous on $\partial\Omega$, the Dirichlet problem, $Lu=f$ in $\Omega$, $u=\varphi$ on $\partial\Omega$, has a (unique) solution $u\in C^0(\overline{\Omega})\cap C^{2,\alpha}(\Omega)$.}

\medskip

It is clear from the above arguments that this result can be modified in various ways (see, for example, Problem 6.4). When the equation is homogeneous and the lower order terms of $L$ are not present, we obtain the following:

\medskip

\noindent\textbf{Corollary 6.24'.} \textit{Let the coefficients $a^{ij}$ of the elliptic equation $a^{ij}D_{ij}u=0$ belong to $C^\alpha(\Omega)$, where $\Omega$ is a bounded strictly convex domain. Then the Dirichlet problem is solvable in $C^0(\overline{\Omega})\cap C^{2,\alpha}(\Omega)$ for arbitrary continuous boundary values.}

\medskip

Although this result is an immediate consequence of Theorem 6.24 its proof can be obtained more directly by observing that because of the strict convexity of $\Omega$ and the special form of the equation, a linear function determines a barrier at each boundary point.

A closer examination of the relation between the coefficients of $L$ and the local curvature properties of the boundary makes it possible to derive other general sufficient conditions for the existence of barriers. Let $f$ and the coefficients of $L$ (with $c\leq 0$) be bounded and assume $\Omega$ is a $C^2$ domain. The minimum eigenvalue of the principal coefficient matrix $[a^{ij}(x)]$ may approach zero on $\partial\Omega$. We shall seek conditions for the existence of a local barrier at $x_0\in \partial\Omega$ with respect to the continuous boundary function $\varphi$ and the bound $M$.

Let $B$ denote a ball centered at $x_0$ and set $G=B\cap\Omega$; $B$ will be specified later.
Let $\psi\in C^2(\overline{G})$ be a fixed function such that $\psi(x)>0$ on $\overline{G}-x_0$ and $\psi(x_0)=0$.
For every $\varepsilon>0$ and a suitable constant $k=k(\varepsilon)$, we can satisfy the inequalities

\[
\varepsilon + k\psi(x) \geq |\varphi(x)-\varphi(x_0)| \qquad \text{on } \partial\Omega\cap B,
\]
\[
\geq M \qquad \text{on } \partial B\cap \Omega.
\]

Let us now define the distance function (see Appendix to Chapter 14).

\[
d(x)=\dist(x,\partial\Omega), \quad x\in\Omega,
\]

% PDF page 131
\source{gilbarg-trudinger:page-0131}

which is in class $C^2$ in some neighborhood

\[
\mathcal{N} = \{x \in \Omega \mid d(x) < d_0\}.
\]

It can be assumed that $G \subset \mathcal{N}$. We propose to find conditions under which functions
$w^\pm(x)$ of the form

\[
w^+ = \varphi(x_0) + \varepsilon + k\psi + Kd, \qquad w^- = \varphi(x_0) - \varepsilon - k\psi - Kd
\]

define a barrier at $x_0$ in $G$, where $K = K(\varepsilon)$ is an appropriate positive constant. This
will be the case if $Lw^+ \leq f$ and $Lw^- \geq f$ in $G$. Thus, the functions $w^\pm$ define a barrier
if the inequality

\[
(6.55) \qquad K(a^{ij}D_{ij}d + b^iD_i d) \leq -k|L\psi| - |f - c\varphi(x_0)|
\]

holds in $G$ for some $K$. This provides a sufficient condition for the existence of a barrier which, in
principle, can be verified by inspection of the given equation and domain.

To realize the condition (6.55) more concretely, let us, for example, assume the coefficients
$a^{ij}(x)$ to be continuous at $x_0$. Choose a coordinate system with origin at $x_0$ and the
$x_n$ axis coincident with the inner normal $Dd$ at $x_0$. A rotation $P$ of the coordinate system
about the $x_n$ axis, in which the new axes are the principal directions of $\partial\Omega$ at $x_0$,
diagonalizes the Hessian matrix $[D_{ij}d(x_0)]$, so that

\[
P^{\!\top}[D_{ij}d(x_0)]P = \diag[-\kappa_1, -\kappa_2, \ldots, -\kappa_{n-1}, 0],
\]

where $\kappa_1, \kappa_2, \ldots, \kappa_{n-1}$ are the principal curvatures of $\partial\Omega$ with respect to the
inner normal at $x_0$; (see Appendix to Chapter 14). If $a_1, a_2, \ldots, a_n$ denote the
corresponding diagonal elements of the matrix $P^{\!\top}[a^{ij}(x_0)]P$, we see that

\[
(6.56) \qquad \left(a^{ij}D_{ij}d\right)_{x=x_0} = -\sum_{i=1}^{n-1} a_i\kappa_i.
\]

Hence, if

\[
(6.57) \qquad \sum_{i=1}^{n-1} a_i\kappa_i > \sup_{\Omega} |b|
\]

it follows by continuity that the inequality (6.55) is satisfied in $G = B \cap \Omega$ for
some ball $B$ about $x_0$, provided $K$ is chosen large enough. If in addition, the coefficients
$b^i$ are continuous at $x_0$, and if $b_\nu$ denotes the normal component of the vector
$b(x_0) = (b^1(x_0), \ldots, b^n(x_0))$ with respect to the inner normal, condition (6.55)
is satisfied in a suitable $G$ provided

\[
(6.58) \qquad \sum_{i=1}^{n-1} a_i\kappa_i - b_\nu > 0.
\]

% PDF page 132
\source{gilbarg-trudinger:page-0132}

Under the stated hypotheses, (6.57) and (6.58) are therefore sufficient conditions
for the existence of a local barrier at $x_0$. If $b(x_0)=0$, (6.58) reduces to the simple
condition,

\[
\sum_{i=1}^{n} a_i \kappa_i > 0,
\]

which involves only the leading coefficients and the principal curvatures and
directions of the boundary. It is not difficult to remove the continuity assumptions
on the coefficients at the boundary and to reformulate (6.57) and (6.58)
appropriately.

We remark that when $\partial\Omega$ is not in $C^2$, the preceding considerations are still
applicable provided a $C^2$ domain $\widetilde{\Omega}$ can be found such that $x_0 \in \partial\Omega \cap \partial\widetilde{\Omega}$ and
$B \cap \Omega \subset B \cap \widetilde{\Omega}$ for some ball $B$ containing $x_0$. The above conditions for the
existence of a barrier at $x_0$ then remain valid if the distance function $d(x)$ is re-
placed by $\widetilde{d}(x)=\operatorname{dist}(x,\partial\widetilde{\Omega})$.

The preceding remarks and those earlier in this section concerning the exterior
sphere condition show the existence of a barrier at the following sets of boundary
points of a $C^2$ domain $\Omega$ where the coefficients $a^{ij},\, b^i$ are continuous:

\[
\Sigma_1 = \{x_0 \in \partial\Omega \mid \text{inequality (6.58) holds}\}
\]
\[
\Sigma_2 = \{x_0 \in \partial\Omega \mid a^{ij}v_i(x_0)v_j(x_0)\neq 0,\ \text{where } v(x_0)\text{ is normal to }\partial\Omega\}.
\]

Taken together with Theorem 6.11 these results yield:

\textbf{Theorem 6.25.} \emph{Let the operator $L$ in (6.1) be elliptic in a $C^2$ domain $\Omega$ with $c \leq 0$ in
$\Omega$ and $a^{ij},\, b^i/\lambda,\, c,\, f/\lambda \in C^\alpha(\Omega)\cap C^0(\overline{\Omega})$. Assume $\Sigma_1 \cup \Sigma_2 = \partial\Omega$. Then the problem}

\[
Lu=f\ \text{in }\Omega,\qquad u=\varphi\ \text{on }\partial\Omega,
\]

\emph{has a unique solution $u\in C^2(\Omega)\cap C^0(\overline{\Omega})$ for all $\varphi\in C^0(\partial\Omega)$.}

We note that under certain conditions on $L$ and $\partial\Omega$, a solution is uniquely
determined by the values of $\varphi$ on $\Sigma_1 \cup \Sigma_2$ even if $\Sigma_1 \cup \Sigma_2 \neq \partial\Omega$, as in the example
(6.54). A maximum principle related to such a boundary value problem is contained
in Problem 6.10.

\subsection*{6.7. Other Boundary Conditions; the Oblique Derivative Problem}

Until now we have been concerned only with Dirichlet boundary conditions. We
now develop an analogue of the Schauder theory for the regular oblique derivative
problem.

% PDF page 133
\source{gilbarg-trudinger:page-0133}

\section*{6.7. Other Boundary Conditions; the Oblique Derivative Problem}

\subsection*{Poisson's Equation}

In extending the Schauder theory to other linear boundary value problems our
point of departure is the theory of Poisson's equation in the half-space $\R_+^n =
\{x \mid x_n>0\}$ with the oblique derivative boundary condition,

\begin{equation}
\tag{6.60}
\N u \equiv a u + \sum_{i=1}^n b_i D_i u = \varphi \quad \text{on } x_n=0,
\end{equation}

in which the coefficients $a, b_i$ are constants. We also write the boundary operator $\N$
in the equivalent forms

\[
\N u \equiv a u + \bvec \cdot D u = a u + b_t D_t u + b_n D_n u
\]

where $\bvec = (b_1,\ldots,b_n) = (b_t,b_n)$ and $D = (D_1,\ldots,D_n) = (D_t,D_n)$, $D_t$ denoting the
tangential gradient. We assume throughout the regular oblique derivative condition
$b_n \neq 0$, and to fix our ideas we take at first

\begin{equation}
\tag{6.61}
b_n>0, \qquad |\bvec| = (|b_t|^2 + b_n^2)^{1/2}=1.
\end{equation}

The latter condition is an inessential normalization which allows us to write

\[
\N u = a u + D_s u
\]

where $D_s u = \partial u/\partial s$ is the directional derivative in the direction of the vector $\bvec$.
We consider initially the homogeneous boundary condition $\N u = 0$ on $x_n = 0$
and construct a harmonic Green's function in $\R_+^n$ satisfying this boundary condition. Letting $\Gamma$ denote the fundamental solution (4.1) of Laplace's equation, we
write for $n \geq 3$ and $a \leq 0$

\begin{equation}
\tag{6.62}
G(x,y)=\Gamma(x-y)-\Gamma(x-y^*)-2b_n\int_0^\infty e^{as} D_n\Gamma(x-y^*+bs)\,ds,
\end{equation}

where $x, y \in \R_+^n$, $D_n = \partial/\partial x_n$, and $y^*=(y_1,\ldots,y_{n-1},-y_n)=(y',-y_n)$. Clearly
$G$ is harmonic in $x$ and $y$ for $x\neq y$ and by direct calculation one finds

\begin{equation}
\tag{6.63}
\N G(x,y)=0 \quad \text{on } x_n=0.
\end{equation}

(Here $\N$ operates on $G$ as a function of $x$ for fixed $y$.) Thus $G$ has the required
properties of a Green's function for the boundary condition (6.63).
The choice of $G$ in (6.62) is motivated by the following considerations. If
$G(x,y)=\Gamma(x-y)+h(x,y)$ is the desired harmonic Green's function satisfying
(6.63), then $\N G$ is also harmonic in $x$ (for $y\neq x$) and vanishes on $x_n = 0$. It follows
from the Schwarz reflection principle (Problem 2.4) that $\N G$ is regular in $\R^n$
except for the singularity

\[
a\Gamma(x-y)+b_t D_t\Gamma(x-y)+b_n D_n\Gamma(x-y)-a\Gamma(x-y^*)
\]
\[
-\, b_t D_t\Gamma(x-y^*)+b_n D_n\Gamma(x-y^*).
\]


% PDF page 134
\source{gilbarg-trudinger:page-0134}
\pagestyle{empty}



Here we have used the fact that $D_i\Gamma(x^* - y^*) = D_i\Gamma(x - y)$ for $i = 1, \ldots, n - 1$,
while $D_n\Gamma(x^* - y^*) = -D_n\Gamma(x - y)$. Thus if we impose the condition that $NG$
vanish at infinity, it follows from the Liouville theorem

\[
\Ds h + ah = -a\Gamma(x - y^*) - b_t D_t\Gamma(x - y^*) + b_n \Dn\Gamma(x - y^*)
\]
\[
= -a\Gamma(x - y^*) - \Ds\Gamma(x - y^*) + 2b_n \Dn\Gamma(x - y^*).
\]

This implies

\[
\Ds\!\left[e^{as} h(x + bs, y)\right]
= -\left[a\Gamma(x - y^* + bs) + \Ds\Gamma(x - y^* + bs)\right] e^{as}
\]
\[
\qquad + 2b_n \Dn\Gamma(x - y^* + bs) e^{as}.
\]

Integrating with respect to $s$ from $s = 0$ to $s = \infty$ and then integrating by parts, we
obtain

\[
h(x, y) = -\Gamma(x - y^*) - 2b_n \int_0^\infty e^{as}\,\Dn\Gamma(x - y^* + bs)\,ds.
\]

This expression for $h$ gives us (6.62), which is valid also when $b_n = 0$.
We now examine $G(x, y)$ more closely with the objective of deriving estimates
analogous to those in Chapter 4 for the Newtonian potential and solutions of
Poisson's equation.

Letting $\xi = (x - y^*)/|x - y^*|$, we have

\[
G(x, y) = \Gamma(x - y) - \Gamma(x - y^*) + \Theta(x, y)
\]

where

\[
\Theta(x, y) = -2b_n \int_0^\infty e^{as}\Dn\Gamma(x - y^* + bs)\,ds,\qquad a \le 0,
\]
\[
= -|x - y^*|^{2-n}\!\left[\frac{2b_n}{n\omega_n}\int_0^\infty e^{a|x-y^*|s}
\frac{(\xi_n + b_n s)\,ds}{\left(1 + 2(\xi\cdot b)s + s^2\right)^{n/2}}\right]
\]
\[
= |x - y^*|^{2-n} g(\xi, |x - y^*|).
\]

One sees that $g$ is regular in its arguments since (by (6.61)) $\xi\cdot b > -|b_t| > -1$
for all $x, y \in \R^n$ and hence the denominator of the integrand is bounded away from
zero. The function $\Theta(x, y) = \Theta(x - y^*)$ satisfies the relations

\begin{equation}
\label{eq:6.64}
D_{x_i}\Theta(x, y) = -D_{y_i}\Theta(x, y),\qquad i = 1, \ldots, n - 1 ;
\end{equation}
\[
D_{x_n}\Theta(x, y) = D_{y_n}\Theta(x, y);
\]
\[
|D^\beta \Theta(x, y)| \le C|x - y^*|^{2-n-|\beta|},\qquad C = C(n, |\beta|, b_n).
\]


% PDF page 135
\source{gilbarg-trudinger:page-0135}



\noindent\textbf{6.7. Other Boundary Conditions; the Oblique Derivative Problem}\hfill 123

\medskip

These relations will suffice in extending the details of Lemma 4.10 for the Newtonian
potential to the analogous integral $\int \Theta(x,y)f(y)\,dy$. If $|b|\neq 1$ we replace $a$ by
$a/|b|$ and $b_n$ by $b_n/|b|$ in the preceding. We note that since $a\leq 0$ the constant $C$ in
(6.64) can be taken independent of $a$.

\medskip

\noindent\textbf{Theorem 6.26.} \textit{Let $B_1=B_R(x_0)$, $B_2=B_{2R}(x_0)$ denote balls with center $x_0\in \bar{\RR}^{\,n}_+$,
and let $B_1^+=B_1\cap \RR^n_+$, $B_2^+=B_2\cap \RR^n_+$, $T=B_2\cap\{x_n=0\}$. Suppose that $u\in$
$C^2(B_2^+)\cap C^1(B_2^+\cup T)$ satisfies $\Delta u=f$ in $B_2^+$, $f\in C^\alpha(\bar B_2^+)$, and the boundary
condition (6.60), $\Nu u=\varphi$ on $T$, with $a\leq 0$, $b_n>0$ and $\varphi\in C^{1,\alpha}(T)$. Then $u\in C^{2,\alpha}(\bar B_1^+)$
and}

\begin{equation}
\tag{6.65}
|u|_{2,\alpha; B_1^+}\leq C\bigl(|u|_{0; B_2^+}+R|\varphi|_{1,\alpha; T}+R^2|f|_{0,\alpha; B_2^+}\bigr),
\end{equation}

\noindent where $C=C(n,\alpha,b_n/|b|)$.

\medskip

\noindent (Here $|\,\cdot\,|$ denotes the weighted norms (4.6)' defined with respect to $R$.)

\medskip

\noindent\textit{Proof.} It is assumed that $T$ is non-empty for otherwise the stated result is already
contained in Theorem 4.6. Suppose first that $\varphi=0$, $|b|=1$ and $n>2$. Consider the
function

\begin{equation*}
w(x)=\int_{B_2^+} G(x,y)f(y)\,dy=w_1(x)+w_2(x),
\end{equation*}

\noindent where

\begin{equation*}
w'_1(x)=\int_{B_2^+}\bigl[\Gamma(x-y)-\Gamma(x-y^*)\bigr]f(y)\,dy
\end{equation*}

\begin{equation*}
w_2(x)=\int_{B_2^+}\Theta(x,y)f(y)\,dy.
\end{equation*}

\noindent We have already seen in (4.26) that $w_1$ satisfies the estimate

\begin{equation}
\tag{6.66}
|D^2 w_1|_{0,\alpha; B_1^+}\leq C|f|_{0,\alpha; B_2^+}, \qquad C=C(n,\alpha).
\end{equation}

\noindent The estimation of $w_2$ is essentially the same as that for the Newtonian potential
in Lemma 4.10. Let $f(x)$ be defined by even reflection with respect to $x_n$, so that
$f(x',-x_n)=f(x',x_n)$. Then a representation analogous to (4.9) is valid for the
second derivatives of $w_2$; namely, for $x\in B_2^+$ and $i,j=1,\dots,n$, we have

\begin{equation*}
D_{ij}w_2(x)=\int_{B_2^+}D_{ij}\Theta(x-y^*)\bigl(f(y^*)-f(x)\bigr)\,dy
 -f(x)\int_{\partial B_2^+}D_i\Theta(x-y^*)\nu_j(y)\,ds_y,
\end{equation*}


% PDF page 136
\source{gilbarg-trudinger:page-0136}



\noindent where $v=(v_1,\dots,v_n)$ is the outward unit normal to $\partial B_2^+$. By virtue of (6.64) the
arguments of Lemmas 4.4 and 4.10 carry over without essential change to give the
estimate

\[
\bigl|D^2 w_2\bigr|_{0,\alpha; B_1^+}' \leq C|f|_{0,\alpha; B_2^+}', \qquad C=C(n,\alpha,b_n).
\]

\noindent Combining this inequality with (6.66), we obtain

\begin{equation}
\tag{6.67}
\bigl|D^2 w\bigr|_{0,\alpha; B_1^+} \leq C|f|_{0,\alpha; B_2^+}', \qquad C=C(n,\alpha,b_n).
\end{equation}

\noindent If $|b|\neq 1$ we replace $b_n$ by $b_n/|b|$ in this estimate.
\par To obtain estimates for $u$ from the preceding, let $\eta \in C_0^2(B_2)$ be a cutoff function
such that $\eta(x)=1$ for $|x-x_0|\leq (3/2)R$ and $|D^\beta \eta|\leq C/R^{|\beta|}$ for $|\beta|\leq 2$. Then, we
have

\[
u(x)\eta(x)=\int_{B_2^+} G(x,y)\Delta[u(y)\eta(y)]\,dy, \qquad x\in B_2^+.
\]

\noindent If $x\in B_1^+$, so that $|x-y|>R/2$ where $D\eta\neq 0$, we obtain

\[
u(x)=\int_{B_2^+} (\eta Gf + u\Delta\eta)\,dy + 2\int_{B_2^+} GDu\cdot D\eta\,dy,
\]

\[
=\int_{B_2^+} (\eta Gf + u\Delta\eta)\,dy - 2\int_{B_2^+} u(DG\cdot D\eta + G\Delta\eta)\,dy.
\]

\noindent The required estimate (6.65) when $\varphi=0$ now follows from (6.67), the bounds
$|D^\beta\eta|\leq C/R^{|\beta|}$, and, by (2.14) and (6.64),

\[
|D^\beta G(x,y)|\leq C|x-y|^{2-n-|\beta|}\leq CR^{2-n-|\beta|},
\]

\noindent where the indicated derivatives may be with respect to both $x$ and $y$ variables.
\par We now remove the restriction that $\varphi=0$. For this purpose we seek a function
$\psi\in C^{2,\alpha}(\overline{B_2^+})$ satisfying $\Nu\psi=\varphi$ on $T$. It can be assumed that $\varphi$ is suitably extended
outside $T$ so that $\varphi\in C^{1,\alpha}(\RR^{n-1})$ on $x_n=0$; (see Lemma 6.38). Choosing a non-
negative function $\eta\in C_0^2(\RR^{n-1})$ such that $\int \eta(y')\,dy'=1$, $y'=(y_1,\dots,y_{n-1})$, we
define

\begin{equation}
\tag{6.68}
\psi(x)=\psi(x',x_n)=b_n^{-1}x_n \int_{\RR^{n-1}} \varphi(x'-x_n y')\eta(y')\,dy'.
\end{equation}

\noindent One verifies easily that $\psi(x',0)=0$, $D_n\psi(x',0)=b_n^{-1}\varphi(x')$, and hence

\begin{equation}
\tag{6.69}
\Nu\psi=\varphi \qquad \text{on } x_n=0.
\end{equation}


% PDF page 137
\source{gilbarg-trudinger:page-0137}

\providecommand{\R}{\mathbb{R}}
\providecommand{\T}{T}
\providecommand{\B}{B}
\providecommand{\cC}{C}
\providecommand{\cN}{N}
\providecommand{\cL}{L}

That $\psi \in C^{2,\alpha}(\R_+^n)$ is seen from the relations
\[
b_n D_{ij}\psi(x)=\int D_i\varphi(x'-x_n y')D_j\eta(y')\,dy',
 \qquad i,j\neq n;
\]
\[
b_n D_i\psi(x)=-\int y'\cdot D\varphi(x'-x_n y')D_i\eta(y')\,dy',
 \qquad i\neq n;
\]
\[
b_n D_{nn}\psi(x)=\int y'\cdot D\varphi(x'-x_n y')\big[(n-2)\eta(y')+y'\cdot D\eta(y')\big]\,dy'.
\]

We observe also that
\begin{equation}\tag{6.70}
b_n|\psi|_{2,x;B_2^+}\le C(R|\varphi|_{0;T}+R^2|D\varphi|_{0;T}+R^{2+\alpha}[D\varphi]_{\alpha;T})
=CR|\varphi|_{1,\alpha;T},
\end{equation}
where the constant $C$ depends only on $n$ and the choice of $\eta$.
We are now able to make the reduction to the case $\varphi=0$. Letting $v=u-\psi$,
it follows that $\Delta v=f-\Delta\psi\in C^\alpha(\bar B_2^+)$, and by virtue of (6.69) we have $Nv=0$ on
$T$. From (6.70) and the proven estimate for $v$ we obtain (6.65). $\square$

We remark that in general the constant $C$ in (6.65) becomes unbounded as
$b_n\to 0$.
The following estimate is a consequence of the preceding theorem and is stated
without proof. The details are similar to the proof of Theorem 4.8.

\textbf{Lemma 6.27.} \textit{Let $\Omega$ be a bounded open set in $\R_+^n$ with boundary portion $T$ on
$x_n=0$, and let $u\in C^2(\Omega)\cap C^1(\Omega\cup T)$ satisfy $\Delta u=f$ in $\Omega$, $f\in C^\alpha(\bar\Omega)$, and the
boundary condition (6.60), $Nu=\varphi$ on $T$, with $a\le 0$, $b_n>0$ and $\varphi\in C^{1,\alpha}(T)$. Then}
\[
|u|_{2,x;\Omega\cup T}^*\le C(|u|_{0;\Omega}+|\varphi|_{1,\alpha;T}+|f|_{0,x;\Omega}),
\]
\textit{where $C=C(n,\alpha,b_n/|b|,\operatorname{diam}\Omega)$.}

The same argument as in Lemma 6.1 provides the following extension to the
oblique derivative problem.

\textbf{Lemma 6.28.} \textit{Under the same hypotheses as in Lemma 6.27 let $u$ satisfy $L_0u=f$
(in place of $\Delta u=f$), where $L_0$ is the constant coefficient operator defined in Lemma
6.1. Then}
\begin{equation}\tag{6.71}
|u|_{2,x;\Omega\cup T}^*\le C(|u|_{0;\Omega}+|\varphi|_{1,\alpha;T}+|f|_{0,x;\Omega})
\end{equation}
\textit{where $C=C(n,\alpha,\lambda,\Lambda,b_n/|b|,\operatorname{diam}\Omega)$.}

\textbf{Variable Coefficients}

We consider now equations with variable coefficients and the corresponding oblique
derivative problem in domains with curved boundaries. The extension of the
estimates of the Schauder theory to these boundary conditions follows closely the
ideas in Sections 6.1 and 6.2. We derive first an analogue of Lemma 6.4.

% PDF page 138
\source{gilbarg-trudinger:page-0138}

\providecommand{\R}{\mathbb{R}}
\providecommand{\cN}{N}
\providecommand{\cL}{L}

Lemma 6.29. \emph{Let $\Omega$ be a bounded open set in $\R_+^n$ with a boundary portion $T$ on
$x_n=0$. Suppose that $u\in C^{2,\alpha}(\Omega\cup T)$ is a solution in $\Omega$ of $Lu=f$ (equation (6.1))
satisfying the boundary condition}

\begin{equation}\tag{6.72}
N(x')u=\gamma(x')u+\sum_{i=1}^n \beta_i(x')D_i u=\varphi(x'), \qquad x'\in T,
\end{equation}

\emph{where $|\beta_n|\ge \kappa>0$ for some constant $\kappa$. It is assumed that $L$ satisfies (6.2) and that
$f\in C^\alpha(\bar\Omega)$, $\varphi\in C^{1,\alpha}(\bar T)$, $a^{ij}, b^i, c\in C^\alpha(\bar\Omega)$ and $\gamma, \beta_i\in C^{1,\alpha}(\bar T)$ with}

\[
|a^{ij}, b^i, c|_{0,\alpha;\Omega},\ |\gamma, \beta_i|_{1,\alpha;T}\le \Lambda,\qquad i,j=1,\ldots,n.
\]

\emph{Then}

\begin{equation}\tag{6.73}
|u|_{2,\alpha;\Omega\cup T}^*\le C(|u|_{0;\Omega}+|\varphi|_{1,\alpha;T}+|f|_{0,\alpha;\Omega}),
\end{equation}

\emph{where $C=C(n,\alpha,\lambda,\Lambda,\kappa,\operatorname{diam}\Omega)$.}

\emph{Proof. We remark first that it can be assumed $\gamma\le 0$ and $\beta_n>0$. For by setting
$v=ue^{kx_n}$, where $k\ge \sup|\gamma|/\kappa$, the boundary condition (6.72) is taken into}

\[
N'v=(\gamma-k\beta_n)v+\sum \beta_i D_i v=\varphi \qquad \text{on } T
\]

\emph{with $\beta_n(\gamma-k\beta_n)\le 0$. At the same time $Lu=f$ is transformed into an equation
$L'v=f'$ satisfying the same hypotheses as in the theorem. The desired estimate (6.73)
is obviously equivalent to the corresponding estimate for $v$.}

\emph{The technique in Theorem 6.2 and Lemma 6.4 of ``freezing'' the coefficients is
once again applicable, with certain modifications due to the boundary condition
(6.72). Thus, let $x_0, y_0$ be any two distinct points in $\Omega$ and suppose $d_{x_0}=\min(\bar d_{x_0}, \bar d_{y_0})$,
where $\bar d_x=\operatorname{dist}(x,\partial\Omega-T)$. Let $\mu\le 1/4$ be a positive constant (to be specified
later), and set $d=\mu \bar d_{x_0},\ B_d=B_d(x_0)$. If $B_d\cap T\ne \varnothing$ let $x_0'$ denote the projection of
$x_0$ on $T$. As in Theorem 6.2 we rewrite the equation $Lu=f$ in the form (6.15) and
the boundary condition (6.72) as}

\begin{equation}\tag{6.74}
N(x_0')u=[N(x_0)-N(x')]u(x')+\varphi(x')\equiv \Phi(x'), \qquad x'\in T,
\end{equation}

\emph{and we consider (6.15), (6.74) as a problem in $B_d\cap\Omega$ with constant coefficients,
disregarding (6.74) if $B_d\cap T=\varnothing$. The argument in Theorem 6.2 and the analogous
one indicated in Lemma 6.4 proceed in essentially the same way, except that
Lemma 6.28 now replaces Lemma 6.1 in the details. Thus in place of (6.16) we now
obtain}

\begin{equation}\tag{6.75}
d_{x_0}^{2+\alpha}\frac{|D^2u(x_0)-D^2u(y_0)|}{|x_0-y_0|^\alpha}\le \frac{C}{\mu^{2+\alpha}}
\Bigl(|u|_{0;\Omega}+|\Phi|_{1,\alpha;B_d\cap T}+|F|_{0,\alpha;B_d\cap\Omega}\Bigr)
\frac{4}{\mu^\alpha}[u]^*_{2;\Omega\cup T}.
\end{equation}

% PDF page 139
\source{gilbarg-trudinger:page-0139}

\setcounter{section}{6}
\setcounter{page}{127}
\noindent 6.7.\ Other Boundary Conditions; the Oblique Derivative Problem \hfill 127

\medskip

The estimation of the right number is essentially as in Theorem 6.2 with the
exception of the additional term $|\Phi|_{1,\alpha;B\cap T}$. Concerning this term, one sees that

\[
|[N(x_0)-N(x') ]u(x')|_{1,\alpha;B\cap T}\le C\mu^{2+\alpha}(C(\mu)|u|_{0;\Omega}+\mu^\alpha[u]^*_{2,\alpha;\Omega\cap T}).
\]

The details are similar to those leading to (6.19) and are not carried out here.
Combining this estimate with that for the other terms in (6.75) we arrive at the
desired estimate (6.73). $\square$

\medskip

The preceding lemma can now be extended to domains with curved boundaries.
A repetition of the arguments in Lemma 6.5 and Theorem 6.6 leads to the following
global estimate for solutions of the oblique derivative problem.

\medskip

\noindent\textbf{Theorem 6.30.} \textit{Let $\Omega$ be a $C^{2,\alpha}$ domain in $\R^n$, and let $u\in C^{2,\alpha}(\overline{\Omega})$ be a solution in
$\Omega$ of $Lu=f$ satisfying the boundary condition}

\[
N(x)u=\gamma(x)u+\sum_{i=1}^n \beta_i(x)D_i u=\varphi(x),\qquad x\in\partial\Omega,
\]

\textit{where the normal component $\beta_\nu$ of the vector $\beta=(\beta_1,\dots,\beta_n)$ is non-zero and}

\begin{equation}\tag{6.76}
|\beta_\nu|\ge \kappa>0\ \text{on }\partial\Omega \qquad (\kappa=\text{const}).
\end{equation}

\textit{It is assumed that $L$ satisfies (6.2) and that $f\in C^\alpha(\overline{\Omega})$, $\varphi\in C^{1,\alpha}(\overline{\Omega})$, $a^{ij}$, $b^i$, $c\in C^\alpha(\overline{\Omega})$
and $\gamma$, $\beta_i\in C^{1,\alpha}(\overline{\Omega})$ with}

\[
|a^{ij},\,b^i,\,c|_{0,\alpha;\Omega}\cdot|\gamma,\beta_i|_{1,\alpha;\Omega}\le \Lambda,\qquad i,j=1,\dots,n.
\]

\textit{Then}

\begin{equation}\tag{6.77}
|u|_{2,\alpha;\Omega}\le C\bigl(|u|_{0;\Omega}+|\varphi|_{1,\alpha;\Omega}+|f|_{0,\alpha;\Omega}\bigr)
\end{equation}

\textit{where $C=C(n,\alpha,\lambda,\Lambda,\kappa,\Omega)$.}

\medskip

\noindent\textit{Remarks.}
1) Condition (6.76) implies that the directional derivative $\beta\cdot \D u$ is
nowhere tangential to $\partial\Omega$. This hypothesis is essential in the present considerations.

2) In the statement of the theorem it is convenient and involves no loss of generality
to assume that $\varphi$, $\gamma$ and $\beta_i$ are defined globally (rather than on $\partial\Omega$), so that the
norm $|\ \ |_{1,\alpha;\Omega}$ is well defined for these functions. In Theorem 6.26 and Lemmas
6.27--6.29 where $T$ was a hyperplane boundary portion, a global extension of
$\varphi$ (and of $\gamma$, $\beta_i$ in Lemma 6.29) was not used since the norm $|\ \ |_{1,\alpha;T}$ was naturally
defined. 3) That $|\varphi|_{1,\alpha}$ appears in the estimate (6.77) is to be expected since $Nu$
involves first order differentiation of the $C^{2,\alpha}$ function $u$. This is in contrast with
the corresponding global estimate (6.36) for the Dirichlet boundary condition
which requires that $\varphi\in C^{2,\alpha}$.


% PDF page 140
\source{gilbarg-trudinger:page-0140}

We have thus far been concerned only with estimates for the oblique derivative problem. The actual solution of the problem for $Lu=f$ can be reduced to that for Poisson's equation by the method of continuity (as in Theorem 6.8), but the method now involves a continuous family of both differential operators and boundary operators. Unique solvability of the oblique derivative problem, under appropriate additional restrictions on the operators $L$ and $N$, is provided in the following theorem.

\noindent\textbf{Theorem 6.31.} \textit{Let $L$ be a strictly elliptic operator with $c\leq 0$ and coefficients in $C^{\alpha}(\overline{\Omega})$ in a $C^{2,\alpha}$ domain $\Omega$. Let $Nu \equiv \gamma u + \beta\cdot Du$ define a boundary operator on $\partial\Omega$ such that $\gamma(\beta\cdot v)>0$ on $\partial\Omega$ if $v$ is the outward unit normal on $\partial\Omega$. Assume that $\gamma$, $\beta \in C^{1,\alpha}(\partial\Omega)$. Then the oblique derivative problem}

\begin{equation}
\tag{6.78}
Lu=f \ \text{in }\Omega, \qquad Nu=\varphi \ \text{on }\partial\Omega
\end{equation}

\textit{has a unique $C^{2,\alpha}(\overline{\Omega})$ solution for all $f\in C^{\alpha}(\overline{\Omega})$ and $\varphi\in C^{1,\alpha}(\partial\Omega)$.}

\noindent\textit{Proof.} We assume without loss of generality that $\gamma>0$ and $\beta\cdot v>0$ on $\partial\Omega$ and also that $\varphi$ and $\beta$ are extended to all of $\overline{\Omega}$ and are contained in $C^{1,\alpha}(\overline{\Omega})$. Consider the family of problems for $0\leq t\leq 1$:

\[
L_tu \equiv tLu+(1-t)\Delta u=f \ \text{in }\Omega,
\]
\begin{equation}
\tag{6.79}
N_tu \equiv tNu+(1-t)\left(\frac{\partial u}{\partial v}+u\right)=\varphi \ \text{on }\partial\Omega.
\end{equation}

We note that $L_1=L$, $L_0=\Delta$, $N_1=N$, $N_0=\partial/\partial v+\text{identity}$, and that for suitable positive constants $\lambda$, $\Lambda$ the coefficients $a^{ij}_t$, $b^i_t$, $c_t$ of $L_t$ satisfy

\[
a^{ij}_t\xi_i\xi_j \geq \lambda_t|\xi|^2=\min(1,\lambda)|\xi|^2 \qquad \forall x\in\Omega,\ \xi\in\mathbb{R}^n
\]

and

\[
|a^{ij}_t|,\,|b^i_t|,\,|c_t|_{0,\alpha} \leq \Lambda_t=\max(1,\Lambda).
\]

Also, since $\beta\cdot v\geq \beta'>0$ and $\gamma\geq \gamma'>0$ for some constants $\beta'$, $\gamma'$, we have

\[
\gamma_t=(1-t)+t\gamma \geq \min(1,\gamma')>0
\]
\[
\beta_t\cdot v=(1-t)+t\beta\cdot v \geq \min(1,\beta')>0 \ \text{on }\partial\Omega
\]

while $|\beta|_{1,\alpha}$ is also bounded independently of $t$.

Consider any solution $u\in C^{2,\alpha}(\overline{\Omega})$ of the problem (6.79) (for some $t$). $|u|_{2,\alpha}$ satisfies the estimate (6.77) with a constant $C$ independent of $t$. By estimating $|u|_0$ in terms of $\varphi$ and $f$ we shall obtain in addition a bound

\begin{equation}
\tag{6.80}
|u|_{2,\alpha}\leq C\left(|\varphi|_{1,\alpha}+|f|_{0,\alpha}\right)
\end{equation}

valid for all $C^{2,\alpha}(\overline{\Omega})$ solutions of (6.79).


% PDF page 141
\source{gilbarg-trudinger:page-0141}

To estimate $|u|_0$ we first make a substitution $v=u/\omega$ where $\omega$ is a fixed $C^0(\overline{\Omega})\cap C^2(\Omega)$ function (independent of $t$) satisfying the conditions:
(i) $\omega\ge \bar{\omega}>0$;
(ii) $L_t\omega\le \bar{c}<0$;
(iii) $\gamma_t+\beta_t\cdot D\omega/\omega\ge \bar{\gamma}>0$ on $\partial\Omega$, where $\bar{\omega}$, $\bar{c}$, $\bar{\gamma}$ are constants.
Such a function $\omega$ can be chosen in the form $\omega(x)=c_1-c_2 e^{\mu x_1}$ if $\mu$ is sufficiently large and $c_1$, $c_2$ are suitable positive constants.
The substitution $v=u/\omega$ transforms (6.79) into another problem:
\[
\bar{L}_t v=\bar{f}=f/\omega \text{ in } \Omega,\qquad \bar{N}_t v=\bar{\varphi}=\varphi/\omega \text{ on } \partial\Omega,
\]
in which the coefficient $L_t\omega/\omega$ of $v$ in $\bar{L}_t v$ satisfies $L_t\omega/\omega\le \bar{c}/\omega<0$ and the coefficient $\gamma_t$ of $v$ in $\bar{N}_t v$ satisfies
\[
\bar{\gamma}_t=\gamma_t+\beta_t\cdot D\omega/\omega\ge \bar{\gamma}>0.
\]
If now sup $|v|=|v(x_0)|$ at some $x_0\in\Omega$, then
\[
\sup_{\Omega}|v|=|v(x_0)|\le |f(x_0)|/\bar{c}|\le \sup_{\Omega}|f|/|\bar{c}|,
\]
and hence
\[
|u|_0\le \sup_{\Omega}\omega\ \sup_{\Omega}|v|\le C|f|_0,
\]
where $C$ is a constant independent of $t$.
On the other hand, if sup $|v|=|v(x_0)|$ for some $x_0\in\partial\Omega$, then either
\[
\sup_{\Omega}|v|=v(x_0)\le \bar{\gamma}^{-1}\bigl(\bar{\varphi}-\beta_t\cdot Dv\bigr)_{x=x_0}\le \bar{\varphi}(x_0)/\bar{\gamma}
\]
or
\[
\sup_{\Omega}|v|=-v(x_0)\le \bar{\gamma}^{-1}\bigl(-\bar{\varphi}+\beta_t\cdot Dv\bigr)_{x=x_0}\le -\bar{\varphi}(x_0)/\bar{\gamma}
\]
and thus $|u|_0\le C\sup_{\partial\Omega}|\varphi|$.
The estimate (6.80) follows at once from (6.77).

The argument proceeds now essentially as in Theorem 6.8. Let
\[
\mathscr{B}_1=C^{2,\alpha}(\overline{\Omega}),\qquad \mathscr{B}_2=C^{\alpha}(\overline{\Omega})\times C^{1,\alpha}(\partial\Omega),
\]
with
\[
\|(f,\varphi)\|_{\mathscr{B}_2}=|f|_{0,\alpha;\Omega}+|\varphi|_{1,\alpha;\partial\Omega},
\]
and consider the operator
\[
\mathcal{L}_t=(L_t,N_t):\mathscr{B}_1\to \mathscr{B}_2.
\]
The solvability of the problem (6.79) for arbitrary $f\in C^{\alpha}(\overline{\Omega})$, $\varphi\in C^{1,\alpha}(\partial\Omega)$ is equivalent to showing that $\mathcal{L}_t$ is one-to-one and onto.
Let $u_t$ denote a solution of this problem for given $f,\varphi$.
It is unique (see Problem 3.1) and from (6.80) we have the bound
\[
|u_t|_{2,\alpha}\le C\bigl(|f|_{0,\alpha}+|\varphi|_{1,\alpha}\bigr)
\]


% PDF page 142
\source{gilbarg-trudinger:page-0142}

or, equivalently,

\begin{equation}
\tag{6.81}
\|u\|_{\mathscr{B}_1}\le C\|\mathcal{L}_t u\|_{\mathscr{B}_2},
\end{equation}

the constant $C$ being independent of $t$. The fact that $\mathcal{L}_0$ is invertible is a consequence
of the solvability in $C^{2,\alpha}(\overline{\Omega})$ of the third boundary value problem:

\[
\Delta u=f\ \text{in }\Omega,\qquad u+\frac{\partial u}{\partial v}=\varphi\ \text{on }\partial\Omega.
\]

We refer to the literature of potential theory for this result; (see [GU], for example).
Assuming this result, we conclude from (6.81) and the method of continuity
(Theorem 5.2) that the theorem is proved. $\square$

If either of the conditions $\gamma>0$ or $c\leq 0$ is not satisfied in the preceding theorem,
unique solvability is no longer assured, and one can assert instead the Fredholm
alternative as in Theorem 6.15. The method of proof is basically the same. An
immediate consequence of the alternative is solvability when $c\leq 0$, $\gamma\geq 0$ and
either $c\neq 0$ or $\gamma\neq 0$, since uniqueness holds under these conditions.

\bigskip

\section*{6.8.\ Appendix 1: Interpolation Inequalities}

We prove here the interpolation inequalities quoted in the course of Chapter 6. We
begin with the inequalities for interior norms and seminorms.

\medskip

\noindent\textbf{Lemma 6.32.} \textit{Suppose $j+\beta<k+\alpha$, where $j$, $k=0$, 1, 2, . . . , and $0\leq\alpha$, $\beta\leq 1$.
Let $\Omega$ be an open subset of $\mathbb{R}^n$ and assume $u\in C^{k,\alpha}(\Omega)$. Then for any $\varepsilon>0$ and some
constant $C=C(\varepsilon, k, j)$ we have}

\begin{equation}
\tag{6.82}
[u]^{*}_{j,\beta;\Omega}\leq C|u|_{0;\Omega}+\varepsilon [u]^{*}_{k,\alpha;\Omega}.
\end{equation}
\[
|u|^{*}_{j,\beta;\Omega}\leq C|u|_{0;\Omega}+\varepsilon [u]^{*}_{k,\alpha;\Omega}.
\]

\textit{Proof.} \textbf{We shall establish (6.82) for the cases $j$, $k=0$, 1, 2 needed in the present
work. A direct extension of the ideas and a suitable induction yield the stated result
for arbitrary $j$, $k$.}

It is assumed that the right member of (6.82) is finite since otherwise the assertion is trivial. For notational convenience we omit the subscript $\Omega$, the domain $\Omega$
being implicitly understood. We consider several cases:

\qquad (i) $j=1$, $k=2$; $\alpha=\beta=0$. We wish to show

\begin{equation}
\tag{6.83}
[u]^*_1\leq C(\varepsilon)|u|_0+\varepsilon [u]^*_2
\end{equation}


% PDF page 143
\source{gilbarg-trudinger:page-0143}

for any $\varepsilon>0$. Let $x$ be any point in $\Omega$, $d_x$ its distance from $\partial\Omega$, and $\mu\le \frac12$ a positive constant to be specified later. We set $d=\mu d_x$ and $B=B_d(x)$. For any $i=1,2,\dots,n$, let $x'$, $x''$ be the endpoints of the segment of length $2d$ parallel to the $x_i$ axis with its center at $x$. Then for some $\bar{x}$ in this segment we have

\[
\lvert D_i u(\bar{x})\rvert=\frac{\lvert u(x')-u(x'')\rvert}{2d}\le \frac{1}{d}\lvert u\rvert_0
\]

and

\[
\lvert D_i u(x)\rvert\le \lvert D_i u(\bar{x})+\int_{\bar{x}}^x D_{ii}u\,dx_i\rvert
\le \frac{1}{d}\lvert u\rvert_0+d\sup_B \lvert D_{ii}u\rvert
\]

\[
\le \frac{1}{d}\lvert u\rvert_0+d\sup_{y\in B} d_y^{-2}\sup_{y\in B} d_y^2\lvert D_{ii}u(y)\rvert.
\]

Since $d_y>d_x-d=(1-\mu)d_x\ge d_x/2$ for all $y\in B$, it follows that

\[
d_x\lvert D_i u(x)\rvert\le \mu^{-1}\lvert u\rvert_0+4\mu\sup_{y\in\Omega} d_y^2\lvert D_i u(y)\rvert\le \mu^{-1}\lvert u\rvert_0+4\mu [u]_2^*.
\]

Hence

\[
[u]_1^*=\sup_{\substack{x\in\Omega\\ i=1,\dots,n}} d_x\lvert D_i u(x)\rvert\le \mu^{-1}\lvert u\rvert_0+4\mu [u]_2^*.
\]

If now $\mu$ is chosen so that $\mu\le \varepsilon/4$ we conclude (6.82) with $C=\mu^{-1}$.

(ii) $j\le k$; $\beta=0$, $\alpha>0$. We proceed in a similar manner. As before let $x\in\Omega$, $0<\mu\le \frac12$, $d=\mu d_x$, $B=B_d(x)$, and let $x'$, $x''$ be the endpoints of the segment of length $2d$ parallel to the $x_i$ axis with its center at $x$. For some $\bar{x}$ on this segment we have

\[
\text{(6.84)}\qquad
\lvert D_{i}u(x)\rvert=\frac{\lvert D_i u(x')-D_i u(x'')\rvert}{2d}\le \frac{1}{d}\sup_B \lvert D_i u\rvert
\]

and

\[
\lvert D_i u(x)\rvert\le \lvert D_i u(\bar{x})\rvert+\lvert D_i u(x)-D_i u(\bar{x})\rvert
\]

\[
\le \frac{1}{d}\sup_{y\in B} d_y^{-1}\sup_{y\in B} d_y\lvert D_i u(y)\rvert
+ d^{\alpha}\sup_{y\in B} d_{x,y}^{-2-\alpha}\sup_{y\in B} d_{x,y}^{2+\alpha}\frac{\lvert D_i u(x)-D_i u(y)\rvert}{\lvert x-y\rvert^{\alpha}}.
\]

Again, since $d_y,d_x > d_x/2$ for all $y\in B$, it follows that

\[
d_x^2\lvert D_i u(x)\rvert\le \frac{2}{\mu}[u]_1^*+2^{2+\alpha}\mu^{\alpha}[u]_{2,\alpha}^*.
\]

% PDF page 144
\source{gilbarg-trudinger:page-0144}

Taking the supremum over $i$, $l$ and $x\in\Omega$, choosing $\mu$ so that $8\mu^{\alpha}\le \varepsilon$ and setting
$C=2/\mu$, we obtain the inequality

\[
\text{(6.85)}\qquad [u]^*_{2}\le C(\varepsilon)[u]^*_{1}+\varepsilon [u]^*_{2,\alpha}.
\]

If $u$ replaces $D_i u$ in (6.84) and the obvious modifications are made in the ensuing details, we obtain (6.82) for $j=k=1$, $\beta=0$, $\alpha>0$:

\[
\text{(6.86)}\qquad [u]^*_{1}\le C(\varepsilon)\lvert u\rvert_{0}+\varepsilon [u]^*_{1,\alpha}.
\]

Combining (6.85) with (6.83) after appropriate choice of $\varepsilon$ in each of these inequalities, we arrive at (6.82) for $k=2$ and $j=1,2$.

(iii) $j<k$; $\beta>0$, $\alpha=0$. Let $x$, $y\in\Omega$ with $d_x\le d_y$, so that $d_x=d_{x,y}$. Let $\mu$, $d$, and $B$ be defined as before. We prove (6.82) for the case $j=0$, first establishing the interpolation inequality

\[
[u]^*_{0,\beta}\le C(\varepsilon)\lvert u\rvert_{0}+\varepsilon [u]^*_{1}
\]

where $\varepsilon>0$ may be arbitrary for $0<\beta<1$. If $y\in B$, we obtain from the theorem of the mean, for $0<\beta\le 1$,

\[
d_x^{\beta}\frac{\lvert u(x)-u(y)\rvert}{\lvert x-y\rvert^{\beta}}
\le \mu^{1-\beta}d_x\lvert Du\rvert_{0;B}\le 2\mu^{1-\beta}[u]^*_{1};
\]

and if $y\notin B$, we have

\[
\text{(6.87)}\qquad d_x^{\beta}\frac{\lvert u(x)-u(y)\rvert}{\lvert x-y\rvert^{\beta}}\le 2\mu^{-\beta}\lvert u\rvert_{0}.
\]

Combining these inequalities, we obtain for $0<\beta\le 1$

\[
\text{(6.88)}\qquad [u]^*_{0,\beta}=\sup_{x,y\in\Omega}d_{x,y}^{\beta}\frac{\lvert u(x)-u(y)\rvert}{\lvert x-y\rvert^{\beta}}
\le 2\mu^{-\beta}\lvert u\rvert_{0}+2\mu^{1-\beta}[u]^*_{1}.
\]

This implies (6.82) when $\beta<1$ and $2\mu^{1-\beta}\le \varepsilon$. Applying (6.83) to the right member of (6.88) and choosing $\mu$ appropriately, we obtain (6.82) for $j=0$, $k=2$, $\alpha=0$, $0<\beta\le 1$. The proof for $j=1$, $k=2$ proceeds in much the same way after replacing $u$ with $D_i u$. There is the following difference, however. In place of (6.87) we now have the inequality,

\[
d_x^{1+\beta}\frac{\lvert D_i u(x)-D_i u(y)\rvert}{\lvert x-y\rvert^{\beta}}
\le \mu^{-\beta}\bigl[d_x\lvert D_i u(x)\rvert+d_y\lvert D_i u(y)\rvert\bigr]\le 2\mu^{-\beta}[u]^*_{1}.
\]

% PDF page 145
\source{gilbarg-trudinger:page-0145}



The conclusion again follows by application of (6.83).

\medskip

\noindent (iv) $j\leq k$; $\alpha,\beta>0$. It suffices to take $j=k$ and hence $\alpha>\beta$. With the same notation as above, we have for $y\in B$

\[
d_x^{\beta}\frac{\abs{u(x)-u(y)}}{\abs{x-y}^{\beta}}\leq \mu^{\alpha-\beta}d_x^{\alpha}\frac{\abs{u(x)-u(y)}}{\abs{x-y}^{\alpha}},
\]

\noindent while if $y\notin B$,

\[
d_x^{\beta}\frac{\abs{u(x)-u(y)}}{\abs{x-y}^{\beta}}\leq 2\mu^{-\beta}\abs{u}_{0}.
\]

Combining these inequalities and taking the supremum over $x$, $y\in\Omega$, we obtain
(6.82) for $j=k=0$ with $\epsilon=\mu^{\alpha-\beta}$ and $C=2/\mu^{\beta}$. The remaining cases when $j=k=1,2$
follow in a similar way with the use of results from case (ii).

The interpolation inequality (6.82) for seminorms immediately implies the norm
inequality

\[
\lvert u\rvert^*_{j,\beta;\Omega}\leq C\lvert u\rvert_{0;\Omega}+\epsilon [u]^*_{k,\alpha;\Omega}.
\]

\medskip

An application of Lemma 6.32 yields the following compactness result.

\medskip

\noindent \textbf{Lemma 6.33.} \textit{Let $\Omega$ be a bounded open set in $\mathbb{R}^n$ and let $S$ be a bounded subset of
the Banach space}

\[
C_*^{k,\alpha}=\{u\in C^{k,\alpha}(\Omega)\mid \lvert u\rvert^*_{k,\alpha;\Omega}<\infty\}, \qquad k=0,1,2,\dots, 0\leq \alpha\leq 1.
\]

\noindent \textit{Suppose the functions of $S$ are also equicontinuous on $\bar{\Omega}$. Then if $k+\alpha>j+\beta$, it
follows that $S$ is precompact in $C_*^{j,\beta}(\Omega)$.}

\medskip

\noindent \textit{Proof.} Since $S$ is equicontinuous on $\bar{\Omega}$ and bounded in $C_*^{k,\alpha}$, it contains a sequence
$\{u_m\}$ that converges uniformly on $\bar{\Omega}$ to a function $u\in C_*^{k,\alpha}$. By hypothesis we may
assume $\lvert u_m\rvert^*_{k,\alpha;\Omega}\leq M$ (independent of $m$). From (6.82) we have for any $\epsilon>0$ and
some constant $C=C(\epsilon)$

\[
\lvert u_m-u\rvert^*_{j,\beta}\leq C\lvert u_m-u\rvert_{0}+\epsilon \lvert u_m-u\rvert^*_{k,\alpha}.
\]

If now $N$ is so large that $\lvert u_m-u\rvert_{0}\leq \epsilon/C$ for all $m>N$, it follows that
$\lvert u_m-u\rvert^*_{j,\beta}\leq \epsilon(1+2M)$ for $m>N$. Thus $\{u_m\}$ converges to $u$ in $C_*^{j,\beta}$, which proves the assertion
of the lemma.

\medskip

We now extend Lemma 6.32 to partially interior norms and seminorms in
domains with a hyperplane boundary portion.


% PDF page 146
\source{gilbarg-trudinger:page-0146}


\noindent \textbf{Lemma 6.34.} \textit{Suppose $j+\beta<k+\alpha$, where $j,k=0,1,2,\dots$, and $0<\alpha,\beta\leq 1$.
Let $\Omega$ be an open subset of $\mathbb{R}_+^n$ with a boundary portion $T$ on $x_n=0$, and assume
$u\in C^{k,\alpha}(\Omega\cup T)$. Then for any $\epsilon>0$ and some constant $C=C(\epsilon,j,k)$ we have}

\begin{equation}
\tag{6.89}
[u]^*_{j;\beta;\Omega\cup T}\leq C|u|_{0;\Omega}+\epsilon [u]^*_{k;\alpha;\Omega\cup T},
\end{equation}
\[
|u|^*_{j;\beta;\Omega\cup T}\leq C|u|_{0;\Omega}+\epsilon [u]^*_{k;\alpha;\Omega\cup T}.
\]

\noindent \textit{Proof.} Again we suppose that the right members are finite. The proof is patterned
after that of Lemma 6.32 and we emphasize only the details in which the proofs
differ. In the following we omit the subscript $\Omega\cup T$, which will be implicitly
understood.

\par We consider first the cases $1\leq j\leq k\leq 2$, $\beta=0$, $\alpha\geq 0$, starting with the inequality

\begin{equation}
\tag{6.90}
[u]^*_{2}\leq C(\epsilon)|u|_{0}+\epsilon [u]^*_{2,\alpha}, \qquad \alpha>0.
\end{equation}

\par Let $x$ be any point in $\Omega$, $d_x$ its distance from $\partial\Omega - T$, and $d=\mu d_x$ where $\mu\leq \frac14$
is a constant to be specified later. If $\operatorname{dist}(x,T)\geq d$, then the ball $B_d(x)\subset\Omega$ and the
argument proceeds as in Lemma 6.32, leading to the inequality

\[
\bar d_x^2|D_{il}u(x)|\leq C(\epsilon)[u]^*_1+\epsilon [u]^*_{2,\alpha}
\]

\noindent provided $\mu=\mu(\epsilon)$ is chosen sufficiently small. If $\operatorname{dist}(x,T)<d$ we consider the ball
$B=B_d(x_0)\subset\Omega$, where $x_0$ is on the perpendicular to $T$ passing through $x$ and
$\operatorname{dist}(x,x_0)=d$. Let $x',x''$ be the endpoints of the diameter of $B$ parallel to the $x_l$
axis. Then we have for some $\bar x$ on this diameter

\[
|D_{il}u(\bar x)|=\frac{|D_{i_l}u(x')-D_{i_l}u(x'')|}{2d}\leq \frac{1}{d}\sup_B |D_i u|
\leq \frac{2}{\mu}\bar d_x^{-2}\sup_{y\in B}\bar d_y |D_i u(y)|
\]

\[
\leq \frac{2}{\mu}\bar d_x^{-2}[u]^*_1, \qquad \text{since }\bar d_y>\bar d_x/2\text{ for all }y\in B;
\]

\noindent and

\[
|D_{il}u(x)|\leq |D_{il}u(\bar x)|+|D_{il}u(x)-D_{il}u(\bar x)|
\]

% [illegible]

\noindent hence

\[
\bar d_x^2|D_{il}u(x)|\leq \frac{2}{\mu}[u]^*_1+16\mu^\alpha [u]^*_{2,\alpha}
\]

\[
\leq C[u]^*_1+\epsilon [u]^*_{2,\alpha},
\]

\noindent provided $16\mu^\alpha\leq \epsilon$, $C=2/\mu$. Choosing the smaller value of $\mu$, corresponding to the
two cases $\operatorname{dist}(x,T)\geq d$ and $\operatorname{dist}(x,T)<d$, and taking the supremum over all


% PDF page 147
\source{gilbarg-trudinger:page-0147}

\pagestyle{empty}


\setcounter{section}{6}
\noindent\textbf{6.8.\ Appendix 1: Interpolation Inequalities}\hfill 135

\medskip

$x \in \Omega$ and $i,\, l = 1, \dots, n$, we obtain (6.90). If $u$ replaces $D_l u$ in the preceding and the details are modified accordingly, we obtain (6.90) for $j = k = 1$.

To prove (6.89) for $j = 1$, $k = 2$, $\alpha = \beta = 0$, we proceed as in Lemma 6.32 with the modifications suggested by the above proof of (6.90). Together with the preceding cases, this gives us (6.89) for $1 \leq j \leq k \leq 2$, $\beta = 0$, $\alpha \geq 0$.

The proof of (6.89) for $\beta > 0$ follows closely that in cases (iii) and (iv) of Lemma 6.32. The principal difference is that the argument for $\beta > 0$, $\alpha = 0$ now requires application of the theorem of the mean in the truncated ball $B_d(x)\cap \Omega$ for points $x$ such that $\operatorname{dist}(x,T) < d$. \hfill $\square$

\medskip

\noindent We conclude this section with the proof of \emph{global interpolation inequalities} in smooth domains.

\medskip

\noindent\textbf{Lemma 6.35.} \emph{Suppose $j+\beta < k+\alpha$, where $j = 0, 1, 2, \dots, k = 1, 2, \dots,$ and $0 \leq \alpha, \beta \leq 1$. Let $\Omega$ be a $C^{k,\alpha}$ domain in $\mathbb{R}^n$, and assume $u \in C^{k,\alpha}(\overline{\Omega})$. Then for any $\varepsilon > 0$ and some constant $C = C(\varepsilon, j, k, \Omega)$ we have}

\begin{equation}\tag{6.91}
|u|_{j,\beta;\Omega} \leq C|u|_{0;\Omega} + \varepsilon |u|_{k,\alpha;\Omega}.
\end{equation}

\noindent\emph{Proof.} The proof is based on a reduction to Lemma 6.34 by means of an argument very similar to that in Lemma 6.5. As in that lemma, at each point $x_0 \in \partial\Omega$ let $B_\rho(x_0)$ be a ball and let $\psi$ be a $C^{k,\alpha}$ diffeomorphism that straightens the boundary in a neighborhood containing $B' = B_\rho(x_0)\cap\Omega$ and $T = B_\rho(x_0)\cap\partial\Omega$. Let $\psi(B') = D' \subset \mathbb{R}_+^n$, $\psi(T) = T' \subset \partial\mathbb{R}_+^n$. Since $T'$ is a hyperplane portion of $\partial D'$, we may apply the interpolation inequality (6.89) in $D'$ to the function $\tilde{u} = u \circ \psi^{-1}$ to obtain

\begin{equation*}
|\tilde{u}|_{j,\beta;D' \cup T'} \leq C(\varepsilon)|\tilde{u}|_{0;D'} + \varepsilon |\tilde{u}|_{k,\alpha;D' \cup T'}.
\end{equation*}

\noindent From (6.30) it follows

\begin{equation*}
|u|^{*}_{j,\beta;B' \cup T} \leq C(\varepsilon)|u|_{0;B'} + \varepsilon |u|^{*}_{k,\alpha;B' \cup T}.
\end{equation*}

\noindent (We recall that the same notation $C(\varepsilon)$ is being used for different functions of $\varepsilon$.) Letting $B'' = B_{\rho/2}(x_0)\cap\Omega$, we infer from (4.17)', (4.17)''

\begin{equation}\tag{6.92}
|u|_{j,\beta;B''} \leq C(\varepsilon)|u|_{0;B'} + \varepsilon |u|_{k,\alpha;B'}
 \leq C(\varepsilon)|u|_{0;\Omega} + \varepsilon |u|_{k,\alpha;\Omega}.
\end{equation}

\noindent Let $B_{\rho_i/4}(x_i)$, $x_i \in \partial\Omega$, $i = 1, \dots, N$, be a finite collection of balls covering $\partial\Omega$, such that the inequality (6.90) holds in each set $B_i'' = B_{\rho_i/2}(x_i)\cap\Omega$, with a constant $C_i(\varepsilon)$. Let $\delta = \min \rho_i/4$ and $C = C(\varepsilon) = \max C_i(\varepsilon)$. Then at every point $x_0 \in \partial\Omega$, we have $B = B_\delta(x_0) \subset B_{\rho_i/2}(x_i)$ for some $i$ and hence

\begin{equation}\tag{6.93}
|u|_{j,\beta;B\cap\Omega} \leq C|u|_{0;\Omega} + \varepsilon |u|_{k,\alpha;\Omega}.
\end{equation}

\noindent\emph{The remainder of the argument is analogous to that in Theorem 6.6 and is left to the reader.} \hfill $\square$


% PDF page 148
\source{gilbarg-trudinger:page-0148}

\begin{flushleft}
The global interpolation inequality (6.91) is valid in more general domains, for
example in $C^{0,1}$ domains; (see Problem 6.7). However, as shown in the example on
page 53, suitable regularity of the domain $\Omega$ is required to insure the inclusion
relation $C^{k,\alpha}(\bar{\Omega}) \subset C^{j,\beta}(\bar{\Omega})$ when $k+\alpha > j+\beta$, and hence the global interpolation
inequalities are not true in arbitrary domains.
\end{flushleft}

\begin{flushleft}
Lemma 6.35 implies the following compactness result.
\end{flushleft}

\begin{flushleft}
\textbf{Lemma 6.36.} \emph{Let $\Omega$ be a $C^{k,\alpha}$ domain in $\R^n$ (with $k \ge 1$) and let $S$ be a bounded
set in $C^{k,\alpha}(\bar{\Omega})$. Then $S$ is precompact in $C^{j,\beta}(\bar{\Omega})$ if $j+\beta < k+\alpha$.}
\end{flushleft}

\begin{flushleft}
The proof is essentially the same as that of Lemma 6.33 and is therefore omitted.
The result is obviously valid for domains in which the global interpolation inequality (6.91) holds, hence in $C^{0,1}$ domains.
\end{flushleft}

\begin{flushleft}
\textbf{6.9. Appendix 2: Extension Lemmas}
\end{flushleft}

\begin{flushleft}
This section establishes some results needed earlier in this chapter and elsewhere
in this work concerning the extension of globally defined functions into larger
domains and the extension of functions defined on the boundary to globally
defined functions.
\end{flushleft}

\begin{flushleft}
We shall use the concept of a partition of unity. Let $\Omega$ be an open set in $\R^n$
covered by a countable collection $\{\Omega_j\}$ of open sets $\Omega_j$. A countable set of functions $\{\eta_i\}$ is a \emph{locally finite partition of unity subordinate to the covering $\{\Omega_j\}$} if:
(i) $\eta_i \in C_0^\infty(\Omega_j)$ for some $j=j(i)$; (ii) $\eta_i \ge 0$, $\sum \eta_i = 1$ in $\Omega$; (iii) at each point of $\Omega$
there is a neighborhood in which only a finite number of the $\eta_i$ are non-zero. For
the proof of existence of such a partition we refer to the literature; (for example,
see [YO], also Problem 6.8). In the following applications the construction of the
partition is relatively simple.
\end{flushleft}

\begin{flushleft}
\textbf{Lemma 6.37.} \emph{Let $\Omega$ be a $C^{k,\alpha}$ domain in $\R^n$ (with $k \ge 1$) and let $\Omega'$ be an open set
containing $\bar{\Omega}$. Suppose $u \in C^{k,\alpha}(\bar{\Omega})$. Then there exists a function $w \in C_0^{k,\alpha}(\Omega')$ such that
$w = u$ in $\Omega$ and}
\end{flushleft}

\[
\tag{6.94}
\lvert w \rvert_{k,\alpha;\Omega'} \le C \lvert u \rvert_{k,\alpha;\Omega},
\]

\begin{flushleft}
where $C = C(k,\Omega,\Omega')$.
\end{flushleft}

\begin{flushleft}
\emph{Proof.} Let $y=\psi(x)$ define a $C^{k,\alpha}$ diffeomorphism that straightens the boundary
near $x_0 \in \partial \Omega$, and let $G$ and $G^+ = G \cap \R^n_+$ be respectively a ball and half-ball in
the image of $\psi$ such that $\psi(x_0) \in G$. Setting $\tilde u(y)=u \circ \psi^{-1}(y)$ and $y=(y_1,\ldots,$
$y_{n-1},y_n)=(y',y_n)$, we define an extension of $\tilde u(y)$ into $y_n<0$ by
\end{flushleft}

\[
\tilde u(y',y_n)=\sum_{i=1}^{k+1} c_i \tilde u(y',-y_n/i), \qquad y_n<0,
\]

% PDF page 149
\source{gilbarg-trudinger:page-0149}

where $c_1,\dots,c_{k+1}$ are constants determined by the system of equations
\[
  \sum_{i=1}^{k+1} c_i (-1/i)^m = 1, \qquad m=0,1,\dots,k.
\]
One verifies readily that the extended function $\tilde u$ is continuous with all derivatives
up to order $k$ in $G$ and that $\tilde u \in C^{k,\alpha}(G)$. Thus $w=\tilde u \circ \psi \in C^{k,\alpha}(\overline{B})$ for some ball
$B=B(x_0)$ and $w=u$ in $B\cap\Omega$, so that $w$ provides a $C^{k,\alpha}$ extension of $u$ into $\Omega\cup B$.
By (6.30) the inequality (6.94) holds (with $\Omega\cup B$ in place of $\Omega'$).

Now consider a finite covering of $\partial\Omega$ by balls $B_i$, $i=1,\dots,N$, such as $B$ in the
preceding, and let $\{w_i\}$ be the corresponding $C^{k,\alpha}$ extensions. We may assume the
balls $B_i$ are so small that their union with $\Omega$ is contained in $\Omega'$. Let $\Omega_0\Subset\Omega$
be an open subset of $\Omega$ such that $\Omega_0$ and the balls $B_i$ provide a finite open cover
of $\Omega$. Let $\{\eta_i\}$, $i=0,1,\dots,N$, be a partition of unity subordinate to this covering,
and set
\[
  w=u\eta_0+\sum w_i\eta_i
\]
with the understanding that $w_i\eta_i=0$ if $\eta_i=0$. One verifies from the above discussion
that $w$ is an extension of $u$ into $\Omega'$ and has the properties asserted in the lemma. $\square$

The following result provides an extension of a boundary function to a globally
defined function in the same regularity class.

\textbf{Lemma 6.38.} \textit{Let $\Omega$ be a $C^{k,\alpha}$ domain in $\mathbb{R}^n$ $(k\ge 1)$ and let $\Omega'$ be an open set containing
$\overline{\Omega}$. Suppose $\varphi\in C^{k,\alpha}(\partial\Omega)$. Then there exists a function $\Phi\in C_0^{k,\alpha}(\Omega')$ such that
$\Phi=\varphi$ on $\partial\Omega$.}

\textit{Proof.} At any point $x_0\in\partial\Omega$ let the mapping $\psi$ and the ball $G$ be defined as in the
preceding lemma, and assume that $\tilde\varphi=\varphi\circ\psi^{-1}\in C^{k,\alpha}(G\cap\partial\mathbb{R}^n_+)$. We define
$\tilde\Phi(y',y_n)=\tilde\varphi(y')$ in $G$ and set $\Phi(x)=\tilde\Phi\circ\psi(x)$ for $x\in\psi^{-1}(G)$. Clearly $\Phi\in C^{k,\alpha}(\overline{B})$
for some ball $B=B(x_0)$ and $\Phi=\varphi$ on $B\cap\partial\Omega$. Now let $\{B_i\}$ be a finite covering of
$\partial\Omega$ by balls such as $B$, and let $\Phi_i$ be the corresponding $C^{k,\alpha}$ functions defined on $B_i$.
The proof can now be completed as in the preceding lemma by use of an appropriate
partition of unity. $\square$

\textit{Remarks.} 1) In the preceding lemma, if $\varphi\in C^0(\partial\Omega)\cap C^{k,\alpha}(T)$ where $T\subset\partial\Omega$,
then the same argument leads to an extension $\Phi\in C^0(\Omega')\cap C^{k,\alpha}(G)$, where $G$ is an
open set containing $T$. A simple modification of the above proof shows that if $\Omega$
is any domain with a $C^{k,\alpha}$ boundary portion $T$ and if $\varphi\in C^{k,\alpha}(T)$, then $\varphi$ can be
extended to a function $\Phi\in C^{k,\alpha}(G)$, where G is an open set containing $T$ and $\Phi=\varphi$
on $T$. A countable covering of $T$ by balls is required for the argument.
If $\varphi\in C^0(\partial\Omega)\cap C^{k,\alpha}(T)$ then the extension $\Phi$ can be determined so that
$\Phi\in C^0(\overline{\Omega})\cap C^{k,\alpha}(G)$. 2) For domains with a simple geometry the construction
of an extended function can often be made directly and easily. Thus if $B=B_R(x_0)$
is a ball in $\mathbb{R}^n$ and $\varphi\in C^0(\partial\Omega)\cap C^{k,\alpha}(T)$, $T\subset\partial B$, then an extension of $\varphi$ into
$\mathbb{R}^n$ can be obtained by setting
\[
  \Phi(x)=\Phi(x_0+r\omega)=\varphi(R\omega)\eta(r),
\]

% PDF page 150
\source{gilbarg-trudinger:page-0150}

where $r = |x-x_0|$, $\omega = (x-x_0)/r$, and $\eta(r)$ is a $C^\infty$ cut-off function such that $\eta(r)=0$ for $0 \le r \le R/4$, $\eta(r)=1$ for $r \ge R/2$. The function $\Phi(x)$ obviously coincides with $\varphi$ on $\partial B$, is in $C^0(\R^n)$ and is of class $C^{k,\alpha}$ in the conical region determined by the rays from the origin through the points of $T$.

\bigskip

\noindent\textbf{Notes}

The apriori estimates and existence theory in Sections 1--3 are, in modified form, the contributions of Schauder [SC 4, 5]. At about the same time, Caccioppoli [CA 1] stated without details similar results, which were elaborated by Miranda [MR 2]. Closely related ideas are contained in the work of Hopf [HO 3] who earlier established the interior regularity theorems of Section 4. The existence theory and general properties of solutions for essentially the same class of problems were previously obtained by Giraud [GR 1--3], who used the method of integral equations based on representing solutions as surface potentials. Further details amplifying on the respective contributions are discussed by Miranda [MR 2]. A development of the Schauder estimates based on methods of Fourier analysis, for equations of arbitrary order, is contained in H\"ormander [HM 3].

The formulation of the interior estimates of Section 1 in terms of interior norms and the method of derivation are patterned after Douglis and Nirenberg [DN], who also extended the interior estimates to elliptic systems. The details of proof of Theorem 6.2 can be simplified somewhat if the estimates are first carried out in balls, using Theorem 4.6, and at the end are converted into the bound (6.14) for the $C^{2,\alpha}$ interior norm. See, for example, the proof of Theorem 9.11.

The global estimates of Section 6.2 and the proof of Theorem 6.8 based on these estimates assume $C^{2,\alpha}$ boundary data. Under weaker regularity hypotheses an existence proof for say $C^{1,\alpha}(\overline{\Omega}) \cap C^2(\Omega)$ solutions does not follow from the Schauder theory in its usual form. Such an existence theorem is implied by regularity results of Widman [WI 1]. Gilbarg and H\"ormander [GH] have extended the global Schauder theory to include conditions of lower regularity of the coefficients, domain and boundary values. We summarize their results, which apply to conditions of higher regularity as well:

\medskip

\noindent If

\[
0 \le k < a = k + \alpha \le k + 1,
\]

\noindent let $H_a(\Omega)$ denote the Hölder space of functions with finite norm $|u|_{a,\Omega} = |u|_{k,\alpha;\Omega}$ (the latter being in the notation of this book; thus $H_a(\Omega)=C^{k,\alpha}(\overline{\Omega})$). Setting

\[
\Omega_\delta = \{x \in \Omega \mid \dist(x,\partial\Omega) > \delta\},
\]

\noindent let $H_a^{(b)}(\Omega)$ denote the set of functions on $\Omega$ belonging to $H_a(\Omega_\delta)$ for all $\delta > 0$ and with finite norm

\[
|u|_{a,\Omega}^{(b)} = |u|_a^{(b)} = \sup_{\delta>0} \delta^{a+b} |u|_{a,\Omega_\delta},
\]


% PDF page 151
\source{gilbarg-trudinger:page-0151}

where $a+b\ge 0$. Since $H_a^{(-b)}\subset H_b^{(-b)}=H_b$ for $a\ge b>0$, $b$ a noninteger, the upper and lower indices in the norm $|u|_a^{(-b)}$ describe, respectively, the global and interior regularity of $u$. Define also $H_a^{(-0)}(\Omega)$ to be the set of functions in $H_a^{(b)}$ such that $\delta^{a+b}|u|_{a,\Omega_\delta}\to 0$ as $\delta\to 0$. Having these spaces, now let $\Omega$ be a bounded $C^\gamma$ domain for some $\gamma\ge 1$, and $a,b$ nonintegers such that $0<b\le a$, $a>2$, $b\le\gamma$.
Let
\[
P=\sum_{|\beta|\le 2} p_\beta(x)D^\beta
\]
be an elliptic second-order differential operator on $\overline{\Omega}$ such that
\[
p_\beta\in H_{a-2}^{(2-b)}(\Omega) \qquad \text{if }|\beta|\le 2,
\]
\[
p_\beta\in C^0(\overline{\Omega}) \qquad \text{if }|\beta|=2,
\]
\[
p_\beta\in H_{a-2}^{(2-|\beta|-0)}(\Omega) \qquad \text{if }b<|\beta|.
\]
(Thus the lower-order coefficients may be unbounded if $b<2$.) Then, if $u\in C^2(\Omega)\cap C^0(\overline{\Omega})$ is a solution of
\[
(6.95)\qquad Pu=f \text{ in }\Omega,\qquad u=\varphi \text{ on }\partial\Omega
\]
where $f\in H_{a-2}^{(2-b)}(\Omega)$, $\varphi\in H_b(\partial\Omega)$, it follows that $u\in H_a^{(-b)}(\Omega)$ and satisfies an estimate
\[
|u|_a^{(-b)}\le C\bigl(|u|_0+|\varphi|_{b,\partial\Omega}+|f|_{a-2}^{(2-b)}\bigr),
\]
where $C$ depends on $\Omega$, $a$, $b$, the norms of the coefficients and their minimum eigenvalue. If $p_0\le 0$ the Dirichlet problem (6.95) has a unique solution in $H_a^{(-b)}$, and a corresponding Fredholm-type theorem holds in general. The case $2+\alpha\le a=b\le\gamma$ is the one treated in this chapter. If $\Omega$ is a Lipschitz domain, analogous results hold for values $b<1$ that depend on the exterior cone condition satisfied at the boundary; and it suffices that $p_\beta\in H_{a-2}^{(0)}$ for $|\beta|=2$, so that the principal coefficients need not be continuous at the boundary.

The conditions under which a regular boundary point for the Laplacian is also a regular point for an elliptic operator, and conversely, have been studied by several authors. The equivalence has been proved for strictly elliptic operators $L$, as in Theorem 6.13, whose coefficients near the boundary are Lipschitz continuous [HR] or Dini continuous [KV 1], [NO 2]; and also for certain classes of discontinuous coefficients [AK] and degenerate elliptic operators [MM], [NO 3]. Capacity and the Wiener criterion (Section 2.9) play an important part in the arguments. If the coefficients of $L$ are only continuous, the equivalence no longer holds in general (see the example in Problem 3.8(a), also [ML 4]). However, in the case of divergence structure equations, there is equivalence when the coefficients are only bounded and measurable [LSW] (see also Chapter 8). For additional results on regular boundary points, see [NO 1], [MZ], [LN 2], [ML 2, 4].

% PDF page 152
\source{gilbarg-trudinger:page-0152}

Hopf [HO 3] proves directly the interior regularity result, Lemma 6.16,
without an existence theorem. His method, based on Korn's device (in [KR 2]) of
perturbation about the constant coefficient equation, provides an extension of his
results in [HO 2] on the regularity of solutions of variational problems and anti-
cipates important aspects of the Schauder theory. A simple direct proof of Lemma
6.16 (and of more general results), based on regularization and interior estimates,
is contained in [ADN 1], p. 723. For another approach see the proof of Lemma 9.16;
also [MY 5], Section 5.6.

A basically simpler proof of the boundary regularity result, Lemma 6.18, can
be obtained as follows. It suffices to prove that $u \in C^1(\Omega \cup T)$, after which the
argument proceeds essentially as in Lemma 6.16. By considering $u-\varphi$ in place
of $u$, we may assume $\varphi \equiv 0$. Let $\Omega' \subset \Omega$ with $\partial\Omega' \cap \partial\Omega = T' \subset T$, and let
\[
\delta = \dist(\Omega',\partial\Omega-T)<1.
\]
For any $x' \in \Omega'$ suppose first
\[
d=\dist(x',\partial\Omega)=|x'-x_0|\le \delta,\qquad x_0 \in \partial\Omega.
\]
Then by Problem 3.6 we have
\[
|u(x)|\le C|x-x_0| \qquad \text{for } x \in \Omega,
\]
and hence $|u(x)|\le Cd$ for $x \in B_d(x')$. From (6.23), in which we set $\Omega' = B_{d/2}(x')$
and $\Omega = B_d(x')$, it follows that for all $x \in B_{d/2}(x')$,
\[
d|Du(x)|\le C\bigl(\sup_{B_d}|u|+d^2|f|_{0,\alpha;B_d}\bigr),
\]
and hence
\[
|Du(x')|\le C(1+|f|_{0,\alpha;\Omega})\le C,
\]
where the constant $C$ depends only on $\delta$ and the given data. If $d>\delta$, then the same
inequality holds, with the constant $C$ now depending on $\delta^{-1}$. We thus have a
bound on $|Du|$ in $\Omega'$ and consequently $u \in C^{0,1}(\Omega \cup T)$.

Section 6.5 is a modification of the ideas of Michael [MI 1] who has shown that
a general existence theory for continuous boundary values can be developed from
interior estimates only. His results apply as well to certain classes of equations with
unbounded coefficients near the boundary; (see Problems 6.5, 6.6).

Section 6.6 considers some cases of nonuniformly elliptic operators for which
the ellipticity degenerates on the boundary. The theory of elliptic operators that
degenerate in the interior is based on essentially different methods from those of
this chapter. For the relevant results the reader is directed to the literature on
hypoellipticity, e.g., [HM 2], [OR], [KJ].

The Schauder theory of the oblique derivative problem in Section 6.7 differs
in some respects from earlier versions. In particular, Fiorenza [FI 1] based his
approach on the representation by surface potentials of solutions of the boundary
value problem (6.60) for Poisson's equation in a half-space, to which he applied
some of the results of Giraud [GR 3]. His Schauder-type estimates for the case of
variable coefficients exhibit a quite precise dependence on the bounds and Hölder
constants of the coefficients. This dependence is used in [LU 4], Chapter 10, and by

% PDF page 153
\source{gilbarg-trudinger:page-0153}


Fiorenza [FI 2] and Ural'tseva [UR] to treat quasilinear equations with nonlinear
boundary conditions. An extension of the Schauder theory to other boundary
conditions for higher order equations and systems appears in Agmon, Douglis and
Nirenberg [ADN 1, 2]. Their method is based on explicit integral representations of
solutions of the constant coefficient problem in a half-space in terms of appropriate
Poisson kernels for the given boundary conditions. Although the details are quite
different, the development in Section 6.7 can be viewed as a special case of these
results. See also Bouligand [BGD]. The nonregular oblique derivative problem, in
which the directional derivative in the boundary condition becomes tangential
$(\beta_\nu=0$ in (6.76)), is essentially deeper than the regular case and the results are
different; see, for example, [HM 1], [EK], [SJ] and [WZ 1, 2].

\medskip

The solution of exterior boundary value problems can be inferred readily from
the results of this chapter. Meyers and Serrin [MS 1] treat boundary value problems
for the equation $Lu=f$ in an exterior domain $\Omega$ (containing the exterior of some
ball), with $c\leq 0$ and coefficients and $f$ Hölder continuous in bounded subsets of $\Omega$.
Under suitable general hypotheses on the behavior of the coefficients at infinity they
prove the existence of solutions in $\Omega$ with a limit at infinity from the convergence of
solutions in expanding domains. They obtain the result, among others, that if $f=0$,
$b^i=0$, $a^{ij}\to a_0^{ij}$ at $\infty$ and the matrix $[a_0^{ij}]$ has rank $\geq 3$, then there is a unique
solution in $\Omega$ of the Dirichlet (and other) problems vanishing at infinity. Under
these conditions, if $n>3$ the operator $L$ may be non-uniformly elliptic at infinity
and the boundary value problem is still well posed. An extension of the Schauder
theory to exterior domains for $n\geq 3$, including Hölder estimates at infinity and a
corresponding treatment of the exterior Dirichlet and Neumann problems, has
been given by Oskolkov [OS 1]. An exterior Neumann problem for a class of
quasilinear equations when $n\geq 3$ is treated in [FG 2].

\medskip

The interpolation inequalities of Appendix 1 can easily be derived from the
general convexity property of Hölder norms:

\[
\lvert u \rvert_{k,\alpha} \leq C(\lvert u \rvert_{k_1,\alpha_1})^t (\lvert u \rvert_{k_2,\alpha_2})^{1-t},
\]
where $0<t<1$,
\[
k+\alpha = t(k_1+\alpha_1)+(1-t)(k_2+\alpha_2),
\]
and the norms may be interior or global. For proof of this inequality, see Hörmander
[HM 3].

\bigskip

\noindent \textbf{Problems}

\medskip

\noindent \textbf{6.1.} \textbf{(a)} Let $u\in C^{k+2,\alpha}(\Omega)$, $k\geq 0$, be a solution of $Lu=f$ in a bounded open set
$\Omega\subset \mathbb{R}^n$, and assume the coefficients of $L$ satisfy (6.2) and $\lvert a^{ij}, b^i, c\rvert_{k,\alpha;\Omega}\leq \Lambda$. If
$\Omega' \Subset \Omega$, show that

\[
\lvert u\rvert_{k+2,\alpha;\Omega'} \leq C(\lvert u\rvert_{0;\Omega}+\lvert f\rvert_{k,\alpha;\Omega})
\]

\noindent where $C = C(n,k,\alpha,\lambda,\Lambda,d)$, $d = \operatorname{dist}(\Omega',\partial\Omega)$.


% PDF page 154
\source{gilbarg-trudinger:page-0154}


(b) In Theorem 6.2 assume $u \in C^{k+2,\alpha}(\Omega)$, $k \geq 0$, and replace (6.13) by the conditions

\[
  \lvert a^{ij}\rvert_{k,\alpha}, \lvert b^i\rvert_{k,\alpha}, \lvert c\rvert_{k,\alpha} \leq \Lambda.
\]

Prove the interior estimate

\[
  \lvert u\rvert_{k+2,\alpha} \leq C(\lvert u\rvert_{0} + \lvert f\rvert_{k,\alpha}^{(2)})
\]

where $C = C(n, k, \alpha, \lambda, \Lambda)$.

\medskip

\noindent \textbf{6.2.} In Theorem 6.6 let $\Omega$ be a $C^{k+2,\alpha}$ domain, $k \geq 0$, and assume $u \in C^{k+2,\alpha}(\overline{\Omega})$,
$\varphi \in C^{k+2,\alpha}(\overline{\Omega})$, $f \in C^{k,\alpha}(\overline{\Omega})$, $\lvert a^{ij}, b^i, c\rvert_{k,\alpha;\Omega} \leq \Lambda$. Prove the global estimate

\[
  \lvert u\rvert_{k+2,\alpha} \leq C(\lvert u\rvert_{0} + \lvert\varphi\rvert_{k+2,\alpha} + \lvert f\rvert_{k,\alpha})
\]

where $C = C(n, k, \alpha, \lambda, \Lambda, \Omega)$.

\medskip

\noindent \textbf{6.3.} Prove Theorem 6.13 for a bounded domain $\Omega$ satisfying an exterior cone
condition. Show that at each point $x_0 \in \partial\Omega$ there is a local barrier determined by
a function of the form $r^\mu f(\theta)$, where $r = \lvert x - x_0\rvert$ and $\theta$ is the angle between the vector
$x - x_0$ and the axis of the exterior cone; (cf. [ML 1, 3]).

\medskip

\noindent \textbf{6.4.} Prove the following extension of Corollary 6.24'. Let $\Omega$ be a bounded strictly
convex domain in $\mathbb{R}^n$ and let the equation

\[
  Lu \equiv a^{ij}D_{ij}u + b^iD_i u = 0
\]

\noindent be elliptic in $\Omega$ with coefficients in $C^\alpha(\Omega)$. For $x_0 \in \partial\Omega$ let $v = v(x_0)$ denote the
unit normal to a supporting plane (directed outward from $\Omega$). At each $x_0$ suppose
$\mathbf{b}\cdot v > 0$ in $B(x_0)\cap\Omega$ for some ball $B(x_0)$. Then the Dirichlet problem for $Lu = 0$
is (uniquely) solvable in $C^{2,\alpha}(\Omega) \cap C^0(\overline{\Omega})$ for arbitrary continuous boundary
values.

\medskip

\noindent \textbf{6.5.} (a) Prove Lemma 6.21 if the ball $B$ is replaced by an arbitrary $C^2$ domain
$\Omega$; (see Problem 4.6).

\noindent \hspace*{1.5em} (b) Extend part (a) to admit coefficients $b^i$ satisfying $\sup_\Omega d_x^{1-\gamma}\lvert b^i(x)\rvert < \infty$ for
some $\gamma \in (0,1)$ and coefficients $c \leq 0$.

\medskip

\noindent \textbf{6.6.} (a) Use the argument in Problem 6.5 to construct a barrier and prove the
solvability of the Dirichlet problem, $Lu=f$ ($c \leq 0$) in $\Omega$, $u = 0$ on $\partial\Omega$, where $L$ is
strictly elliptic in the $C^2$ domain $\Omega$, and the following conditions are satisfied:
the coefficients $a^{ij}$ are bounded; $a^{ij}, b^i, c, f \in C^\alpha(\Omega)$; $\sup_\Omega d_x^{1-\beta}\lvert b^i(x)\rvert < \infty$ and
\[
  \sup_\Omega d_x^{2-\beta}\lvert f(x)\rvert < \infty
\]
for some $\beta \in (0,1)$.


% PDF page 155
\source{gilbarg-trudinger:page-0155}

{\bf Problems}\hfill 143

\medskip

(b) Under the additional hypothesis, $\sup d_x^{2-\beta}|c(x)|<\infty$, extend part (a) to include the boundary condition $u=\varphi$ on $\partial\Omega$, where $\varphi$ is continuous.

\medskip

{\bf 6.7.} Prove the global interpolation inequality (6.91) if $\Omega$ is a $C^{0,1}$ (Lipschitz) domain. The result can be obtained as follows:

\medskip

\noindent
(i) Prove there is a constant $K$ depending only on $\Omega$ such that every pair of points $x, y$ in $\Omega$ can be connected by an arc $\gamma(x,y)$ in $\Omega$ whose length $|\gamma(x,y)| \le K|x-y|$.

\noindent
(ii) Show that for some constants $\rho_0$ and $L$ depending only on $\Omega$, if $y\in\Omega$ and $\dist(y,\partial\Omega)<\rho_0$, then for all $\rho<\rho_0$ there is a point $x\in B_\rho(y)$ such that $B_{\rho/L^2}(x)\subset \Omega$.

\noindent
(iii) Use parts (i) and (ii) in a modification of the proof of Lemma 6.34 to establish (6.91).

\medskip

{\bf 6.8.} Let $\{\Omega_i\}$ be a countable open covering of an open set $\Omega$ in $\mathbb{R}^n$. If either
(a) $\Omega$ is bounded and $\overline{\Omega}\subset \bigcup \Omega_i$, or (b) $\Omega_i$ is bounded and $\overline{\Omega}_i\subset\Omega$, $i=1,2,\dots$, prove the existence of a partition of unity $\{\eta_i\}$ such that $\eta_i\in C_0^\infty(\Omega_i)$.

\medskip

{\bf 6.9.} (a) Use a partition of unity and the definition of $C^{k,\alpha}$ domains in Section 6.2 to show that any such domain $\Omega$ with $k\ge1$ can be defined by a function $F\in C^{k,\alpha}(\overline{\Omega})$, such that $F>0$ in $\Omega$, $F=0$ on $\partial\Omega$ and $\grad F\ne0$ on $\partial\Omega$.

(b) Use part (a) and approximation by smooth functions to show that any $C^{k,\alpha}$ domain $\Omega(k\ge1)$, defined by $F>0$, can be exhausted by arbitrarily smooth domains $\Omega_\nu$, defined by $F_\nu>0$, such that

\[
\partial\Omega_\nu \to \partial\Omega_\nu,\quad\text{and}\quad |F_\nu|_{k,\alpha;\Omega}\le C|F|_{k,\alpha;\Omega}\qquad\text{as }\nu\to\infty,
\]

where $C$ is a constant independent of $k$.

\medskip

{\bf 6.10.} Let $L$ be elliptic in a $C^2$ domain $\Omega$, with $c\le0$ and $a^{ij}$, $b^i\in C^0(\overline{\Omega})$. Let $\Sigma=\Sigma_1\cup\Sigma_2\subset\partial\Omega$ be defined by (6.59). Suppose $u\in C^0(\overline{\Omega})\cap C^2(\overline{\Omega}-\Sigma)$ satisfies $Lu\ge0$ in $\Omega$. Then, if either $c<0$, or $a^{ii}>0$, for some $i$, on $\partial\Omega-\Sigma$, prove that

\[
\sup_\Omega u \le \sup_\Sigma u^+.
\]

(Use the distance function to flatten the boundary near an assumed maximum on $\partial\Omega-\Sigma$, and consider the differential equation there. Cf. [OR].)

\medskip

{\bf 6.11.} Under the hypotheses of Theorem 6.30, but assuming $|\beta_i|_{0,\alpha}\le\Lambda$ in place of $|\beta_i|_{1,\alpha}\le\Lambda$, prove the estimate

\[
|u|_{2,\alpha}\le C\bigl(|u|_0+|\varphi|_{1,\alpha}+|f|_{0,\alpha}+|Du|_0\cdot [D\beta]_\alpha\bigr),
\]

where $C=C(n,\alpha,\lambda,\Lambda,\kappa,\Omega)$.

% PDF page 156
\source{gilbarg-trudinger:page-0156}

Chapter 7

\begin{center}
\Large Sobolev Spaces
\end{center}

To motivate the theory of this chapter we now consider a different approach to
Poisson's equation from that of Chapter 4. By the divergence theorem (equation
(2.3)) a $C^2(\Omega)$ solution of $\Delta u = f$ satisfies the integral identity

\begin{equation}
(7.1)\qquad \int_{\Omega} Du \cdot D\varphi \, dx = - \int_{\Omega} f\varphi \, dx
\end{equation}

for all $\varphi \in C_0^1(\Omega)$. The bilinear form

\begin{equation}
(7.2)\qquad (u,\varphi) = \int_{\Omega} Du \cdot D\varphi \, dx
\end{equation}

is an inner product on the space $C_0^1(\Omega)$ and the completion of
$C_0^1(\Omega)$ under the metric induced by (7.2) is consequently a Hilbert
space, which we call $W_0^{1,2}(\Omega)$.

Furthermore, for appropriate $f$ the linear functional $F$ defined by
$F(\varphi) = - \int_{\Omega} f\varphi \, dx$
may be extended to a bounded linear functional on $W_0^{1,2}(\Omega)$. Hence, by
the Riesz representation theorem (Theorem 5.7), there exists an element
$u \in W_0^{1,2}(\Omega)$ satisfying $(u,\varphi) = F(\varphi)$ for all
$\varphi \in C_0^1(\Omega)$. Thus, the existence of a generalized solution to the
Dirichlet problem, $\Delta u = f$, $u = 0$ on $\partial \Omega$, is readily
established. The question of classical existence is accordingly transformed into
the question of regularity of generalized solutions under appropriately smooth
boundary conditions. In the next chapter, the Lax-Milgram theorem (Theorem 5.8)
will be applied to linear elliptic equations in divergence form in a similar
manner to the above application of the Riesz representation theorem and by means
of various arguments based on integral identities, regularity results will be
established. But before we can so proceed we need to examine the class of
Sobolev spaces, that is, the $W^{k,p}(\Omega)$ and $W_0^{k,p}(\Omega)$ spaces of
which the space $W_0^{1,2}(\Omega)$ is a member. Some of the inequalities we
treat will also be necessary for the development of the theory of quasilinear
equations in Part II.
